% Lutte contre la mauvaise alimentation et l'inactivité physique — SVT 1ère TI — Cours signé OnBuch+
% Programme officiel MINESEC. Compiler avec : tectonic NCT06-lutte-mauvaise-alimentation-inactivite.tex
\newcommand{\DOCMATIERE}{SVT}
\newcommand{\DOCNIVEAU}{1ère TI}
\newcommand{\DOCMODULE}{Module 2 : L'Éducation à la Santé}
\newcommand{\DOCLECON}{NCT06}
\newcommand{\DOCTITRE}{Lutte contre la mauvaise alimentation et l'inactivité physique}
% OnBuch+ — préambule « cours signés » ENRICHI (compilé avec tectonic / XeLaTeX)
% Système visuel « Pop » d'OnBuch+ : fond crème (jamais blanc), bordures encre
% épaisses, ombres « dures » décalées, titres Archivo Black en capitales.
% Version MATHÉMATIQUES : graphes (pgfplots/tikz), boîtes pédagogiques
% (définition, propriété, méthode, attention, le savais-tu, exercice résolu,
% exercice, TP, bilan), page de garde et signature OnBuch+ ancrée en bas.
%
% Macros attendues dans main.tex AVANT \input{preamble} :
%   \DOCMATIERE  \DOCNIVEAU  \DOCMODULE  \DOCLECON (numéro)  \DOCTITRE
\documentclass[11pt]{article}

\usepackage{fontspec}
\usepackage[french]{babel}
\usepackage[a4paper,margin=2cm,top=3.4cm,bottom=2.8cm,headheight=1.4cm,footskip=1.3cm]{geometry}
\usepackage{amsmath,amssymb,amsfonts}
\usepackage{mathtools} % \xmapsto, \coloneqq... souvent utilisés par les modèles
\usepackage{cancel} % simplifications barrées (bilans de forces)
% \abs{z} (module, valeur absolue) et \norm{u} (norme), utilisés par les modèles
\providecommand{\abs}{}\let\abs\relax\DeclarePairedDelimiter{\abs}{\lvert}{\rvert}
\providecommand{\norm}{}\let\norm\relax\DeclarePairedDelimiter{\norm}{\lVert}{\rVert}
\usepackage[table]{xcolor} % [table] : fournit aussi \rowcolors (alternance de lignes)
\usepackage{pagecolor}
\usepackage{tikz}
\usetikzlibrary{arrows.meta,positioning,calc,shapes.geometric,shapes.misc,shapes.arrows,decorations.pathmorphing,decorations.pathreplacing,decorations.markings,patterns,shadows,mindmap,trees,fit,backgrounds,matrix,intersections,angles,quotes}
\usepackage{pgfplots}
\usepackage[version=4]{mhchem} % équations nucléaires \ce{^{235}_{92}U}
\usepackage[european,siunitx]{circuitikz} % schémas électriques
\pgfplotsset{compat=1.18}
\usepgfplotslibrary{fillbetween,groupplots} % aires entre courbes (calcul intégral)
\usepackage{enumitem}
\usepackage{fancyhdr}
% [most] charge toutes les bibliothèques tcolorbox (listings, minted, raster,
% documentation...) et gonfle inutilement la mémoire TeX sur un document long
% (~300 boîtes) : on ne charge que ce qui est réellement utilisé.
\usepackage[skins,breakable]{tcolorbox}
\usepackage{graphicx}
\usepackage{colortbl}
\usepackage{booktabs}
\usepackage{tabularx}
\usepackage{diagbox} % case d'en-tête diagonale des tableaux à double entrée
% Colonnes extensibles centrée / gauche / droite (C, L, R), souvent utilisées par les modèles
\newcolumntype{C}{>{\centering\arraybackslash}X}
\newcolumntype{L}{>{\raggedright\arraybackslash}X}
\newcolumntype{R}{>{\raggedleft\arraybackslash}X}
\usepackage{array}
\usepackage{multicol}
\usepackage{siunitx}
% parse-numbers=false : accepte aussi \SI{2,5 \times 10^{-3}}{...}
\sisetup{locale=FR,output-decimal-marker={,},group-minimum-digits=4,parse-numbers=false}
\DeclareSIUnit\ohm{\ensuremath{\symup{\Omega}}} % U+2126 absent de la police
\DeclareSIUnit\division{div} % réglages d'oscilloscope (ms/div, V/div)
\DeclareSIUnit\tr{tr} % tours (vitesse de rotation, ex. \SI{1500}{\tr\per\minute})
\DeclareSIUnit\ch{ch} % cheval-vapeur (puissance moteur, unité usuelle hors SI)
\DeclareSIUnit\franccfa{FCFA} % franc CFA (coûts, contexte camerounais)
\DeclareSIUnit\dioptre{\ensuremath{\delta}} % vergence d'une lentille (dioptrie)
\usepackage{qrcode}
% \degree : commande fréquemment utilisée par les modèles pour un angle ou une
% température en dehors de \SI{}{} (non fournie par les paquets chargés).
\providecommand{\degree}{^{\circ}}
% \qed (fin de démonstration), utilisé par les modèles sans amsthm
\providecommand{\qed}{\hfill$\square$}
% \barre : double barre des valeurs interdites dans un tableau de variations (tkz-tab)
\providecommand{\barre}{\ensuremath{\|}}
% environnement proof (amsthm non chargé) utilisé par les modèles
\ifdefined\proof\else\newenvironment{proof}[1][Démonstration]{\par\noindent\textit{#1.}\ }{\qed\par}\fi
\usepackage{lastpage}
\usepackage{hyperref}
\hypersetup{hidelinks}

% ============================================================
% Palette « Pop » OnBuch+ — mêmes hex que la classe P (plus_ui.dart)
% ============================================================
\definecolor{poporange}{HTML}{F59321}
\definecolor{popdark}{HTML}{E07A0C}
\definecolor{popcream}{HTML}{FFF6E8}
\definecolor{popink}{HTML}{1C1714}
\definecolor{poppurple}{HTML}{7A5AE0}
\definecolor{popgreen}{HTML}{2E9E6A}
\definecolor{poppink}{HTML}{E0457B}
\definecolor{popblue}{HTML}{2F7DE1}
\definecolor{popgold}{HTML}{C98A12}
\definecolor{popmuted}{HTML}{8A7F76}
\definecolor{popline}{HTML}{EADFD2}
\definecolor{popgray}{HTML}{8A7F76}
\colorlet{popgrey}{popgray}
% Alias tolérants (le modèle nomme parfois les couleurs différemment)
\colorlet{popwhite}{white}
\colorlet{popblack}{popink}
\colorlet{popred}{poppink}
\colorlet{popyellow}{popgold}
% Teintes claires pour fonds de tableaux / zones
\colorlet{popblueL}{popblue!10!white}
\colorlet{poporangeL}{poporange!14!white}
\colorlet{popgreenL}{popgreen!12!white}
\colorlet{poppurpleL}{poppurple!12!white}
\colorlet{poppinkL}{poppink!12!white}
\colorlet{popgoldL}{popgold!14!white}

\pagecolor{popcream}

% ============================================================
% Typographie
% ============================================================
\defaultfontfeatures{Ligatures=TeX}
\setmainfont{PlusJakartaSans-Medium.ttf}[Path=../../../fonts/,BoldFont=PlusJakartaSans-Bold.ttf]
% Maths en sans-serif (Fira Math) avec chiffres et lettres droites de Plus
% Jakarta, pour que les formules se fondent dans le texte.
\usepackage[math-style=ISO,bold-style=ISO]{unicode-math}
\setmathfont{FiraMath-Regular.otf}[Path=../../../fonts/]
\setmathfont{PlusJakartaSans-Medium.ttf}[Path=../../../fonts/,range={up/num,up/{Latin,latin},\mathup}]
\usepackage{icomma} % virgule décimale sans espace en mode math
% Caractères mathématiques Unicode absents des polices de texte (Plus Jakarta,
% Archivo Black) : rendus par leur équivalent mathématique, en texte comme en
% formule — sinon ils s'affichent en carré vide (ex. « Arithmétique dans ℤ »).
\usepackage{newunicodechar}
\newunicodechar{ℝ}{\ensuremath{\mathbb{R}}}
\newunicodechar{ℤ}{\ensuremath{\mathbb{Z}}}
\newunicodechar{ℂ}{\ensuremath{\mathbb{C}}}
\newunicodechar{ℕ}{\ensuremath{\mathbb{N}}}
\newunicodechar{ℚ}{\ensuremath{\mathbb{Q}}}
\newunicodechar{𝔻}{\ensuremath{\mathbb{D}}}
\newunicodechar{α}{\ensuremath{\alpha}}
\newunicodechar{β}{\ensuremath{\beta}}
\newunicodechar{γ}{\ensuremath{\gamma}}
\newunicodechar{δ}{\ensuremath{\delta}}
\newunicodechar{ε}{\ensuremath{\varepsilon}}
\newunicodechar{θ}{\ensuremath{\theta}}
\newunicodechar{λ}{\ensuremath{\lambda}}
\newunicodechar{μ}{\ensuremath{\mu}}
\newunicodechar{σ}{\ensuremath{\sigma}}
\newunicodechar{τ}{\ensuremath{\tau}}
\newunicodechar{φ}{\ensuremath{\varphi}}
\newunicodechar{ψ}{\ensuremath{\psi}}
\newunicodechar{ω}{\ensuremath{\omega}}
\newunicodechar{Σ}{\ensuremath{\Sigma}}
% majuscules produites par \MakeUppercase dans les titres (α-aminés -> Α)
\newunicodechar{Α}{\ensuremath{\alpha}}
\newunicodechar{Β}{\ensuremath{\beta}}
\newunicodechar{Γ}{\ensuremath{\Gamma}}
\newunicodechar{Δ}{\ensuremath{\Delta}}
\newunicodechar{Ε}{\ensuremath{\varepsilon}}
\newunicodechar{Θ}{\ensuremath{\Theta}}
\newunicodechar{Λ}{\ensuremath{\Lambda}}
\newunicodechar{Μ}{\ensuremath{\mu}}
\newunicodechar{Τ}{\ensuremath{\tau}}
\newunicodechar{Φ}{\ensuremath{\Phi}}
\newunicodechar{Ψ}{\ensuremath{\Psi}}
\newunicodechar{Ω}{\ensuremath{\Omega}}
\newunicodechar{↦}{\ensuremath{\mapsto}}
\newunicodechar{∈}{\ensuremath{\in}}
\newunicodechar{∉}{\ensuremath{\notin}}
\newunicodechar{∧}{\ensuremath{\wedge}}
\newunicodechar{⇔}{\ensuremath{\Leftrightarrow}}
\newunicodechar{⟺}{\ensuremath{\Longleftrightarrow}}
\newunicodechar{⇒}{\ensuremath{\Rightarrow}}
\newunicodechar{⟹}{\ensuremath{\Longrightarrow}}
\newunicodechar{∘}{\ensuremath{\circ}}
\newunicodechar{∪}{\ensuremath{\cup}}
\newunicodechar{∩}{\ensuremath{\cap}}
\newunicodechar{≡}{\ensuremath{\equiv}}
\newunicodechar{⊂}{\ensuremath{\subset}}
\newunicodechar{⊥}{\ensuremath{\perp}}
\newunicodechar{∃}{\ensuremath{\exists}}
\newunicodechar{∀}{\ensuremath{\forall}}
\newunicodechar{∛}{\ensuremath{\sqrt[3]{\,}}}
\newunicodechar{∜}{\ensuremath{\sqrt[4]{\,}}}
\newunicodechar{⃗}{\ensuremath{{}^{\to}}}
\newunicodechar{ⁿ}{\ifmmode^{n}\else\textsuperscript{\ensuremath{n}}\fi}
\newunicodechar{ˣ}{\ifmmode^{x}\else\textsuperscript{\ensuremath{x}}\fi}
\newunicodechar{ᵇ}{\ifmmode^{b}\else\textsuperscript{\ensuremath{b}}\fi}
\newunicodechar{ʳ}{\ifmmode^{r}\else\textsuperscript{\ensuremath{r}}\fi}
\newunicodechar{ᵘ}{\ifmmode^{u}\else\textsuperscript{\ensuremath{u}}\fi}
\newunicodechar{ʸ}{\ifmmode^{y}\else\textsuperscript{\ensuremath{y}}\fi}
\newunicodechar{ᵃ}{\ifmmode^{a}\else\textsuperscript{\ensuremath{a}}\fi}
\newunicodechar{ᵉ}{\ifmmode^{e}\else\textsuperscript{\ensuremath{e}}\fi}
\newunicodechar{ᵏ}{\ifmmode^{k}\else\textsuperscript{\ensuremath{k}}\fi}
\newunicodechar{ᵖ}{\ifmmode^{p}\else\textsuperscript{\ensuremath{p}}\fi}
\newunicodechar{ᴺ}{\ifmmode^{N}\else\textsuperscript{\ensuremath{N}}\fi}
\newunicodechar{ᶜ}{\ifmmode^{c}\else\textsuperscript{\ensuremath{c}}\fi}
\newunicodechar{ʰ}{\ifmmode^{h}\else\textsuperscript{\ensuremath{h}}\fi}
\newunicodechar{ᵠ}{\ifmmode^{q}\else\textsuperscript{\ensuremath{q}}\fi}
\newunicodechar{ₙ}{\ifmmode_{n}\else\textsubscript{\ensuremath{n}}\fi}
\newunicodechar{ₐ}{\ifmmode_{a}\else\textsubscript{\ensuremath{a}}\fi}
\newunicodechar{ᵢ}{\ifmmode_{i}\else\textsubscript{\ensuremath{i}}\fi}
\newunicodechar{ᵣ}{\ifmmode_{r}\else\textsubscript{\ensuremath{r}}\fi}
\newunicodechar{ₖ}{\ifmmode_{k}\else\textsubscript{\ensuremath{k}}\fi}
\newunicodechar{ᵦ}{\ifmmode_{\beta}\else\textsubscript{\ensuremath{\beta}}\fi}
\newunicodechar{ₓ}{\ifmmode_{x}\else\textsubscript{\ensuremath{x}}\fi}
\newfontfamily\popdisplay{ArchivoBlack-Regular.ttf}[Path=../../../fonts/]
\newfontfamily\poplogo{SpaceGrotesk-Bold.ttf}[Path=../../../fonts/]
\newfontfamily\popsemi{PlusJakartaSans-SemiBold.ttf}[Path=../../../fonts/]
\newfontfamily\popxbold{PlusJakartaSans-ExtraBold.ttf}[Path=../../../fonts/]
\newfontfamily\popxboldls{PlusJakartaSans-ExtraBold.ttf}[Path=../../../fonts/,LetterSpace=3.0]

\setlist[enumerate]{leftmargin=1.7em,itemsep=5pt,topsep=4pt}
\setlist[itemize]{leftmargin=1.7em,itemsep=4pt,topsep=4pt,label={\color{poporange}\textbullet}}
\setlength{\parindent}{0pt}
\setlength{\parskip}{5pt}
\renewcommand{\arraystretch}{1.25}

% pgfplots : style Pop par défaut
\pgfplotsset{
  every axis/.append style={axis line style={popink,line width=1pt},tick style={popink},
    grid style={popline,line width=0.5pt},label style={font=\small},tick label style={font=\footnotesize}},
  popcurve/.style={poporange,line width=1.8pt,no markers,smooth},
}
% Même style disponible dans un \draw TikZ simple (hors pgfplots)
\tikzset{popcurve/.style={poporange,line width=1.8pt}}
\tikzset{popgrid/.style={popline,very thin}}
\colorlet{popcurve}{poporange} % le nom du style est parfois utilisé comme couleur

% ============================================================
% Pill / squiggle
% ============================================================
\newcommand{\poppill}[3]{%
  \tikz[baseline=-0.55ex]\node[draw=popink,line width=1.2pt,rounded corners=2.6pt,%
    fill=#2,inner xsep=6.5pt,inner ysep=3pt]{%
    \popxboldls\fontsize{7}{7}\selectfont\color{#3}\MakeUppercase{#1}};%
}
\newcommand{\popsquiggle}[1]{%
  \tikz[baseline=-0.4ex]{
    \draw[#1,line width=2.6pt,line cap=round]
      plot [smooth] coordinates {(0,0.09) (0.34,0.01) (0.68,0.21) (1.02,0.01) (1.36,0.10)};
  }%
}
% Mot-clé surligné (marqueur orange sous le texte)
\newcommand{\cle}[1]{{\popxbold\color{popdark}#1}}

% ============================================================
% En-tête / pied
% ============================================================
\pagestyle{fancy}
\fancyhf{}
\renewcommand{\headrulewidth}{0pt}
\renewcommand{\footrulewidth}{0pt}
\lhead{\raisebox{-0.15ex}{\poplogo\fontsize{13}{13}\selectfont\color{popink}OnBuch\color{poporange}+}}
\rhead{\poppill{\DOCMATIERE\ \textbullet\ \DOCNIVEAU\ \textbullet\ Leçon \DOCLECON}{white}{popink}}
\lfoot{\popxbold\fontsize{7.6}{7.6}\selectfont\color{popmuted}\MakeUppercase{Cours signé OnBuch+}\ \textbullet\ {\popsemi\DOCTITRE}}
\rfoot{\tikz[baseline=-0.6ex]\node[rounded corners=8.5pt,fill=popink,minimum height=17pt,minimum width=17pt,inner xsep=5pt,inner sep=0]{\popxbold\fontsize{8}{8}\selectfont\color{popcream}\thepage\,/\,\pageref*{LastPage}};}

% ============================================================
% Titres
% ============================================================
\newcommand{\chaptitre}[2]{%
  \begin{tcolorbox}[enhanced,colback=poporange,colframe=popink,boxrule=2.6pt,arc=11pt,
      drop shadow=popink,left=15pt,right=15pt,top=12pt,bottom=11pt,before skip=3pt,after skip=18pt]
    \poppill{#2}{popink}{popcream}\par\vspace{9pt}
    {\raggedright\hyphenpenalty=10000\exhyphenpenalty=10000\popdisplay\fontsize{22}{24}\selectfont\color{popcream}\MakeUppercase{#1}\par}
  \end{tcolorbox}
}
% Section numérotée : pastille orange avec le numéro + titre Archivo
\newcounter{popsec}
\newcommand{\coursec}[1]{%
  \stepcounter{popsec}\setcounter{popsub}{0}%
  \par\vspace{16pt}\noindent
  \begin{minipage}{\linewidth}
  \tikz[baseline=-0.7ex]\node[draw=popink,line width=1.6pt,fill=poporange,rounded corners=4pt,minimum size=22pt,inner sep=0]{\popdisplay\fontsize{12}{12}\selectfont\color{popcream}\thepopsec};%
  \hspace{8pt}{\popdisplay\fontsize{15}{17}\selectfont\color{popink}\MakeUppercase{#1}}\par
  \vspace{4pt}\noindent{\color{poporange}\rule{36pt}{3.6pt}}
  \end{minipage}\par\nopagebreak\vspace{8pt}
}
\newcounter{popsub}
\newcommand{\courssub}[1]{%
  \stepcounter{popsub}%
  \par\vspace{9pt}\noindent{\popxbold\fontsize{11.5}{13}\selectfont\color{popink}{\color{poporange}\thepopsec.\thepopsub}\hspace{6pt}#1}\par\nopagebreak\vspace{4pt}
}

% ============================================================
% Boîtes pédagogiques — même recette : fond blanc, bordure épaisse colorée,
% ombre dure encre, badge plein. Titre optionnel [#1].
% ============================================================
\tcbset{popbase/.style={breakable,enhanced,colback=white,boxrule=2.3pt,arc=9pt,
  shadow={3pt}{-3pt}{0pt}{popink},left=14pt,right=14pt,top=11pt,bottom=11pt,before skip=11pt,after skip=15pt,
  fontupper=\popsemi\fontsize{10.6}{14.5}\selectfont\color{popink},
  fontlower=\popsemi\fontsize{10.6}{14.5}\selectfont\color{popink}}}
% Coche dessinée (le glyphe ✓ n'existe pas dans les polices du document)
\newcommand{\popcheck}[1]{\tikz[baseline=-0.5ex]\draw[#1,line width=1.6pt,line cap=round,line join=round](-0.13,0.0)--(-0.04,-0.09)--(0.13,0.1);}
\providecommand{\coche}{\popcheck{popgreen}}
\providecommand{\resultat}[1]{\par\smallskip\noindent\fcolorbox{popgreen}{popgreenL}{\parbox{\dimexpr\linewidth-2\fboxsep-2\fboxrule\relax}{\textbf{Résultat.} #1}}\par\smallskip}
\newcommand{\popboxhead}[3]{\poppill{#1}{#2}{white}\ifx\relax#3\relax\else\hspace{6pt}{\popxbold\fontsize{10.6}{12}\selectfont\color{popink}#3}\fi\par\vspace{7pt}}

\newtcolorbox{aretenir}[1][]{popbase,colframe=popgreen,before upper={\popboxhead{À retenir}{popgreen}{#1}}}
\newtcolorbox{exemplebox}[1][]{popbase,colframe=poporange,before upper={\popboxhead{Exemple}{poporange}{#1}}}
\newtcolorbox{definition}[1][]{popbase,colframe=popblue,before upper={\popboxhead{Définition}{popblue}{#1}}}
\newtcolorbox{propriete}[1][]{popbase,colframe=poppurple,before upper={\popboxhead{Propriété}{poppurple}{#1}}}
\newtcolorbox{methode}[1][]{popbase,colframe=popgold,before upper={\popboxhead{Méthode}{popgold}{#1}}}
\newtcolorbox{attention}[1][]{popbase,colframe=poppink,before upper={\popboxhead{Attention}{poppink}{#1}}}
\newtcolorbox{experience}[1][]{popbase,colframe=popblue,colback=popblueL,before upper={\popboxhead{Expérience}{popblue}{#1}}}
% « Le savais-tu ? » : carte violette pleine (contexte, culture, Cameroun)
\newtcolorbox{savaistu}[1][]{popbase,colback=poppurple,colframe=popink,
  fontupper=\popsemi\fontsize{10.4}{14.2}\selectfont\color{white},
  before upper={\poppill{Le savais-tu\,?}{white}{poppurple}\ifx\relax#1\relax\else\hspace{6pt}{\popxbold\color{white}#1}\fi\par\vspace{7pt}}}
% Exercice résolu : énoncé en haut, \tcblower puis la solution sur fond crème
\newcounter{popexo}\newcounter{popexr}
\newtcolorbox{exoresolu}[1][]{popbase,colframe=poporange,colbacklower=popcream,
  segmentation style={popink,line width=1.2pt,dashed},
  before upper={\stepcounter{popexr}\popboxhead{Exercice résolu \thepopexr}{poporange}{#1}},
  before lower={\poppill{Solution}{popink}{popcream}\par\vspace{6pt}}}
% Exercice (non résolu en place) — la correction est dans la partie Corrigés
\newtcolorbox{exercice}[1][]{popbase,colframe=popink,
  before upper={\stepcounter{popexo}\popboxhead{Exercice \thepopexo}{popink}{#1}}}
% Corrigé
\newtcolorbox{corrige}[1][]{popbase,colframe=popgreen,colback=popgreenL,
  before upper={\popboxhead{Corrigé}{popgreen}{#1}}}
% Objectifs / prérequis / bilan
\newtcolorbox{objectifs}{popbase,colframe=popink,colback=poporangeL,
  before upper={\popboxhead{Objectifs de la leçon}{popink}{}}}
\newtcolorbox{prerequis}{popbase,colframe=popink,colback=popblueL,
  before upper={\popboxhead{Prérequis}{popblue}{}}}
\newtcolorbox{bilan}{popbase,colframe=popink,colback=white,boxrule=2.8pt,
  before upper={\popboxhead{Fiche bilan}{poporange}{L'essentiel à connaître pour l'examen}}}
% Figure encadrée avec légende
\newtcolorbox{popfigure}[1][]{enhanced,breakable=false,colback=white,colframe=popline,boxrule=1.4pt,arc=8pt,
  left=10pt,right=10pt,top=10pt,bottom=8pt,before skip=10pt,after skip=12pt,halign=center,
  lower separated=false,halign lower=center,
  fontlower=\popsemi\fontsize{9}{11.5}\selectfont\color{popmuted},
  #1}
\newcommand{\legende}[1]{\par\vspace{6pt}{\popsemi\fontsize{9}{11.5}\selectfont\color{popmuted}#1}}

% ============================================================
% Signature OnBuch+
% ============================================================
\newcommand{\poplogoinline}[1]{{\poplogo\fontsize{#1}{#1}\selectfont\color{popink}OnBuch\color{poporange}+}}
\newcommand{\signatureonbuch}{%
  \begin{tcolorbox}[enhanced,colback=popink,colframe=popink,boxrule=2.2pt,arc=10pt,
      drop shadow=poporange,left=14pt,right=14pt,top=10pt,bottom=10pt,before skip=0pt,after skip=0pt]
    \tikz[baseline=-0.7ex]\node[circle,fill=popgreen,draw=popcream,line width=1.3pt,minimum size=22pt,inner sep=0]{\popcheck{white}};%
    \hspace{9pt}\begin{minipage}[c]{0.62\linewidth}
      {\popxbold\fontsize{12}{13}\selectfont\color{popcream}Signé\ \textbullet\ {\poplogo OnBuch\color{poporange}+}}\par\vspace{2pt}
      {\raggedright\popsemi\fontsize{8}{10.5}\selectfont\color{popline}\MakeUppercase{Équipe pédagogique OnBuch+}\\\MakeUppercase{Conforme au programme officiel MINESEC}\par}
    \end{minipage}\hfill
    {\poplogo\fontsize{22}{22}\selectfont\color{popcream}OnBuch\color{poporange}+}
  \end{tcolorbox}%
}

% ============================================================
% Page de garde
% #1 titre · #2 matière · #3 niveau · #4 module · #5 n° leçon · #6 durée ·
% #7 description / objectifs (texte ou itemize)
% ============================================================
\newcommand{\pagegarde}[7]{%
  \thispagestyle{empty}
  \newgeometry{a4paper,margin=1.7cm,top=1.6cm,bottom=1.4cm}%
  \begin{tikzpicture}[remember picture,overlay]
    \fill[poporange!18] ([xshift=-3.2cm,yshift=-3.6cm]current page.north east) circle (4.6cm);
    \fill[poppurple!14] ([xshift=2.4cm,yshift=7.6cm]current page.south west) circle (3.8cm);
  \end{tikzpicture}%
  {\poplogo\fontsize{30}{30}\selectfont\color{popink}OnBuch\color{poporange}+}\hfill
  \raisebox{0.6ex}{\poppill{Cours signé \textbullet\ Premium}{popink}{popcream}}\par
  \vspace{1pt}\popsquiggle{poppurple}\par
  \vspace{8pt}
  \begin{tcolorbox}[enhanced,colback=poporange,colframe=popink,boxrule=2.8pt,arc=13pt,
      drop shadow=popink,left=18pt,right=18pt,top=12pt,bottom=13pt,after skip=10pt]
    \poppill{#2 \textbullet\ #3}{popink}{popcream}\hspace{5pt}\poppill{#4}{white}{popink}\par\vspace{8pt}
    {\popdisplay\fontsize{14}{14}\selectfont\color{popink}LEÇON #5}\par\vspace{4pt}
    {\raggedright\hyphenpenalty=10000\exhyphenpenalty=10000\popdisplay\fontsize{25}{27}\selectfont\color{popcream}\MakeUppercase{#1}\par}
    \vspace{8pt}
    {\popxbold\fontsize{9.5}{9.5}\selectfont\color{popink}\MakeUppercase{Durée conseillée : #6}}
  \end{tcolorbox}
  \begin{tcolorbox}[enhanced,colback=white,colframe=popink,boxrule=2.2pt,arc=10pt,
      drop shadow=popink,left=16pt,right=16pt,top=10pt,bottom=10pt,after skip=8pt]
    \poppill{Dans cette leçon}{poppurple}{white}\par\vspace{5pt}
    \popsemi\fontsize{9.3}{12.6}\selectfont\color{popink}#7
  \end{tcolorbox}
  \noindent\begin{minipage}[t]{0.487\linewidth}
    \begin{tcolorbox}[enhanced,colback=popgreen,colframe=popink,boxrule=2.2pt,arc=10pt,
        drop shadow=popink,left=11pt,right=11pt,top=10pt,bottom=10pt,height=3.1cm,valign=center]
      \tcbox[colback=white,colframe=popink,boxrule=1.3pt,arc=5pt,boxsep=0pt,nobeforeafter,
        left=3pt,right=3pt,top=3pt,bottom=3pt,valign=center]{\qrcode[height=1.9cm]{https://play.google.com/store/apps/details?id=com.luvvix.onbuch&hl=fr}}%
      \hspace{8pt}\begin{minipage}[c]{0.5\linewidth}
        {\popdisplay\fontsize{11}{12.5}\selectfont\color{white}\MakeUppercase{Télécharge OnBuch}}\par\vspace{4pt}
        {\raggedright\popsemi\fontsize{8.4}{11}\selectfont\color{white}Gratuit \textbullet\ Google Play\\Tuteur IA, annales, concours, résultats\par}
      \end{minipage}
    \end{tcolorbox}
  \end{minipage}\hfill
  \begin{minipage}[t]{0.487\linewidth}
    \begin{tcolorbox}[enhanced,colback=poppurple,colframe=popink,boxrule=2.2pt,arc=10pt,
        drop shadow=popink,left=11pt,right=11pt,top=10pt,bottom=10pt,height=3.1cm,valign=center]
      \tikz[baseline=-0.7ex]\node[circle,fill=white,draw=popink,line width=1.3pt,minimum size=1.6cm,inner sep=0]{\popdisplay\fontsize{24}{24}\selectfont\color{poppurple}+};%
      \hspace{8pt}\begin{minipage}[c]{0.6\linewidth}
        {\popdisplay\fontsize{11}{12.5}\selectfont\color{white}\MakeUppercase{Passe à OnBuch+}}\par\vspace{4pt}
        {\raggedright\popsemi\fontsize{8.4}{11}\selectfont\color{white}Tous les cours signés, Tuteur IA illimité, fiches PDF — onglet OnBuch+ de l'app\par}
      \end{minipage}
    \end{tcolorbox}
  \end{minipage}
  \vspace{8pt}
  \popcontenu
  \vfill
  \signatureonbuch
  \restoregeometry
  \newpage
}

% Rangée « Ce cours contient » (page de garde)
\newcommand{\poptile}[3]{% #1 couleur #2 gros texte #3 légende
  \begin{tcolorbox}[enhanced,colback=#1,colframe=popink,boxrule=2pt,arc=9pt,shadow={2.5pt}{-2.5pt}{0pt}{popink},
      left=6pt,right=6pt,top=9pt,bottom=9pt,halign=center,valign=center,height=1.75cm,width=\linewidth,nobeforeafter]
    {\popdisplay\fontsize{12.5}{13}\selectfont\color{white}\MakeUppercase{#2}}\par\vspace{3pt}
    {\centering\popsemi\fontsize{7.8}{9.5}\selectfont\color{white}#3\par}
  \end{tcolorbox}}
\newcommand{\popcontenu}{%
  \par\noindent{\popxboldls\fontsize{7.5}{7.5}\selectfont\color{popmuted}CE COURS SIGNÉ CONTIENT}\par\vspace{7pt}
  \noindent\begin{minipage}[t]{0.235\linewidth}\poptile{popblue}{Cours}{complet, illustré, pas à pas}\end{minipage}\hfill
  \begin{minipage}[t]{0.235\linewidth}\poptile{poporange}{Méthodes}{et exercices résolus}\end{minipage}\hfill
  \begin{minipage}[t]{0.235\linewidth}\poptile{poppink}{Exercices}{type examen + corrigés}\end{minipage}\hfill
  \begin{minipage}[t]{0.235\linewidth}\poptile{popgreen}{Fiche bilan}{l'essentiel à retenir}\end{minipage}\par}

% Sommaire
\newtcolorbox{sommaire}{enhanced,colback=white,colframe=popink,boxrule=2.2pt,arc=10pt,
  drop shadow=popink,left=18pt,right=18pt,top=13pt,bottom=13pt,before skip=6pt,after skip=20pt,
  before upper={\poppill{Sommaire}{poppurple}{white}\par\vspace{9pt}\popsemi\fontsize{10.6}{14.5}\selectfont\color{popink}}}

% Signature de fin de cours — poussée tout en bas de la dernière page
\newcommand{\findecours}{%
  \par\vfill\nopagebreak
  \begin{center}\popsquiggle{poporange}\end{center}\vspace{4pt}
  \signatureonbuch
}
\AtBeginDocument{\def\llbracket{\mathopen{[\mkern-3.5mu[}}\def\rrbracket{\mathclose{]\mkern-3.5mu]}}} % intervalles d'entiers (glyphe ⟦ absent de la police)
% Glyphes absents de Fira Math : redéfinis à partir de glyphes disponibles
\AtBeginDocument{%
  \def\setminus{\mathbin{\backslash}}\let\smallsetminus\setminus
  \def\vdots{\mathord{\vbox{\baselineskip3.2pt\lineskiplimit0pt\kern4pt\hbox{.}\hbox{.}\hbox{.}}}}%
  \def\ddots{\mathinner{\mkern1mu\raise6.5pt\hbox{.}\mkern2mu\raise3.6pt\hbox{.}\mkern2mu\raise0.7pt\hbox{.}\mkern1mu}}%
  \def\triangle{\mathord{\scalebox{0.9}{$\symup{\Delta}$}}}%
  \def\nabla{\mathord{\raisebox{\depth}{\scalebox{1}[-1]{$\symup{\Delta}$}}}}%
  \def\checkmark{\popcheck{popdark}}%
  \def\star{\mathbin{\ast}}\def\bigstar{\mathord{\ast}}%
  \def\diamond{\mathbin{\scalebox{0.6}{$\Diamond$}}}%
  \def\bigcup{\mathop{\mathchoice{\raisebox{-0.3em}{\scalebox{1.7}{$\cup$}}}{\scalebox{1.25}{$\cup$}}{\cup}{\cup}}\displaylimits}%
  \def\bigcap{\mathop{\mathchoice{\raisebox{-0.3em}{\scalebox{1.7}{$\cap$}}}{\scalebox{1.25}{$\cap$}}{\cap}{\cap}}\displaylimits}%
  \def\lesseqqgtr{\mathrel{\lessgtr}}\def\bigoplus{\mathop{\oplus}\displaylimits}%
}
\providecommand{\ssi}{si et seulement si} % abréviation fréquente
\AtBeginDocument{\providecommand{\wideparen}{\overparen}} % arc de cercle

%% FIGFIX-BEGIN v5 — correctifs globaux des figures (docs/onbuch-generation/tools/figfix)
\usepackage{adjustbox}\usepackage{etoolbox}\tcbuselibrary{raster}
% toute figure TikZ/pgfplots plus large que la ligne est réduite à la largeur disponible
\BeforeBeginEnvironment{tikzpicture}{\begin{adjustbox}{max width=\linewidth}}
\AfterEndEnvironment{tikzpicture}{\end{adjustbox}}
% légendes pgfplots lisibles par-dessus les courbes
\pgfplotsset{every axis/.append style={legend style={fill=white,fill opacity=0.92,text opacity=1,draw=black!15,inner sep=3pt}}}
% étiquettes d'arêtes (syntaxe edge["..."]) avec fond blanc
\tikzset{every edge quotes/.append style={fill=white,inner sep=1.2pt}}
%% FIGFIX-END
\begin{document}

\pagegarde{Lutte contre la mauvaise alimentation et l'inactivité physique}{SVT}{1ère TI}{Module 2 : L'Éducation à la Santé}{NCT06}{4 h}{Cette leçon étudie l'équilibre alimentaire et les besoins nutritionnels de l'adolescent, les conséquences sanitaires de la sédentarité et de la malbouffe (obésité, diabète de type 2, maladies cardiovasculaires), ainsi que les recommandations officielles en matière d'alimentation et d'activité physique. Elle forme l'élève à évaluer une ration alimentaire, à calculer et interpréter l'IMC, et à élaborer un programme simple de prévention.
\vspace{4pt}
\begin{multicols}{2}\small
\begin{itemize}
\item À la fin de cette leçon, je sais définir l'équilibre alimentaire et citer les grandes familles d'aliments avec leurs rôles.
\item À la fin de cette leçon, je sais décrire les besoins nutritionnels spécifiques de l'adolescent (énergie, protéines, fer, calcium).
\item À la fin de cette leçon, je sais évaluer une ration alimentaire quotidienne par rapport aux apports recommandés.
\item À la fin de cette leçon, je sais calculer l'IMC d'une personne et l'interpréter à l'aide des seuils de l'OMS.
\item À la fin de cette leçon, je sais relier les valeurs d'IMC aux risques sanitaires (obésité, diabète de type 2, maladies cardiovasculaires).
\item À la fin de cette leçon, je sais expliquer les mécanismes par lesquels la sédentarité et la mauvaise alimentation provoquent ces maladies.
\end{itemize}\end{multicols}}

\begin{sommaire}
\begin{multicols}{2}
\begin{enumerate}
\item Équilibre alimentaire et besoins de l'adolescent
\item Évaluation d'une ration alimentaire
\item L'IMC : calcul et interprétation
\item Conséquences de la sédentarité et de la mauvaise alimentation
\item Recommandations : alimentation et activité physique
\item Élaborer un programme de prévention
\item Activité d'intégration
\item Méthodes et exercices résolus
\item Exercices
\item Corrigés des exercices
\item Fiche bilan
\end{enumerate}
\end{multicols}
\end{sommaire}

\begin{objectifs}
\begin{itemize}
\item À la fin de cette leçon, je sais définir l'équilibre alimentaire et citer les grandes familles d'aliments avec leurs rôles.
\item À la fin de cette leçon, je sais décrire les besoins nutritionnels spécifiques de l'adolescent (énergie, protéines, fer, calcium).
\item À la fin de cette leçon, je sais évaluer une ration alimentaire quotidienne par rapport aux apports recommandés.
\item À la fin de cette leçon, je sais calculer l'IMC d'une personne et l'interpréter à l'aide des seuils de l'OMS.
\item À la fin de cette leçon, je sais relier les valeurs d'IMC aux risques sanitaires (obésité, diabète de type 2, maladies cardiovasculaires).
\item À la fin de cette leçon, je sais expliquer les mécanismes par lesquels la sédentarité et la mauvaise alimentation provoquent ces maladies.
\item À la fin de cette leçon, je sais citer les recommandations officielles en matière d'activité physique et d'alimentation.
\item À la fin de cette leçon, je sais élaborer un programme hebdomadaire simple combinant alimentation équilibrée et activité physique.
\end{itemize}
\end{objectifs}

\begin{prerequis}
\begin{itemize}
\item Notions de nutrition vues au collège : nutriments (glucides, lipides, protéines, vitamines, sels minéraux, eau) et leurs rôles.
\item Notion d'énergie et unités (joule, kilocalorie) vues en Physique.
\item Calculs élémentaires : quotient, pourcentage, lecture de tableaux et de diagrammes.
\item Notions sur le système cardiovasculaire et la digestion (SVT de Seconde).
\item Hygiène de vie et notions de prévention abordées dans les leçons précédentes du Module 2.
\end{itemize}
\end{prerequis}

\newpage

\coursec{Équilibre alimentaire et besoins de l'adolescent}

\textbf{Situation d'accroche.} Il est 6 h 30 à Yaoundé, quartier Mokolo. Comme beaucoup de ses camarades, Junior, 16 ans, achète au coin de la rue son petit-déjeuner : des beignets-haricots bien gras, avalés en vitesse. À midi, au réfectoire improvisé près du lycée, c'est un plat de friture de pommes de terre arrosé d'une boisson sucrée bien fraîche. L'après-midi et le soir ? Assis, le nez sur son téléphone ou devant la télévision. Lors de la visite médicale scolaire, le médecin constate que Junior — comme un nombre croissant d'élèves de son âge — présente déjà un \cle{surpoids} et se plaint d'une \cle{fatigue chronique}. Pourtant, Junior « mange bien », affirme-t-il : il mange beaucoup, et gras !

Cette situation, familière à tant de lycéens camerounais, pose des questions essentielles : \emph{comment savoir si notre alimentation quotidienne couvre bien nos besoins ? Que signifie réellement « bien manger » ?} C'est à ces questions que répond la notion d'\cle{équilibre alimentaire}, fondation de toute cette leçon.

\courssub{Qu'est-ce que l'équilibre alimentaire ?}

\textbf{Observation et intuition.} Une moto a besoin d'essence pour rouler, mais aussi d'huile pour le moteur et d'un mécanicien pour l'entretenir. Si l'on met uniquement de l'essence, en grande quantité, la moto finira par tomber en panne. Notre corps fonctionne de la même manière : il ne suffit pas de le « remplir », il faut lui fournir \emph{les bons carburants, dans les bonnes proportions}.

\begin{definition}[Équilibre alimentaire]
L'\cle{équilibre alimentaire} est la correspondance entre les \textbf{apports alimentaires} (ce que l'on mange et boit) et les \textbf{besoins de l'organisme} (ce dont le corps a besoin), à la fois :
\begin{itemize}
\item en \textbf{quantité} : l'énergie apportée doit couvrir les dépenses, sans excès ni déficit ;
\item en \textbf{qualité} : tous les nutriments indispensables (glucides, lipides, protéines, vitamines, minéraux, eau) doivent être présents dans des \textbf{proportions adaptées} ;
\item dans le \textbf{temps} : les apports doivent être \textbf{répartis sur la journée} (idéalement 3 repas, éventuellement une collation).
\end{itemize}
Un régime peut être abondant et pourtant déséquilibré : c'est le cas de Junior, qui mange beaucoup mais presque uniquement des aliments gras et sucrés.
\end{definition}

\begin{attention}[Erreur fréquente : « manger beaucoup » $\neq$ « bien manger »]
Dans beaucoup de familles, « bien manger » évoque un plat copieux et gras. C'est une idée reçue dangereuse. L'équilibre alimentaire ne se mesure ni à la quantité ni à la richesse en graisse d'un plat, mais à la \textbf{variété} des aliments consommés et à leurs \textbf{proportions}. Un énorme plat de friture apporte un excès de lipides et d'énergie, mais presque pas de vitamines, de calcium ou de protéines de qualité : c'est une alimentation déséquilibrée, même si le ventre est plein.
\end{attention}

\courssub{Les sept familles d'aliments et leurs rôles}

\textbf{Observation.} Au marché de Mokolo, on trouve des centaines d'aliments différents. Pourtant, les nutritionnistes les rangent en seulement \textbf{sept familles}, chacune ayant un rôle dominant dans l'organisme. Aucune famille ne peut être supprimée sans conséquence : c'est la complémentarité qui fait l'équilibre.

\begin{center}
\begin{tabularx}{\linewidth}{|l|X|X|}
\hline
\rowcolor{popblueL} \textbf{Famille d'aliments} & \textbf{Rôle principal} & \textbf{Exemples camerounais} \\
\hline
1. Féculents & \textbf{Énergétique} : sucres lents (glucides complexes), carburant durable des muscles et du cerveau & Maïs, manioc (bâton, gari), igname, taro, plantain, riz, patate douce \\
\hline
2. Viandes, poissons, œufs, légumineuses & \textbf{Bâtisseur} : protéines pour construire et réparer les tissus (muscles, peau, sang) ; fer & Poisson fumé, viande de bœuf, poulet bicyclette, œufs, haricots niébé, arachides, soja \\
\hline
3. Lait et produits laitiers & \textbf{Bâtisseur} : protéines et surtout \textbf{calcium} pour l'ossification & Lait caillé, yaourt, fromage, lait en poudre \\
\hline
4. Fruits et légumes & \textbf{Protecteur} : vitamines, minéraux, fibres ; défenses immunitaires & Ndolé, folong, eru, gombo, carotte, mangue, banane douce, orange, goyave, papaye \\
\hline
5. Matières grasses & \textbf{Énergétique} concentrée ; transport des vitamines A, D, E, K ; constituants des membranes cellulaires & Huile de palme, huile d'arachide, beurre de karité, avocat \\
\hline
6. Produits sucrés & \textbf{Énergétique} rapide mais « calories vides » : à limiter fortement & Sucre, boissons sucrées, bonbons, beignets, gâteaux \\
\hline
7. Eau & \textbf{Indispensable} : milieu de toutes les réactions chimiques du corps, transport, régulation de la température & Eau de boisson (1,5 à 2 L/jour) \\
\hline
\end{tabularx}
\end{center}

\begin{aretenir}[Les trois grands rôles des aliments]
\begin{itemize}
\item Rôle \cle{énergétique} : fournir les calories (féculents, matières grasses, produits sucrés) ;
\item Rôle \cle{bâtisseur} (ou plastique) : construire et réparer le corps (viandes-poissons-œufs-légumineuses, laitages) ;
\item Rôle \cle{protecteur} (ou fonctionnel) : assurer le bon fonctionnement et les défenses (fruits, légumes, eau).
\end{itemize}
Une alimentation équilibrée puise \textbf{chaque jour} dans chacune des sept familles, en privilégiant les quatre premières.
\end{aretenir}

\courssub{Les besoins énergétiques de l'adolescent}

\textbf{Observation et démarche.} Pourquoi un adolescent de 16 ans mange-t-il souvent plus que son père de 45 ans ? \emph{Hypothèse} : la croissance rapide de l'adolescence exige des apports supplémentaires. \emph{Données} : entre 12 et 18 ans, un garçon peut grandir de 8 à 10 cm par an et prendre 5 à 8 kg de masse musculaire et osseuse par an ; cette construction corporelle massive consomme de l'énergie et des matériaux. \emph{Conclusion} : les besoins de l'adolescent sont parmi les plus élevés de toute la vie.

\begin{propriete}[Besoins énergétiques de l'adolescent]
Les besoins énergétiques quotidiens d'un adolescent varient approximativement entre \textbf{2 200 et 3 000 kcal/jour}, selon :
\begin{itemize}
\item le \textbf{sexe} : les garçons ont en moyenne des besoins supérieurs (environ 2 500 à 3 000 kcal/j) aux filles (environ 2 200 à 2 400 kcal/j), en raison d'une masse musculaire plus importante ;
\item le \textbf{niveau d'activité physique} : un élève sportif ou qui marche plusieurs kilomètres par jour dépense bien plus qu'un élève sédentaire ;
\item la \textbf{vitesse de croissance} : en pleine poussée de croissance, les besoins sont maximaux.
\end{itemize}
\end{propriete}

\begin{exemplebox}[Exemple chiffré]
Un adolescent de 16 ans très actif (sport régulier, déplacements à pied) a besoin d'environ \textbf{2 800 kcal/jour}, soit environ \textbf{90 g de protéines} et \textbf{1 300 mg de calcium} par jour — autant de calcium que dans un litre de lait ! Ces chiffres montrent que l'adolescence est une période où aucune famille d'aliments ne peut être négligée.
\end{exemplebox}

\courssub{Les besoins spécifiques de l'adolescent}

Au-delà de l'énergie, certains nutriments sont \emph{critiques} pendant l'adolescence :

\begin{itemize}
\item \textbf{Les protéines} (viandes, poissons, œufs, haricots, arachides) : elles fournissent les \cle{acides aminés}, briques de la croissance musculaire et du renouvellement des tissus. Un déficit freine la croissance et affaiblit les défenses immunitaires.
\item \textbf{Le calcium} (lait et produits laitiers, petits poissons consommés avec les arêtes, certains légumes verts) : c'est le minéral de l'\cle{ossification}. Près de la moitié de la masse osseuse adulte se construit pendant l'adolescence ; un déficit à cet âge fragilise le squelette pour toute la vie.
\item \textbf{Le fer} (viande, poisson, ndolé, folong, haricots) : il entre dans la composition de l'\cle{hémoglobine} qui transporte l'oxygène dans le sang. Les \textbf{adolescentes} ont des besoins accrus à cause des pertes de sang liées aux règles ; le déficit en fer provoque l'\cle{anémie} (fatigue, pâleur, difficultés de concentration — peut-être une des causes de la fatigue de Junior !).
\item \textbf{Les vitamines} (fruits et légumes) : vitamine C (défenses, absorption du fer), vitamine A (vision, peau), vitamines du groupe B (utilisation de l'énergie), etc.
\end{itemize}

\begin{savaistu}[Un truc de grand-mère... validé par la science]
Le fer des légumes verts comme le ndolé ou le folong est du fer « non héminique », moins bien absorbé que le fer de la viande. Mais il est \textbf{beaucoup mieux absorbé} lorsqu'on accompagne le repas d'un agrume riche en \textbf{vitamine C} : une orange ou un jus de goyave après le plat de ndolé peut multiplier par deux ou trois l'absorption du fer ! Voilà une habitude simple et gratuite pour lutter contre l'anémie, particulièrement chez les adolescentes.
\end{savaistu}

\courssub{La répartition conseillée de la ration quotidienne}

Manger équilibré, ce n'est pas seulement choisir les bonnes familles : c'est aussi respecter des \textbf{proportions} entre les trois grands nutriments énergétiques.

\begin{propriete}[Répartition énergétique recommandée]
Sur l'énergie totale apportée chaque jour, la répartition recommandée est :
\begin{itemize}
\item \textbf{Glucides : 50 à 55 \%} de l'énergie totale (principalement des sucres lents : féculents) ;
\item \textbf{Lipides : 30 à 35 \%} (en limitant les graisses de friture) ;
\item \textbf{Protéines : 12 à 15 \%}.
\end{itemize}
\end{propriete}

\begin{popfigure}
\begin{center}
\begin{tikzpicture}[scale=1.05]
% Secteurs : protéines 13.5% (48.6°), lipides 32.5% (117°), glucides 54% (194.4°)
\fill[popgreen!70] (0,0) -- (90:2.6) arc (90:138.6:2.6) -- cycle;
\fill[poporange!75] (0,0) -- (138.6:2.6) arc (138.6:255.6:2.6) -- cycle;
\fill[popblue!70] (0,0) -- (255.6:2.6) arc (255.6:450:2.6) -- cycle;
\draw[popdark, thick] (0,0) circle (2.6);
% Étiquettes
\node[white, font=\bfseries] at (114:1.6) {Protéines};
\node[white, font=\bfseries] at (114:1.15) {12--15 \%};
\node[white, font=\bfseries] at (197:1.6) {Lipides};
\node[white, font=\bfseries] at (197:1.15) {30--35 \%};
\node[white, font=\bfseries] at (-8:1.7) {Glucides};
\node[white, font=\bfseries] at (-8:1.25) {50--55 \%};
% Légende latérale
\fill[popblue!70] (3.4,1.6) rectangle (3.8,2.0); \node[right, font=\small] at (3.9,1.8) {Glucides : féculents (maïs, manioc, igname)};
\fill[poporange!75] (3.4,0.7) rectangle (3.8,1.1); \node[right, font=\small] at (3.9,0.9) {Lipides : huiles, avocat, arachide};
\fill[popgreen!70] (3.4,-0.2) rectangle (3.8,0.2); \node[right, font=\small] at (3.9,0.0) {Protéines : viande, poisson, œufs, haricots};
\end{tikzpicture}
\end{center}
\legende{Diagramme circulaire de la répartition énergétique recommandée pour un adolescent : les glucides (sucres lents surtout) doivent fournir environ la moitié de l'énergie, les lipides environ un tiers, et les protéines 12 à 15 \%.}
\end{popfigure}

\begin{methode}[Comment évaluer une ration alimentaire]
Pour juger si une ration quotidienne est équilibrée, on procède en trois étapes :
\begin{enumerate}
\item \textbf{Calculer l'apport énergétique total} de la journée (somme des calories de tous les repas et boissons) et le comparer au besoin (2 200 à 3 000 kcal/j selon le sexe et l'activité) ;
\item \textbf{Calculer la part de chaque nutriment} : énergie des glucides $\times 100 /$ énergie totale, et de même pour lipides et protéines (rappel : 1 g de glucides = 4 kcal ; 1 g de protéines = 4 kcal ; 1 g de lipides = 9 kcal) ;
\item \textbf{Comparer aux valeurs recommandées} : 50--55 \% de glucides, 30--35 \% de lipides, 12--15 \% de protéines, et vérifier la présence des sept familles d'aliments.
\end{enumerate}
\end{methode}

\courssub{Le petit-déjeuner et l'hydratation : deux piliers souvent négligés}

\textbf{Le petit-déjeuner.} Après 10 à 12 heures de jeûne nocturne, les réserves d'énergie du cerveau sont basses. Un élève qui saute le petit-déjeuner — ou qui le réduit à des beignets gras — arrive en classe avec une hypoglycémie relative : somnolence, difficultés de concentration, maux de tête en fin de matinée. Le petit-déjeuner devrait fournir environ \textbf{25 \% de l'énergie de la journée} et combiner un féculent (pain, bouillie de maïs), une source de protéines (œuf, lait, arachide) et un fruit.

\textbf{L'hydratation.} Le corps d'un adolescent est constitué d'environ 60 \% d'eau. Sous le climat camerounais, les pertes par la transpiration sont importantes : il faut boire \textbf{1,5 à 2 litres d'eau par jour}, davantage en cas de sport ou de forte chaleur. L'eau est la seule boisson indispensable : les boissons sucrées n'hydratent pas mieux et apportent un excès de sucre.

\begin{popfigure}
\begin{center}
\begin{tikzpicture}[scale=1]
% Pyramide : 4 couches
\fill[popblue!60] (-4,0) -- (4,0) -- (2.8,1.5) -- (-2.8,1.5) -- cycle;
\fill[popgreen!60] (-2.8,1.5) -- (2.8,1.5) -- (1.7,3) -- (-1.7,3) -- cycle;
\fill[popgold!80] (-1.7,3) -- (1.7,3) -- (0.7,4.5) -- (-0.7,4.5) -- cycle;
\fill[poporange!80] (-0.7,4.5) -- (0.7,4.5) -- (0,6) -- cycle;
% Labels intérieurs
\node[font=\small\bfseries, popdark] at (0,0.75) {Fruits et légumes + féculents (base de chaque repas)};
\node[font=\small\bfseries, popdark] at (0,2.25) {Viandes, poissons, œufs, haricots + laitages};
\node[font=\small\bfseries, popdark] at (0,3.75) {Matières grasses};
\node[font=\footnotesize\bfseries, popdark] at (0,5.15) {Sucreries};
% Flèches et annotations latérales
\draw[->, popmuted] (4.1,0.75) -- (5.1,0.75); \node[right, font=\footnotesize, align=left] at (5.1,0.75) {À volonté, à chaque repas :\\ ndolé, folong, mangue, orange,\\ maïs, manioc, igname, plantain};
\draw[->, popmuted] (1.8,2.25) -- (5.1,2.25); \node[right, font=\footnotesize, align=left] at (5.1,2.25) {Chaque jour, en quantité modérée :\\ poisson fumé, œufs, arachides,\\ lait caillé, haricots};
\draw[->, popmuted] (0.8,3.75) -- (5.1,3.75); \node[right, font=\footnotesize, align=left] at (5.1,3.75) {En petite quantité :\\ huile de palme, huile d'arachide};
\draw[->, popmuted] (0.15,5.4) -- (5.1,5.4); \node[right, font=\footnotesize, align=left] at (5.1,5.4) {Occasionnellement :\\ boissons sucrées, beignets, bonbons};
% Eau à la base
\node[font=\small\bfseries, popblue] at (0,-0.6) {+ EAU : 1,5 à 2 litres par jour, tout au long de la journée};
\end{tikzpicture}
\end{center}
\legende{Pyramide alimentaire simplifiée adaptée aux aliments camerounais : la base (fruits, légumes, féculents) doit occuper la plus grande part des repas ; le sommet (produits sucrés) doit rester exceptionnel. L'eau accompagne toute la pyramide.}
\end{popfigure}

\begin{exoresolu}[La journée alimentaire de Junior est-elle équilibrée ?]
Énoncé. Voici la journée alimentaire typique de Junior, 16 ans, plutôt sédentaire (besoin estimé : 2 400 kcal/j) : petit-déjeuner = beignets-haricots (environ 600 kcal, très riches en lipides) ; déjeuner = grande friture de pommes de terre + boisson sucrée (environ 900 kcal, dont 55 \% de lipides) ; dîner = riz + sauce grasse (environ 800 kcal) ; aucun fruit, aucun légume, aucun laitage, très peu d'eau. 1) Calculer l'apport énergétique total. 2) Identifier au moins trois déséquilibres. 3) Proposer une amélioration du petit-déjeuner.
\tcblower
Solution détaillée.
\begin{enumerate}
\item Apport total : $600 + 900 + 800 = 2\,300$ kcal. L'apport énergétique global est proche du besoin (2 400 kcal) : le problème n'est donc \emph{pas} la quantité totale.
\item Déséquilibres :
\begin{itemize}
\item \textbf{Excès de lipides} : beignets frits, friture, sauce grasse — les lipides dépassent largement 35 \% de l'énergie (probablement 45--50 \%), au détriment des glucides complexes ;
\item \textbf{Absence totale des familles protectrices et bâtisseuses} : ni fruit, ni légume, ni laitage, ni source solide de protéines — d'où des carences probables en vitamine C, calcium et fer (ce qui peut expliquer la fatigue chronique) ;
\item \textbf{Excès de sucres rapides} (boisson sucrée) et \textbf{hydratation insuffisante} ;
\item \textbf{Petit-déjeuner inadapté} : uniquement des aliments frits.
\end{itemize}
\item Petit-déjeuner amélioré (environ 600 kcal, soit 25 \% de la ration) : bouillie de maïs au lait + un œuf ou une poignée d'arachides grillées + une orange + un verre d'eau. On obtient ainsi des sucres lents, des protéines, du calcium, de la vitamine C et une bonne hydratation dès le matin.
\end{enumerate}
Conclusion : une ration peut couvrir les besoins en énergie tout en étant profondément déséquilibrée en qualité. L'équilibre alimentaire porte sur la \textbf{variété} et les \textbf{proportions}, pas sur la seule quantité.
\end{exoresolu}

\begin{aretenir}[L'essentiel de la section]
\begin{itemize}
\item L'\cle{équilibre alimentaire} = apports adaptés aux besoins, en \textbf{quantité}, en \textbf{qualité} et \textbf{répartis sur la journée}.
\item Les \textbf{sept familles d'aliments} assurent trois rôles complémentaires : énergétique, bâtisseur, protecteur. Aucune ne doit manquer.
\item L'adolescent a des besoins élevés : \textbf{2 200 à 3 000 kcal/jour}, avec des besoins spécifiques en \textbf{protéines} (croissance), \textbf{calcium} (os), \textbf{fer} (sang, surtout chez l'adolescente) et \textbf{vitamines}.
\item Répartition recommandée : \textbf{50--55 \% de glucides, 30--35 \% de lipides, 12--15 \% de protéines}.
\item Un vrai \textbf{petit-déjeuner} (environ 25 \% de la ration) et \textbf{1,5 à 2 L d'eau par jour} sont indispensables.
\end{itemize}
\end{aretenir}

\coursec{Évaluation d'une ration alimentaire}

Dans la section précédente, nous avons vu que l'adolescent a des besoins nutritionnels importants : environ \num{2400} à \num{2800}~kcal par jour selon le sexe et l'activité physique, réparties entre glucides, lipides et protéines. Mais une question pratique se pose immédiatement : comment savoir si ce que mange réellement un élève de Première au cours d'une journée couvre ces besoins, ni trop, ni trop peu ? C'est tout l'objet de cette section : apprendre à \cle{évaluer une ration alimentaire}, c'est-à-dire à calculer ce qu'apporte concrètement une journée de repas, puis à comparer ce résultat aux besoins recommandés. C'est une compétence exigible au Probatoire et, surtout, un outil de prévention très concret : avant de lutter contre la mauvaise alimentation, il faut d'abord savoir la mesurer.

\courssub{La notion de ration alimentaire}

\textbf{Situation concrète.} À la cantine du lycée de Bafoussam, deux élèves déjeunent. Le premier prend une assiette de riz avec de la sauce arachide et un morceau de poisson ; le second prend deux beignets et une boisson gazeuse. À l'œil, le second a « moins mangé ». Pourtant, à la fin de la journée, qui a réellement absorbé le plus d'énergie ? Pour répondre, il faut une notion précise : celle de ration alimentaire.

\begin{definition}[Ration alimentaire]
La \cle{ration alimentaire} est l'ensemble des aliments et des boissons consommés par un individu au cours d'une journée (24 heures), y compris les collations, les boissons sucrées et les matières grasses utilisées pour la cuisson.
\end{definition}

\begin{attention}[Une ration, ce n'est pas un repas]
Une erreur fréquente consiste à confondre ration et repas. La ration comprend \textbf{tout} ce qui est ingéré en 24~h : le petit-déjeuner, le déjeuner, le dîner, mais aussi les arachides grillées grignotées à la récréation, le jus de fruit, la bouteille de boisson gazeuse et l'huile de friture absorbée par les beignets. Oublier un seul de ces éléments fausse l'évaluation.
\end{attention}

Pour évaluer une ration, on la décompose en ses constituants : chaque aliment apporte des \cle{glucides}, des \cle{lipides} et des \cle{protéines} en quantités mesurables. Reste à savoir combien d'énergie chaque gramme de nutriment fournit.

\courssub{Les valeurs énergétiques des nutriments}

\textbf{D'où viennent les chiffres ?} L'énergie contenue dans un aliment se mesure expérimentalement à l'aide d'une \cle{bombe calorimétrique} : on brûle complètement une masse connue de nutriment dans une enceinte remplie d'oxygène, et on mesure la chaleur dégagée par l'élévation de température d'une masse d'eau connue. Ces mesures de calorimétrie, réalisées dès la fin du XIX\textsuperscript{e} siècle, ont établi des valeurs de référence que nous admettons.

\begin{propriete}[Valeurs énergétiques de référence (admises)]
L'oxydation complète des nutriments fournit, en moyenne :
\begin{itemize}
\item \num{1}~g de \cle{glucides} $\rightarrow$ \num{4}~kcal ;
\item \num{1}~g de \cle{protéines} $\rightarrow$ \num{4}~kcal ;
\item \num{1}~g de \cle{lipides} $\rightarrow$ \num{9}~kcal.
\end{itemize}
Ces valeurs sont admises ; elles proviennent des mesures de calorimétrie et sont données dans les énoncés. On retiendra que les lipides sont plus de deux fois plus énergétiques que les glucides ou les protéines, à masse égale.
\end{propriete}

\begin{aretenir}[Unité d'énergie]
L'unité utilisée en nutrition est la \cle{kilocalorie} (kcal), souvent appelée abusivement « calorie » dans le langage courant. L'unité légale est le joule : $1~\text{kcal} = \num{4,18}~\text{kJ}$. Au Probatoire, les calculs se font en kcal.
\end{aretenir}

\courssub{Lire une table de composition des aliments}

Pour connaître la teneur d'un aliment en nutriments, on utilise une \cle{table de composition des aliments}, établie par des analyses de laboratoire. Elle donne, pour 100~g d'aliment, les masses de glucides, de lipides, de protéines et la valeur énergétique correspondante.

\begin{center}
\begin{tabularx}{\linewidth}{|l|X|X|X|X|}
\hline
\rowcolor{popblueL} \textbf{Aliment (pour 100 g)} & \textbf{Glucides (g)} & \textbf{Protéines (g)} & \textbf{Lipides (g)} & \textbf{Énergie (kcal)} \\
\hline
Riz blanc cuit & 28 & 2,7 & 0,3 & 130 \\
\hline
Arachides grillées & 21 & 26 & 49 & 629 \\
\hline
Huile végétale & 0 & 0 & 100 & 900 \\
\hline
Poisson (capitaine) grillé & 0 & 22 & 5 & 133 \\
\hline
Banane plantain mûre & 32 & 1,3 & 0,4 & 137 \\
\hline
\end{tabularx}
\end{center}

\begin{exemplebox}[Lecture de la table]
Lisons la ligne « arachides grillées » : 100~g d'arachides contiennent 21~g de glucides, 26~g de protéines et 49~g de lipides, soit une valeur énergétique de 629~kcal pour 100~g. Remarquons que l'huile, constituée de lipides presque purs, est l'aliment le plus énergétique de la table : 900~kcal pour 100~g. C'est la conséquence directe de la valeur 9~kcal/g des lipides.
\end{exemplebox}

\courssub{Méthode de calcul de l'apport énergétique}

\begin{methode}[Calculer l'apport énergétique d'un aliment]
Pour calculer l'énergie $E$ (en kcal) fournie par une portion d'aliment :
\begin{enumerate}
\item Relever dans la table les masses de glucides $G$, de protéines $P$ et de lipides $L$ contenues dans la portion (attention : si la portion n'est pas de 100~g, faire une règle de trois).
\item Appliquer la formule :
\[ E = 4 \times (G + P) + 9 \times L \quad \text{(masses en grammes)} \]
\item Pour une ration complète, additionner les apports de tous les aliments et boissons de la journée.
\end{enumerate}
\end{methode}

\begin{exemplebox}[Exemple chiffré : 100 g d'arachides grillées]
Appliquons la formule aux données de la table :
\[ E = 4 \times (21 + 26) + 9 \times 49 = 4 \times 47 + 441 = 188 + 441 = 629~\text{kcal} \]
100~g d'arachides grillées fournissent donc \textbf{629~kcal}, soit environ un quart des besoins quotidiens d'un adolescent. C'est un aliment très nourrissant, mais dont l'excès fait rapidement basculer la ration vers le surplus énergétique.
\end{exemplebox}

\begin{attention}[Les calories « cachées »]
Deux pièges classiques faussent l'évaluation d'une ration :
\begin{itemize}
\item les \textbf{boissons sucrées} (jus industriels, boissons gazeuses), qui n'ont pas de valeur nutritive réelle mais apportent beaucoup de glucides simples ;
\item l'\textbf{huile de friture} : les beignets, les beignets-haricots (accra) et les pommes frites absorbent une grande quantité d'huile, qui peut doubler leur apport énergétique.
\end{itemize}
Ne jamais les oublier dans le calcul : c'est souvent là que se cache l'excès.
\end{attention}

\begin{savaistu}[Le poids énergétique d'une boisson gazeuse]
Une bouteille de 50~cl de boisson gazeuse sucrée contient environ 53~g de sucre, soit environ \textbf{220~kcal} : près de 10~\% des besoins quotidiens d'un adolescent, absorbés en quelques minutes, sans aucune sensation de satiété durable. Boire une telle bouteille chaque jour sans activité physique supplémentaire correspond théoriquement à une prise de masse grasse de plusieurs kilogrammes par an.
\end{savaistu}

\courssub{Étude de cas guidée : la ration d'un élève de Première}

\textbf{Problème posé.} Junior, élève de Première C à Yaoundé, mène une vie très sédentaire. Voici sa journée alimentaire typique. Évaluons sa ration.

\begin{center}
\begin{tabularx}{\linewidth}{|l|X|X|X|X|}
\hline
\rowcolor{popgreenL} \textbf{Repas} & \textbf{Glucides (g)} & \textbf{Protéines (g)} & \textbf{Lipides (g)} & \textbf{Énergie (kcal)} \\
\hline
Matin : 2 beignets + bouillie sucrée & 75 & 8 & 22 & 530 \\
\hline
Midi : riz + sauce arachide & 85 & 25 & 35 & 755 \\
\hline
Soir : fufu + ndolé & 90 & 30 & 25 & 705 \\
\hline
Boisson gazeuse (50 cl) & 53 & 0 & 0 & 212 \\
\hline
\textbf{Total} & \textbf{303} & \textbf{63} & \textbf{82} & \textbf{\num{2202}} \\
\hline
\end{tabularx}
\end{center}

\begin{exoresolu}[Calcul de la ration de Junior]
\textbf{Énoncé.} À partir du tableau ci-dessus : 1) calculer l'apport énergétique de chaque repas ; 2) calculer l'apport total de la journée ; 3) comparer cet apport aux besoins recommandés pour un adolescent sédentaire (\num{2400}~kcal/jour) et conclure.
\tcblower
\textbf{Solution détaillée.}

\textbf{1) Apport de chaque repas.} On applique $E = 4 \times (G + P) + 9 \times L$.
\begin{align*}
E_{\text{matin}} &= 4 \times (75 + 8) + 9 \times 22 = 4 \times 83 + 198 = 332 + 198 = 530~\text{kcal} \\
E_{\text{midi}} &= 4 \times (85 + 25) + 9 \times 35 = 4 \times 110 + 315 = 440 + 315 = 755~\text{kcal} \\
E_{\text{soir}} &= 4 \times (90 + 30) + 9 \times 25 = 4 \times 120 + 225 = 480 + 225 = 705~\text{kcal} \\
E_{\text{boisson}} &= 4 \times (53 + 0) + 9 \times 0 = 212~\text{kcal}
\end{align*}

\textbf{2) Apport total.}
\[ E_{\text{total}} = 530 + 755 + 705 + 212 = \num{2202}~\text{kcal} \]
Vérification globale à partir des totaux : $E = 4 \times (303 + 63) + 9 \times 82 = 4 \times 366 + 738 = \num{1464} + 738 = \num{2202}$~kcal. Les deux méthodes concordent.

\textbf{3) Comparaison aux besoins.} Les besoins de Junior (adolescent sédentaire) sont d'environ \num{2400}~kcal/jour. Son apport est de \num{2202}~kcal : il est donc légèrement inférieur à ses besoins en énergie totale. \textbf{Mais} l'analyse ne s'arrête pas là : l'énergie apportée par les lipides est de $9 \times 82 = 738$~kcal, soit $\dfrac{738}{2202} \approx 33~\%$ de l'énergie totale (limite haute recommandée) ; les glucides sont presque exclusivement des glucides rapides ou raffinés (beignets, boisson sucrée, fufu), et la boisson gazeuse seule apporte 212~kcal « vides ». De plus, Junior est sédentaire : sans activité physique, même une ration proche de l'équilibre énergétique favorise le stockage des graisses. La ration de Junior est donc \textbf{quantitativement proche de l'équilibre mais qualitativement déséquilibrée}.
\end{exoresolu}

\begin{popfigure}
\begin{center}
\begin{tikzpicture}
\begin{axis}[
    ybar,
    width=0.85\linewidth,
    height=7cm,
    bar width=1.6cm,
    ymin=0, ymax=3200,
    ylabel={Énergie (kcal/jour)},
    symbolic x coords={Matin, Midi, Soir, Boisson, Total ration, Besoins},
    xtick=data,
    x tick label style={font=\small, rotate=15, anchor=east},
    nodes near coords,
    nodes near coords style={font=\small},
    every node near coord/.append style={color=popdark},
    axis lines=left,
    enlarge x limits=0.12,
]
\addplot[fill=popblue, draw=popdark] coordinates {(Matin,530) (Midi,755) (Soir,705) (Boisson,212) (Total ration,2202)};
\addplot[fill=poporange, draw=popdark] coordinates {(Besoins,2400)};
\end{axis}
\end{tikzpicture}
\end{center}
\legende{Histogramme comparant l'apport énergétique de chaque prise alimentaire de la journée de Junior (en bleu) à ses besoins quotidiens recommandés (en orange, \num{2400}~kcal). Le total de la ration (\num{2202}~kcal) est légèrement inférieur aux besoins, mais la boisson sucrée apporte 212~kcal sans valeur nutritive.}
\end{popfigure}

\courssub{Comparer la ration aux besoins recommandés : déficit, équilibre ou excès}

Une fois l'apport total calculé, la dernière étape de l'évaluation consiste à le comparer aux besoins recommandés, vus dans la section précédente (environ \num{2400} à \num{2800}~kcal/jour pour un adolescent selon le sexe et l'activité physique).

\begin{definition}[Les trois situations d'une ration]
\begin{itemize}
\item \cle{Déficit} : l'apport est inférieur aux besoins. À long terme : amaigrissement, fatigue, baisse des performances scolaires et sportives, retard de croissance chez l'adolescent.
\item \cle{Équilibre} : l'apport couvre les besoins, avec une répartition correcte des nutriments (environ 50--55~\% de glucides, 30--35~\% de lipides, 12--15~\% de protéines).
\item \cle{Excès} : l'apport dépasse durablement les besoins. L'énergie non utilisée est stockée sous forme de graisse dans le tissu adipeux : c'est la voie vers le surpoids puis l'obésité.
\end{itemize}
\end{definition}

\begin{attention}[Ne pas se limiter au total énergétique]
Comme le montre le cas de Junior, une ration peut être proche de l'équilibre en calories tout en étant \textbf{qualitativement mauvaise} : excès de sucres rapides, excès de graisses de friture, carence en fruits, légumes et protéines de qualité. Au Probatoire, une évaluation complète commente toujours deux aspects : la \textbf{quantité} (comparaison aux besoins en kcal) et la \textbf{qualité} (répartition et nature des nutriments).
\end{attention}

\begin{exercice}[À toi de jouer]
Une élève consomme au déjeuner 150~g de riz cuit, 80~g de poisson grillé et 10~g d'huile. En utilisant la table de composition de cette leçon : 1) calcule l'apport énergétique de ce repas ; 2) ce repas peut-il convenir pour un déjeuner représentant environ 35~\% des besoins quotidiens d'une adolescente (\num{2400}~kcal/jour) ? Justifie.
\end{exercice}

\begin{corrige}[Exercice « À toi de jouer »]
\textbf{1)} D'après la table (valeurs pour 100~g) :
\begin{itemize}
\item Riz : $150~\text{g} \rightarrow 1,5 \times 130 = 195$~kcal ;
\item Poisson : $80~\text{g} \rightarrow 0,8 \times 133 = 106,4$~kcal ;
\item Huile : $10~\text{g} \rightarrow 0,1 \times 900 = 90$~kcal.
\end{itemize}
Total : $E = 195 + 106,4 + 90 = 391,4 \approx 391$~kcal.

\textbf{2)} 35~\% de \num{2400}~kcal représentent $0,35 \times 2400 = 840$~kcal. Le repas n'apporte que 391~kcal : il est nettement insuffisant pour un déjeuner (déficit d'environ 449~kcal). Il faudrait augmenter la portion de riz, ajouter une sauce (arachide, légumes) et un fruit pour atteindre un apport adapté.
\end{corrige}

\begin{aretenir}[L'essentiel de la section]
\begin{itemize}
\item La \cle{ration alimentaire} est l'ensemble des aliments et boissons consommés en 24~h.
\item Valeurs énergétiques admises : 4~kcal/g pour les glucides et les protéines, 9~kcal/g pour les lipides.
\item Formule de calcul : $E = 4 \times (G + P) + 9 \times L$ (masses en grammes).
\item L'évaluation se conclut par la comparaison aux besoins : déficit, équilibre ou excès, sans oublier l'analyse qualitative de la ration.
\end{itemize}
\end{aretenir}

Savoir évaluer une ration est une première étape. Mais comment juger rapidement, à l'échelle d'une population ou d'un individu, si le poids est adapté à la taille ? C'est l'objet de la section suivante : l'indice de masse corporelle (IMC), son calcul et son interprétation.

\coursec{L'IMC : calcul et interprétation}

Dans la section précédente, nous avons appris à évaluer une ration alimentaire : nous savons désormais comparer les apports énergétiques et nutritionnels d'un repas aux besoins recommandés pour un adolescent. Mais une question demeure : comment savoir si, globalement, l'alimentation et le mode de vie d'une personne ont conduit à une corpulence compatible avec une bonne santé ? Comment un médecin, un infirmier scolaire ou un agent de santé communautaire peut-il repérer rapidement qu'un élève est en surpoids ou, au contraire, en dénutrition ? C'est précisément le rôle d'un indicateur simple, calculable avec seulement deux mesures : la masse et la taille. Cet indicateur s'appelle l'\cle{Indice de Masse Corporelle}, ou \cle{IMC}.

\courssub{Qu'est-ce que l'IMC ? Définition et intuition}

\textbf{Une situation concrète pour commencer.}\par

Lors de la visite médicale scolaire au lycée de Nkongsamba, l'infirmier pèse et mesure chaque élève, puis note un nombre dans le dossier : 19,8 pour Amina, 27,3 pour Rodrigue, 16,9 pour Sandrine. Pourtant, Rodrigue n'est pas le plus lourd des trois en valeur absolue : c'est simplement que, rapportée à sa taille, sa masse est élevée. L'infirmier ne juge donc pas la masse seule : il la \emph{rapporte à la taille}. C'est toute l'idée de l'IMC.

\textbf{Pourquoi ne pas utiliser la masse seule ?}\par

Une masse de 80 kg n'a pas la même signification selon qu'elle est portée par une personne mesurant 1,60 m ou 1,90 m. Pour comparer les corpulences entre individus de tailles différentes, il faut un indice qui « neutralise » l'effet de la taille. Les études épidémiologiques ont montré que le rapport de la masse au \emph{carré} de la taille est celui qui varie le moins avec la taille elle-même : il permet donc des comparaisons fiables entre personnes.

\begin{definition}[Indice de Masse Corporelle (IMC)]
L'\cle{Indice de Masse Corporelle} (IMC), appelé aussi indice de Quételet, est le rapport de la masse d'une personne au carré de sa taille :
\[ \mathrm{IMC} = \dfrac{m}{t^{2}} \]
où $m$ est la masse exprimée en \cle{kilogrammes} (kg) et $t$ la taille exprimée en \cle{mètres} (m). L'IMC s'exprime en \cle{kg/m$^{2}$} (souvent écrit kg.m$^{-2}$).
\end{definition}

\begin{attention}[Les deux erreurs classiques de calcul]
Deux erreurs reviennent très souvent dans les copies :
\begin{itemize}
\item \textbf{Oublier d'élever la taille au carré} : calculer $m/t$ au lieu de $m/t^{2}$. Pour 70 kg et 1,70 m, cela donnerait 41,2 au lieu de 24,2 : une conclusion totalement fausse !
\item \textbf{Oublier de convertir la taille en mètres} : si la taille est donnée en centimètres (170 cm), il faut d'abord convertir : $t = 1,70$ m. Calculer $70/170^{2}$ donnerait un résultat absurde ($\approx 0,0024$).
\end{itemize}
Réflexe à acquérir : \emph{taille en mètres, taille au carré}.
\end{attention}

\courssub{Calcul guidé de l'IMC : exemples concrets}

\begin{methode}[Calculer et interpréter un IMC en trois étapes]
\begin{enumerate}
\item \textbf{Calculer} : vérifier que la masse est en kg et la taille en m (convertir si nécessaire), élever la taille au carré, puis effectuer la division $m/t^{2}$.
\item \textbf{Situer} : placer la valeur obtenue dans les intervalles de référence de l'OMS (voir tableau ci-dessous) pour déterminer la catégorie de corpulence.
\item \textbf{Conclure} : relier la catégorie au risque sanitaire éventuel, en gardant à l'esprit les limites de l'indice (adolescent, sportif musclé, femme enceinte...).
\end{enumerate}
\end{methode}

\begin{exemplebox}[L'élève de 70 kg pour 1,70 m]
Un élève de Première C pèse $m = 70$ kg et mesure $t = 1,70$ m. Calculons son IMC pas à pas :
\begin{align*}
t^{2} &= 1,70 \times 1,70 = 2,89 \ \mathrm{m^{2}}\\
\mathrm{IMC} &= \dfrac{70}{2,89} \approx 24,2 \ \mathrm{kg/m^{2}}
\end{align*}
La valeur 24,2 se situe dans l'intervalle $[18,5 \ ;\ 24,9]$ : cet élève a une \textbf{corpulence normale}. On remarque cependant qu'il est proche de la limite supérieure (24,9) : un conseil de prévention (activité physique régulière, limitation des boissons sucrées) est pertinent pour éviter de basculer dans le surpoids.
\end{exemplebox}

\begin{exoresolu}[Deux camarades, deux situations différentes]
\textbf{Énoncé.} Lors de la même visite médicale, on relève : Béatrice, $m = 82$ kg, $t = 1,60$ m ; Charles, $m = 48$ kg, $t = 1,75$ m. (a) Calculer l'IMC de chacun. (b) Indiquer la catégorie de corpulence selon l'OMS. (c) Formuler un conseil de santé adapté à chacun.
\tcblower
\textbf{Solution détaillée.}

(a) \textbf{Calculs.}
\begin{align*}
\text{Béatrice : } \mathrm{IMC} &= \dfrac{82}{1,60^{2}} = \dfrac{82}{2,56} \approx 32,0 \ \mathrm{kg/m^{2}}\\
\text{Charles : } \mathrm{IMC} &= \dfrac{48}{1,75^{2}} = \dfrac{48}{3,0625} \approx 15,7 \ \mathrm{kg/m^{2}}
\end{align*}

(b) \textbf{Catégories OMS.} Béatrice : $\mathrm{IMC} \approx 32,0 \geq 30$ : \textbf{obésité} (classe I). Charles : $\mathrm{IMC} \approx 15,7 < 18,5$ : \textbf{insuffisance pondérale} (dénutrition modérée à marquée).

(c) \textbf{Conseils.} Pour Béatrice : consultation médicale pour rechercher une complication (glycémie à jeun pour dépister un diabète de type 2, mesure de la tension artérielle), réduction progressive des apports énergétiques (fritures, boissons sucrées) et mise en place d'une activité physique régulière (au moins 30 à 60 minutes de marche rapide par jour). Pour Charles : enquête alimentaire (ration insuffisante ? maladie chronique ? parasitose intestinale ?), enrichissement de la ration en énergie et en protéines (haricots, arachides, poisson, œufs), et suivi médical pour dépister une anémie ou une infection. Dans les deux cas, l'IMC n'est qu'un \emph{signal d'alerte} : il doit être confirmé par un examen clinique.
\end{exoresolu}

\courssub{Les seuils d'interprétation de l'OMS}

L'Organisation mondiale de la santé (OMS) a défini, à partir de vastes études épidémiologiques reliant l'IMC à la mortalité et aux maladies, des intervalles de référence valables pour l'\textbf{adulte} (homme ou femme, à partir d'environ 18-19 ans).

\begin{center}
\begin{tabularx}{\linewidth}{|l|c|X|}
\hline
\rowcolor{popblueL} \textbf{Catégorie} & \textbf{IMC (kg/m$^{2}$)} & \textbf{Risque sanitaire associé} \\
\hline
Insuffisance pondérale & $< 18,5$ & Dénutrition, anémie, baisse des défenses immunitaires \\
\hline
Corpulence normale & $18,5$ à $24,9$ & Risque minimal (zone de référence) \\
\hline
Surpoids & $25$ à $29,9$ & Risque accru de diabète de type 2, hypertension \\
\hline
Obésité (classe I) & $30$ à $34,9$ & Risque élevé de maladies cardiovasculaires \\
\hline
Obésité (classe II) & $35$ à $39,9$ & Risque très élevé \\
\hline
Obésité massive (classe III) & $\geq 40$ & Risque extrêmement élevé \\
\hline
\end{tabularx}
\end{center}

\begin{aretenir}[Les quatre seuils à connaître]
Pour le Probatoire, retiens les quatre grandes catégories : $\mathrm{IMC} < 18,5$ : \cle{insuffisance pondérale} ; $18,5 \leq \mathrm{IMC} \leq 24,9$ : \cle{corpulence normale} ; $25 \leq \mathrm{IMC} \leq 29,9$ : \cle{surpoids} ; $\mathrm{IMC} \geq 30$ : \cle{obésité}.
\end{aretenir}

La figure ci-dessous visualise ces zones : la courbe représente schématiquement le niveau de risque sanitaire en fonction de l'IMC. On observe que le risque est minimal dans la zone verte (corpulence normale) et augmente de part et d'autre : c'est une relation « en cuvette », car la dénutrition est aussi dangereuse que l'obésité.

\begin{popfigure}
\begin{center}
\begin{tikzpicture}
\begin{axis}[
width=\linewidth, height=7cm,
xlabel={IMC (kg/m$^{2}$)}, ylabel={Risque sanitaire relatif},
xmin=12, xmax=42, ymin=0, ymax=5.2,
xtick={15,18.5,25,30,35,40},
ytick=\empty,
axis lines=left, clip=false,
every axis x label/.style={at={(ticklabel* cs:1.02)}, anchor=west},
every axis y label/.style={at={(ticklabel* cs:1.02)}, anchor=south},
]
% Zones colorées
\fill[poporangeL] (axis cs:12,0) rectangle (axis cs:18.5,5);
\fill[popgreenL] (axis cs:18.5,0) rectangle (axis cs:25,5);
\fill[popgoldL] (axis cs:25,0) rectangle (axis cs:30,5);
\fill[poppinkL] (axis cs:30,0) rectangle (axis cs:42,5);
% Courbe de risque schématique (en cuvette)
\addplot[popcurve,domain=14:41,samples=200,very thick] {0.035*(x-21.5)^2 + 0.6};
% Lignes de seuil
\draw[dashed,popdark] (axis cs:18.5,0) -- (axis cs:18.5,5);
\draw[dashed,popdark] (axis cs:25,0) -- (axis cs:25,5);
\draw[dashed,popdark] (axis cs:30,0) -- (axis cs:30,5);
% Étiquettes des zones
\node[popdark,font=\small] at (axis cs:15.4,4.5) {Insuffisance};
\node[popdark,font=\small] at (axis cs:15.4,4.1) {pondérale};
\node[popdark,font=\small] at (axis cs:21.7,4.5) {Corpulence};
\node[popdark,font=\small] at (axis cs:21.7,4.1) {normale};
\node[popdark,font=\small] at (axis cs:27.5,4.5) {Surpoids};
\node[popdark,font=\small] at (axis cs:35.8,4.5) {Obésité};
% Seuils annotés
\node[popdark,font=\footnotesize,below] at (axis cs:18.5,0) {18,5};
\node[popdark,font=\footnotesize,below] at (axis cs:25,0) {25};
\node[popdark,font=\footnotesize,below] at (axis cs:30,0) {30};
\end{axis}
\end{tikzpicture}
\end{center}
\legende{Zones d'interprétation de l'IMC selon l'OMS (adulte). La courbe schématise l'évolution du risque sanitaire : il est minimal dans la zone de corpulence normale (18,5 à 24,9 kg/m$^{2}$) et augmente aussi bien vers les IMC faibles (dénutrition) que vers les IMC élevés (surpoids, obésité).}
\end{popfigure}

\courssub{Les limites de l'IMC : un indicateur à manier avec discernement}

L'IMC est un outil de \emph{dépistage} rapide et peu coûteux, mais il n'est pas un diagnostic. Il comporte des limites qu'un bon élève de Première doit connaître.

\textbf{1. L'IMC ne distingue pas la masse musculaire de la masse grasse.}\par

Le muscle est plus dense que la graisse : à volume égal, il est plus lourd. Un joueur de rugby très musclé peut afficher un IMC de 28 (classé « surpoids ») alors que sa masse grasse est faible et sa santé excellente. Inversement, une personne âgée ou sédentaire peut avoir un IMC « normal » mais une proportion de graisse élevée et des muscles fondus : on parle d'\emph{obésité sarcopénique}. Pour affiner l'évaluation, on mesure le \cle{tour de taille} (la graisse abdominale est la plus dangereuse : risque élevé au-delà de 94 cm chez l'homme et 80 cm chez la femme) ou le pourcentage de masse grasse (balance impédancemètre, pli cutané).

\textbf{2. L'IMC ne renseigne pas sur la répartition de la graisse.}\par

La graisse localisée autour de l'abdomen (obésité dite « androïde » ou « en pomme ») est fortement associée au diabète de type 2 et aux maladies cardiovasculaires, alors que la graisse des hanches et des cuisses (obésité « gynoïde » ou « en poire ») l'est beaucoup moins. Deux personnes de même IMC n'ont donc pas forcément le même risque.

\textbf{3. Chez l'adolescent, les seuils de l'adulte ne s'appliquent pas directement.}\par

C'est un point essentiel pour vous, élèves de Première. Pendant la croissance, la masse et la taille évoluent rapidement, et à des rythmes différents selon l'âge et le sexe (la puberté modifie la composition corporelle : les filles stockent physiologiquement plus de graisse, les garçons développent plus de muscle). Un IMC de 24 peut être parfaitement normal à 18 ans mais signifier un surpoids à 13 ans. C'est pourquoi, chez l'enfant et l'adolescent, on utilise les \cle{courbes de croissance} (courbes d'IMC-pour-âge de l'OMS ou du carnet de santé) : on reporte l'IMC de l'adolescent sur la courbe correspondant à son âge et à son sexe, et l'on regarde dans quel \emph{percentile} il se situe. Au-delà du 97\textsuperscript{e} percentile, on parle d'obésité ; entre le 85\textsuperscript{e} et le 97\textsuperscript{e}, de surpoids ; en dessous du 3\textsuperscript{e} percentile, d'insuffisance pondérale.

\begin{attention}[Ne pas appliquer mécaniquement les seuils adultes à un adolescent]
Dans un exercice portant sur un élève de 16-17 ans, les seuils OMS adultes donnent une \emph{première indication} (c'est ce qu'on te demandera généralement), mais la conclusion rigoureuse mentionne que l'interprétation exacte chez l'adolescent se fait sur les courbes d'IMC-pour-âge. Cette nuance est appréciée dans les copies.
\end{attention}

\courssub{IMC et risques sanitaires : ce que disent les données}

\textbf{Démarche d'investigation : d'où viennent ces seuils ?}\par

\begin{itemize}
\item \textbf{Observation} : les assureurs et les médecins constatent depuis longtemps que les personnes très corpulentes meurent plus souvent d'infarctus ou de diabète, et que les personnes très maigres meurent plus souvent d'infections.
\item \textbf{Hypothèse} : il existerait une corpulence « optimale » associée à la mortalité la plus faible.
\item \textbf{Données} : des études de cohorte portant sur des centaines de milliers de personnes suivies pendant des décennies (études américaines, européennes, puis mondiales coordonnées par l'OMS) ont mesuré l'IMC des participants puis enregistré les maladies et les décès.
\item \textbf{Interprétation} : la mortalité et la fréquence du diabète de type 2, de l'hypertension et des maladies cardiovasculaires sont minimales pour un IMC compris entre 18,5 et 25, puis augmentent progressivement au-delà, et fortement au-delà de 30. Elles augmentent aussi en dessous de 18,5.
\item \textbf{Conclusion} : les seuils 18,5 / 25 / 30 ne sont pas arbitraires : ils correspondent aux valeurs où le risque sanitaire commence à s'élever significativement. Attention cependant : la \emph{corrélation} statistique ne prouve pas à elle seule la \emph{causalité} ; c'est la convergence avec les mécanismes physiologiques (voir ci-dessous) qui rend le lien solide.
\end{itemize}

\textbf{Les mécanismes qui relient un IMC élevé aux maladies.}\par

Un excès de masse grasse, surtout abdominale, n'est pas un simple stock inerte : le tissu adipeux sécrète des substances (adipokines, médiateurs de l'inflammation) qui perturbent le métabolisme. Trois conséquences majeures :
\begin{itemize}
\item \textbf{Diabète de type 2} : l'excès de graisse rend les cellules (muscles, foie) moins sensibles à l'\cle{insuline} — c'est l'\emph{insulinorésistance}. Le pancréas doit produire de plus en plus d'insuline pour maintenir une glycémie normale, puis s'épuise : la glycémie s'élève durablement (hyperglycémie chronique), définissant le diabète de type 2, qui endommage à terme les vaisseaux, les reins, les yeux et les nerfs.
\item \textbf{Hypertension artérielle} : la masse grasse augmente le volume sanguin nécessaire et perturbe le fonctionnement des vaisseaux ; le cœur doit pomper contre une pression plus élevée, ce qui use prématurément les artères.
\item \textbf{Maladies cardiovasculaires} : l'excès de graisse s'accompagne souvent d'une élévation des lipides sanguins (LDL-cholestérol, triglycérides) qui favorise le dépôt de plaques d'\cle{athérome} sur la paroi des artères (athérosclérose). Une plaque peut obstruer une artère coronaire (infarctus du myocarde) ou cérébrale (accident vasculaire cérébral).
\end{itemize}

\textbf{Les risques d'un IMC trop faible.}\par

L'insuffisance pondérale ($\mathrm{IMC} < 18,5$) traduit le plus souvent une ration alimentaire insuffisante en quantité ou en qualité, ou une maladie chronique (tuberculose, parasitoses intestinales, VIH/Sida non traité). Ses conséquences :
\begin{itemize}
\item la \textbf{dénutrition} : fonte musculaire, fatigue, retard de croissance et troubles de la concentration chez l'adolescent (conséquences directes sur les résultats scolaires !) ;
\item l'\textbf{anémie} : le déficit en fer et en protéines diminue la production d'hémoglobine, réduisant le transport du dioxygène vers les organes ;
\item la \textbf{baisse des défenses immunitaires} : la synthèse des anticorps et des cellules immunitaires exige des protéines et des micronutriments ; l'organisme dénutri devient plus vulnérable aux infections (paludisme plus sévère, infections respiratoires et digestives plus fréquentes).
\end{itemize}

\courssub{Comprendre la variation de l'IMC : la balance énergétique}

Pourquoi l'IMC augmente-t-il ou diminue-t-il ? La réponse tient en un principe simple : la \cle{balance énergétique}, c'est-à-dire la comparaison entre l'énergie apportée par les aliments et l'énergie dépensée par l'organisme.

\begin{propriete}[La balance énergétique (connaissance, pas de démonstration exigible au Probatoire)]
L'énergie dépensée par l'organisme comprend le \cle{métabolisme de base} (énergie des fonctions vitales au repos : cœur, respiration, température corporelle — environ 60 à 70 \% des dépenses), la thermogenèse alimentaire (digestion) et l'\cle{activité physique}. Trois cas :
\begin{itemize}
\item Apports $=$ Dépenses : masse stable, IMC constant (équilibre).
\item Apports $>$ Dépenses : l'excès est stocké, principalement sous forme de triglycérides dans le tissu adipeux : la masse augmente, l'IMC monte.
\item Apports $<$ Dépenses : l'organisme puise dans ses réserves (glycogène puis graisses) : la masse diminue, l'IMC baisse.
\end{itemize}
\end{propriete}

\begin{popfigure}
\begin{center}
\begin{tikzpicture}[font=\small, >=Stealth, node distance=0.4cm]
% Plateau central
\node[draw=popdark, fill=popcream, rounded corners, minimum width=3.4cm, minimum height=1.6cm, align=center] (bilan) at (0,0) {\textbf{Bilan énergétique}\\ de l'organisme};
% Apports
\node[draw=popgreen, fill=popgreenL, rounded corners, minimum width=3.2cm, minimum height=1.3cm, align=center, left=1.6cm of bilan] (apports) {\textbf{APPORTS}\\ Aliments et boissons\\ (énergie en kcal)};
% Dépenses
\node[draw=poporange, fill=poporangeL, rounded corners, minimum width=3.2cm, minimum height=1.3cm, align=center, right=1.6cm of bilan] (depenses) {\textbf{DÉPENSES}\\ Métabolisme de base\\ + activité physique};
\draw[->, very thick, popgreen] (apports) -- (bilan);
\draw[->, very thick, poporange] (bilan) -- (depenses);
% Conséquences
\node[draw=poppink, fill=poppinkL, rounded corners, align=center, below left=1.1cm and -1.2cm of bilan] (stock) {Apports $>$ Dépenses\\ $\Rightarrow$ stockage de graisses\\ $\Rightarrow$ IMC $\nearrow$};
\node[draw=popblue, fill=popblueL, rounded corners, align=center, below=1.1cm of bilan] (equilibre) {Apports $=$ Dépenses\\ $\Rightarrow$ masse stable\\ $\Rightarrow$ IMC constant};
\node[draw=popgold, fill=popgoldL, rounded corners, align=center, below right=1.1cm and -1.2cm of bilan] (destock) {Apports $<$ Dépenses\\ $\Rightarrow$ utilisation des réserves\\ $\Rightarrow$ IMC $\searrow$};
\draw[->, thick, popdark] (bilan.south) -- (equilibre.north);
\draw[->, thick, popdark] (bilan.south west) -- (stock.north);
\draw[->, thick, popdark] (bilan.south east) -- (destock.north);
\end{tikzpicture}
\end{center}
\legende{Schéma de la balance énergétique. La masse corporelle (et donc l'IMC) dépend de la comparaison entre l'énergie apportée par l'alimentation et l'énergie dépensée par le métabolisme de base et l'activité physique. Un déséquilibre prolongé dans un sens ou dans l'autre fait sortir l'IMC de la zone de corpulence normale.}
\end{popfigure}

Ce schéma explique pourquoi la lutte contre la mauvaise alimentation \emph{et} l'inactivité physique sont indissociables : agir sur les apports sans agir sur les dépenses (ou l'inverse) est deux fois moins efficace. C'est le fondement des recommandations que nous étudierons dans les sections suivantes.

\begin{savaistu}[Une épidémie mondiale qui n'épargne pas l'Afrique]
Selon l'OMS, la prévalence mondiale de l'obésité a presque \textbf{triplé depuis 1975}. En 2022, plus d'un milliard de personnes dans le monde vivaient avec l'obésité, dont environ 160 millions d'enfants et d'adolescents. Ce phénomène, longtemps considéré comme réservé aux pays riches, touche désormais les villes africaines : à Yaoundé, Douala ou Bamenda, la sédentarité (transport motorisé, écrans) combinée à la consommation croissante d'aliments ultratransformés, de fritures et de boissons sucrées fait progresser le surpoids chez les adolescents, alors même que la dénutrition persiste dans certaines zones rurales. Les pays comme le Cameroun font ainsi face à un « double fardeau nutritionnel » : dénutrition et obésité coexistent, parfois dans la même ville, parfois dans la même famille. Savoir calculer et interpréter un IMC est donc une compétence de santé publique directement utile dans ton environnement.
\end{savaistu}

\begin{exercice}[À toi de jouer]
Trois personnes se présentent au centre de santé : Mme Ngo Bell, 65 kg pour 1,55 m ; M. Etoundi, 95 kg pour 1,80 m ; la petite Larissa, 17 ans, 44 kg pour 1,62 m. (a) Calcule l'IMC de chacun (arrondi au dixième). (b) Classe-les selon les catégories de l'OMS. (c) Pour Larissa, quelle précaution d'interprétation faut-il prendre ? (d) Propose pour chacun un premier conseil de santé.
\end{exercice}

\begin{corrige}[Exercice : À toi de jouer]
(a) Mme Ngo Bell : $\mathrm{IMC} = 65/1,55^{2} = 65/2,4025 \approx 27,1$ kg/m$^{2}$. M. Etoundi : $\mathrm{IMC} = 95/1,80^{2} = 95/3,24 \approx 29,3$ kg/m$^{2}$. Larissa : $\mathrm{IMC} = 44/1,62^{2} = 44/2,6244 \approx 16,8$ kg/m$^{2}$.
(b) Mme Ngo Bell : surpoids ($25 \leq 27,1 \leq 29,9$). M. Etoundi : surpoids ($29,3$), mais très proche du seuil d'obésité (30) : vigilance renforcée. Larissa : insuffisance pondérale ($16,8 < 18,5$).
(c) Larissa a 17 ans : c'est encore une adolescente. Les seuils adultes donnent une indication, mais l'interprétation rigoureuse exige de reporter son IMC sur les courbes d'IMC-pour-âge (courbes de croissance OMS) pour son sexe et son âge.
(d) Mme Ngo Bell : activité physique régulière (marche rapide 30 min/jour), réduction des fritures et boissons sucrées, contrôle annuel de la glycémie et de la tension. M. Etoundi : mêmes conseils, plus une consultation médicale rapprochée (risque proche de l'obésité : dépistage du diabète de type 2 et de l'hypertension). Larissa : enquête alimentaire et recherche d'une cause (ration insuffisante, parasitose, infection chronique), enrichissement de la ration en énergie, protéines et fer, suivi médical.
\end{corrige}

\begin{aretenir}[L'essentiel de la section]
\begin{itemize}
\item $\mathrm{IMC} = m/t^{2}$, avec $m$ en kg et $t$ en m ; l'IMC s'exprime en kg/m$^{2}$.
\item Seuils OMS (adulte) : $< 18,5$ insuffisance pondérale ; 18,5-24,9 corpulence normale ; 25-29,9 surpoids ; $\geq 30$ obésité.
\item Méthode : calculer (taille en mètres, au carré !), situer dans les intervalles, conclure sur le risque.
\item Limites : l'IMC ne distingue pas muscle et graisse ; chez l'adolescent, on utilise les courbes d'IMC-pour-âge.
\item IMC élevé : risque de diabète de type 2, hypertension, maladies cardiovasculaires. IMC faible : dénutrition, anémie, baisse des défenses immunitaires.
\item La variation de l'IMC s'explique par la balance énergétique : apports alimentaires contre dépenses (métabolisme de base + activité physique).
\end{itemize}
\end{aretenir}

Maintenant que tu sais mesurer et interpréter la corpulence, il reste à comprendre en détail \emph{comment} la sédentarité et la mauvaise alimentation font basculer l'organisme vers ces maladies : c'est l'objet de la section suivante, consacrée aux conséquences de la sédentarité et de la mauvaise alimentation.

\coursec{Conséquences de la sédentarité et de la mauvaise alimentation}

Dans la section précédente, tu as appris à calculer l'IMC et à l'interpréter : un IMC supérieur à 25 signale un surpoids, supérieur à 30 une obésité. Mais pourquoi ces chiffres sont-ils si importants ? Parce que derrière l'IMC se cachent des maladies graves, souvent silencieuses pendant des années. C'est ce que nous allons découvrir maintenant : comprendre \emph{comment} la sédentarité et la mauvaise alimentation détruisent progressivement l'organisme.

\courssub{Deux comportements à risque : la sédentarité et la malnutrition par excès}

\textbf{Situation concrète.} À Yaoundé, deux frères de 16 ans ont des modes de vie très différents. L'aîné marche \SI{40}{min} par jour pour aller au lycée, joue au football le week-end et mange des repas familiaux équilibrés. Le cadet se déplace en moto-taxi, passe ses soirées devant son téléphone, grignote des beignets et boit deux boissons sucrées par jour. En trois ans, le cadet a pris \SI{15}{kg} et se fatigue au moindre effort. Comment expliquer une telle différence ?

\begin{definition}[La sédentarité]
La \cle{sédentarité} est un mode de vie caractérisé par une dépense physique quotidienne insuffisante : moins de \SI{30}{min} d'activité physique modérée (marche rapide, vélo, sport, travaux physiques) par jour. Elle se traduit par une faible \cle{dépense énergétique} quotidienne.
\end{definition}

\begin{definition}[La malnutrition par excès]
La \cle{malnutrition par excès} (ou malbouffe) est une alimentation déséquilibrée caractérisée par une surabondance de \cle{sucres rapides} (boissons sucrées, friandises) et de \cle{graisses} (fritures, huiles en excès), souvent associée à une carence en fibres, vitamines et minéraux. Contrairement à la sous-nutrition (déficit), c'est une malnutrition « par trop-plein ».
\end{definition}

\begin{attention}[Malnutrition ne veut pas dire seulement « manque de nourriture »]
Au Probatoire, beaucoup d'élèves définissent la malnutrition uniquement comme une insuffisance alimentaire. Erreur ! La malnutrition recouvre \textbf{deux} situations : la sous-nutrition (carences) \emph{et} la malnutrition par excès (surcharge en sucres et graisses). Un adolescent obèse souffre de malnutrition, même s'il mange beaucoup.
\end{attention}

\courssub{Le mécanisme de l'obésité : une question de balance énergétique}

\textbf{Observation.} Pourquoi grossit-on ? L'intuition « on grossit parce qu'on mange trop » est incomplète : un sportif de haut niveau mange énormément sans grossir. Le vrai critère est la \emph{comparaison} entre ce qui entre et ce qui sort.

\begin{definition}[La balance énergétique]
L'organisme fonctionne comme un bilan comptable : les \cle{apports énergétiques} (aliments, exprimés en kilocalories, kcal) sont comparés aux \cle{dépenses énergétiques} (métabolisme de base + activité physique + thermogenèse post-prandiale).
\begin{itemize}
\item Si apports = dépenses : la masse corporelle reste stable.
\item Si apports $>$ dépenses de façon durable : l'excédent est transformé en graisses et \cle{stocké dans le tissu adipeux} (sous la peau, autour des organes) : c'est la prise de masse grasse, qui mène au surpoids puis à l'\cle{obésité}.
\item Si apports $<$ dépenses : l'organisme puise dans ses réserves : perte de masse.
\end{itemize}
\end{definition}

\begin{propriete}[Prise de masse grasse liée à un excès énergétique chronique — admise]
Un excès chronique d'environ \SI{100}{kcal} par jour par rapport aux dépenses (soit à peine un petit beignet ou un verre de boisson sucrée en trop) entraîne une prise de masse grasse de l'ordre de \textbf{4 à 5 kg par an}. Cette propriété, admise, montre qu'un petit excès \emph{quotidien} suffit à provoquer une obésité en quelques années.
\end{propriete}

\begin{exemplebox}[Le calcul qui fait réfléchir]
Un adolescent boit chaque jour une boisson sucrée de \SI{33}{cL} (environ \SI{140}{kcal}) sans augmenter son activité physique. Excès annuel approximatif : $140 \times 365 \approx 51\,000$ kcal. Sachant qu'un kilogramme de tissu adipeux correspond à environ \SI{7\,700}{kcal}, cela représente $\dfrac{51\,000}{7\,700} \approx 6,6$ kg de graisse potentielle en un an, si rien ne compense. Une simple habitude quotidienne peut donc transformer un corpulent normal en obèse en 2 ou 3 ans.
\end{exemplebox}

\courssub{Le diabète de type 2 : quand l'insuline n'agit plus}

\textbf{Observation.} Au Centre Hospitalier d'Essos à Yaoundé, les consultations de diabétologie reçoivent désormais des patients de plus en plus jeunes, parfois des adolescents. Pourtant, le diabète de type 2 était autrefois inconnu avant 40 ans. Que s'est-il passé ?

\textbf{Hypothèse.} L'excès de masse grasse perturbe l'action de l'insuline, l'hormone qui régule le glucose sanguin.

\textbf{Rappel physiologique.} Après un repas, la glycémie (taux de glucose dans le sang) augmente. Le pancréas sécrète alors de l'\cle{insuline}, hormone qui permet aux cellules (muscles, foie, tissu adipeux) de capter le glucose : la glycémie redescend vers sa valeur normale (environ \SI{1}{g/L} à jeun).

\begin{definition}[Le diabète de type 2]
Le \cle{diabète de type 2} est une maladie chronique caractérisée par une \cle{hyperglycémie chronique} (glycémie à jeun $\geq$ \SI{1,26}{g/L}) due à une \cle{résistance à l'insuline} : les cellules répondent mal à l'insuline, le glucose s'accumule dans le sang. Cette résistance est fortement favorisée par l'excès de masse grasse, surtout abdominale.
\end{definition}

\textbf{Mécanisme en chaîne.} L'excès de graisse corporelle rend les cellules musculaires et hépatiques moins sensibles à l'insuline. Le pancréas compense en produisant davantage d'insuline (hyperinsulinémie compensatrice), puis s'épuise (insulinopénie relative). Le glucose reste alors durablement en excès dans le sang.

\textbf{Les complications de l'hyperglycémie chronique.} Le glucose en excès endommage progressivement les petits vaisseaux sanguins (microangiopathie) et les nerfs (neuropathie) :
\begin{itemize}
\item \textbf{Les yeux} : atteinte de la rétine (rétinopathie diabétique), pouvant conduire à la cécité ;
\item \textbf{Les reins} : atteinte des filtres rénaux (néphropathie), pouvant conduire à l'insuffisance rénale et à la dialyse ;
\item \textbf{Les nerfs} : neuropathie (fourmillements, pertes de sensibilité des pieds), favorisant les plaies infectées (pied diabétique), parfois jusqu'à l'amputation.
\end{itemize}

\begin{attention}[Un diabète qui rajeunit]
Le diabète de type 2, autrefois appelé « diabète de l'adulte » ou « diabète gras », apparaît désormais chez des \textbf{adolescents en surpoids}, y compris au Cameroun. Ce n'est plus une maladie réservée aux personnes âgées : c'est une maladie liée au mode de vie, qui se prépare dès l'adolescence.
\end{attention}

\courssub{Les maladies cardiovasculaires : des artères qui se bouchent}

\textbf{Observation.} Les infarctus et les accidents vasculaires cérébraux (AVC) frappent de plus en plus de Camerounais, parfois jeunes. Ces accidents surviennent brutalement, mais ils sont le résultat d'un processus silencieux commencé des années plus tôt.

\begin{definition}[L'athérosclérose]
L'\cle{athérosclérose} est une maladie des artères caractérisée par le dépôt de \cle{cholestérol} (principalement LDL-cholestérol) et de graisses dans la paroi artérielle (intima), formant une \cle{plaque d'athérome}. Cette plaque rétrécit le calibre de l'artère (sténose) et diminue le flux sanguin. Elle peut se rompre et provoquer un caillot (thrombus) qui bouche totalement l'artère (thrombose).
\end{definition}

\textbf{Conséquences selon l'artère atteinte :}
\begin{itemize}
\item Artère coronaire (cœur) bouchée $\rightarrow$ \cle{infarctus du myocarde} (nécrose d'une partie du muscle cardiaque par ischémie) ;
\item Artère cérébrale bouchée (ischémie) ou rompue (hémorragie) $\rightarrow$ \cle{AVC} (paralysie, troubles du langage, décès) ;
\item Artères des membres inférieurs $\rightarrow$ artérite des membres inférieurs (douleurs à la marche, risque de gangrène).
\end{itemize}

\begin{definition}[L'hypertension artérielle]
L'\cle{hypertension artérielle} (HTA) est une élévation durable de la pression du sang dans les artères (au-delà de \SI{140}{\mmHg} pour la systolique et/ou \SI{90}{\mmHg} pour la diastolique). Favorisée par l'excès de sel, le surpoids et la sédentarité, elle fatigue le cœur (hypertrophie ventriculaire gauche) et fragilise les parois artérielles, multipliant le risque d'infarctus et d'AVC. On l'appelle le « tueur silencieux » car elle ne provoque souvent aucun symptôme.
\end{definition}

\begin{exemplebox}[Un chiffre camerounais qui alerte]
L'hypertension artérielle touche environ \textbf{1 adulte sur 3 au Cameroun}, souvent favorisée par l'excès de sel (cubes d'assaisonnement, poisson fumé très salé) et par le surpoids. Beaucoup d'hypertendus l'ignorent : d'où l'importance de faire mesurer régulièrement sa tension.
\end{exemplebox}

\begin{popfigure}
\begin{center}
\begin{tikzpicture}[scale=1]
% Artère saine (gauche)
\draw[fill=poppinkL,draw=poppink,thick] (0,0) circle (1.5);
\draw[fill=white,draw=poppink,thick] (0,0) circle (1.1);
\draw[->,popblue,thick] (0,0) -- (0.7,0.55);
\node[popdark] at (0,-2.2) {\textbf{Artère saine}};
\node[popdark,align=center,font=\small] at (0,2.1) {Sang circulant librement};
% Légendes artère saine
\draw[thin,popdark] (-1.3,0.4) -- (-2.7,1.1) node[left,align=right,font=\small] {Paroi artérielle\\ (intima + média)};
\draw[thin,popdark] (0.8,-0.8) -- (2.4,-1.5) node[right,align=left,font=\small] {Lumière artérielle\\ (passage du sang)};

% Artère athéromateuse (droite)
\begin{scope}[xshift=7cm]
\draw[fill=poppinkL,draw=poppink,thick] (0,0) circle (1.5); % Adventice + Média externe
\draw[fill=popgoldL,draw=popgold,thick] (0,0) circle (1.1); % Plaque d'athérome (intima épaissie)
\draw[fill=white,draw=popgold,thick] (0,0) circle (0.45); % Lumière rétrécie
\draw[->,popblue,thick] (0,0) -- (0.28,0.35); % Flux réduit
\node[popdark] at (0,-2.2) {\textbf{Artère athéromateuse}};
\node[popdark,align=center,font=\small] at (0,2.1) {Passage du sang rétréci};
% Légendes pointées artère malade
\draw[thin,popdark] (1.1,-0.7) -- (2.7,-1.3) node[right,align=left,font=\small] {Plaque d'athérome\\(dépôt cholestérol, cellules)};
\draw[thin,popdark] (-1.4,0.4) -- (-2.8,1.1) node[left,align=right,font=\small] {Paroi artérielle\\ (adventice/média)};
\end{scope}
\end{tikzpicture}
\end{center}
\legende{Comparaison entre une artère saine (à gauche, lumière large, paroi fine) et une artère rétrécie par une plaque d'athérome (à droite, lumière réduite, paroi épaissie par le dépôt de cholestérol). L'athérosclérose diminue le flux sanguin et peut se compliquer de thrombose (infarctus, AVC).}
\end{popfigure}

\begin{savaistu}[Une maladie qui commence à ton âge]
La plaque d'athérome se constitue \textbf{silencieusement dès l'adolescence} lorsque l'alimentation est trop riche en graisses saturées (fritures répétées, viandes grasses) et en sucres rapides. Des études d'autopsies de jeunes adultes (notamment soldats tués au combat) ont montré des stries graisseuses dans les artères dès 15-20 ans. L'infarctus de l'adulte de 50 ans est souvent le résultat des habitudes de l'adolescent de 15 ans. Choisir son alimentation aujourd'hui, c'est protéger ses artères pour dans 40 ans !
\end{savaistu}

\courssub{Les autres conséquences : un corps et un moral mis à l'épreuve}

La sédentarité et la mauvaise alimentation ne se contentent pas de provoquer diabète et maladies cardiovasculaires. Leurs effets sont multiples et touchent l'ensemble de l'organisme :
\begin{itemize}
\item \textbf{Troubles articulaires} : l'excès de masse surcharge les genoux, les hanches et la colonne vertébrale (rachis lombaire), provoquant douleurs mécaniques et arthrose précoce (gonarthrose, coxarthrose) ;
\item \textbf{Apnées du sommeil} : la graisse du cou (dépôt parapharyngé) comprime les voies respiratoires supérieures, provoquant des arrêts respiratoires nocturnes (SAOS), un sommeil non réparateur et une somnolence diurne excessive ;
\item \textbf{Fatigue chronique et intolérance à l'effort} : le cœur (déconditionnement cardiovasculaire) et les muscles (fonte musculaire, sarcopénie relative), peu entraînés, s'épuisent vite au moindre effort ;
\item \textbf{Retentissement psychologique} : moqueries (grossophobie), isolement social, perte d'estime de soi, anxiété, parfois dépression ;
\item \textbf{Retentissement scolaire} : fatigue, difficultés de concentration, absentéisme (maladies répétées, rendez-vous médicaux) dégradent les résultats scolaires.
\end{itemize}

\begin{popfigure}
\begin{center}
\begin{tikzpicture}[
grow=right,
level 1/.style={sibling distance=3cm, level distance=4cm},
level 2/.style={sibling distance=1.8cm, level distance=4.2cm},
every node/.style={draw, rounded corners, align=center, font=\small, inner sep=4pt},
edge from parent/.style={draw=popmuted, thick, ->}
]
\node[fill=poporangeL, draw=poporange, thick, align=center]{Sédentarité\\+ Malnutrition\\par excès}
child { node[fill=poppinkL, draw=poppink, thick, align=center]{Déséquilibre\\énergétique durable\\$\rightarrow$ Obésité\\(excès masse grasse)}
  child { node[fill=popblueL, draw=popblue, align=center]{Troubles\\mécaniques\\Articulations, Apnées\\du sommeil} }
  child { node[fill=popblueL, draw=popblue, align=center]{Maladies\\métaboliques\\Diabète type 2,\\Dyslipidémie, Stéatose} }
  child { node[fill=popblueL, draw=popblue, align=center]{Maladies\\cardiovasculaires\\Athérosclérose, HTA,\\Infarctus, AVC} }
  child { node[fill=poppurpleL, draw=poppurple, align=center]{Retentissements\\psychosociaux\\et scolaires} }
};
\end{tikzpicture}
\end{center}
\legende{Arbre des conséquences : la sédentarité associée à la malnutrition par excès crée un déséquilibre énergétique durable, source d'obésité, elle-même à l'origine de pathologies métaboliques (diabète type 2), cardiovasculaires (athérosclérose, HTA), mécaniques et psychosociales.}
\end{popfigure}

\courssub{Facteurs aggravants en milieu urbain camerounais}

Pourquoi ces maladies progressent-elles si vite dans nos villes ? Plusieurs changements récents du mode de vie urbain jouent un rôle majeur :
\begin{itemize}
\item \textbf{Les fritures de rue omniprésentes} : beignets, puff-puff, pommes frites, « koki » très huilé, poisson braisé, vendus partout et bon marché, apportent des graisses saturées et des acides gras trans en excès ;
\item \textbf{Les boissons sucrées industrialisées} : sodas et jus « néctar » consommés quotidiennement apportent des sucres rapides (saccharose, sirop de glucose-fructose) sans pouvoir rassasiant (absence de fibres) ;
\item \textbf{Les écrans prolongés} : téléphone, télévision et jeux vidéo remplacent les jeux de plein air ; certains adolescents passent plus de \SI{6}{h} par jour assis (temps sédentaire) ;
\item \textbf{La réduction drastique de la marche quotidienne} : motos-taxis (benskin) et voitures remplacent la marche active, même pour de courts trajets (\textless \SI{1}{km}) ;
\item \textbf{L'excès de sel caché} : cubes d'assaisonnement (glutamate, sel), poisson fumé/salé, sauces industrielles favorisent la rétention d'eau et l'hypertension artérielle.
\end{itemize}

\begin{exoresolu}[Analyser une situation de vie]
\textbf{Énoncé.} Rodrigue, 17 ans, en classe de Première C à Douala, mesure \SI{1,70}{m} et pèse \SI{92}{kg}. Il prend la moto-taxi pour \SI{800}{m}, boit 2 sodas par jour, déjeune souvent de fritures et passe \SI{5}{h} par jour sur son téléphone. Sa glycémie à jeun est de \SI{1,30}{g/L}. 
\begin{enumerate}
\item Calcule son IMC et interprète-le.
\item Identifie les facteurs de risque présents.
\item Quelle maladie menace déjà Rodrigue ? Justifie.
\end{enumerate}
\tcblower
\textbf{Solution.}
\begin{enumerate}
\item $\text{IMC} = \dfrac{92}{1,70^2} = \dfrac{92}{2,89} \approx 31,8$. Un IMC supérieur à 30 correspond à une \textbf{obésité} (classe I).
\item Facteurs de risque identifiés : 
\begin{itemize}
    \item \textbf{Sédentarité marquée} : moto-taxi pour \SI{800}{m} (supprime marche active), \SI{5}{h}/jour écrans (temps assis), absence de sport déclaré.
    \item \textbf{Malnutrition par excès} : 2 sodas/jour $\approx$ \SI{280}{kcal} de sucres rapides (sans fibres) ; fritures riches en graisses saturées et calories vides ; excès de sel probable.
\end{itemize}
\item Sa glycémie à jeun de \SI{1,30}{g/L} est supérieure au seuil diagnostique de \SI{1,26}{g/L} (confirmé sur deux dosages) : Rodrigue présente une \textbf{hyperglycémie chronique}, signe d'un diabète de type 2 débutant. Son obésité abdominale (viscérale) a favorisé une résistance à l'insuline. Sans changement de mode de vie, il s'expose aux complications micro/macro-vasculaires (yeux, reins, nerfs, artères) et à l'aggravation de l'obésité. La bonne nouvelle : à son âge, une alimentation équilibrée (réduction sucres/graisses) et une activité physique régulière (\textgreater \SI{60}{min}/jour) peuvent encore améliorer la sensibilité à l'insuline et faire régresser l'hyperglycémie — ce sera l'objet de la prochaine section.
\end{enumerate}
\end{exoresolu}

\begin{aretenir}[L'essentiel de la section]
\begin{itemize}
\item La \cle{sédentarité} (moins de \SI{30}{min} d'activité modérée/jour pour un adulte, \textless \SI{60}{min} pour un ado) et la \cle{malnutrition par excès} créent un déséquilibre durable entre apports et dépenses énergétiques.
\item L'excédent est stocké dans le \cle{tissu adipeux} (adipocytes) : un excès de \SI{100}{kcal}/jour $\approx$ 4 à 5 kg de graisse par an.
\item L'obésité (surtout androïde/abdominale) favorise le \cle{diabète de type 2} (résistance à l'insuline $\rightarrow$ hyperglycémie chronique $\geq$ \SI{1,26}{g/L}, complications : rétinopathie, néphropathie, neuropathie).
\item Le cholestérol (LDL) se dépose dans l'intima artérielle (\cle{athérosclérose}, \cle{plaque d'athérome}) ; avec l'\cle{hypertension artérielle} (\textgreater \SI{140/90}{mmHg}), c'est le risque majeur d'\cle{infarctus du myocarde} et d'\cle{AVC}.
\item S'y ajoutent troubles articulaires (arthrose), apnées du sommeil (SAOS), fatigue, retentissements psychologiques (estime de soi) et scolaires.
\item En milieu urbain camerounais, fritures de rue, boissons sucrées, écrans, réduction de la marche et excès de sel (cubes) aggravent le risque précocement.
\end{itemize}
\end{aretenir}

\coursec{Recommandations : alimentation et activité physique}

Dans la section précédente, nous avons vu que la sédentarité et la mauvaise alimentation sont des facteurs de risque majeurs : elles favorisent l'obésité, le diabète de type 2 et les maladies cardiovasculaires, parfois dès l'adolescence. La bonne nouvelle, c'est que ces maladies liées au mode de vie sont en grande partie \cle{évitables}. Il ne suffit pas de savoir ce qui rend malade : il faut savoir \emph{quoi faire}, concrètement, chaque jour, pour rester en bonne santé. C'est précisément l'objet de cette section : transformer les connaissances en recommandations pratiques, mesurables et applicables dans la vie quotidienne d'un élève camerounais.

\begin{exemplebox}[Une situation pour ouvrir la réflexion]
Deux camarades de Première, Amina et Brice, ont le même âge (16 ans). Amina marche 25 minutes chaque matin pour aller au lycée, joue au handball le mercredi et le samedi, mange un petit-déjeuner avant de partir et prend des fruits au goûter. Brice se fait déposer en moto-taxi, passe ses soirées devant son téléphone (environ 4 h par jour), saute souvent le petit-déjeuner et grignote des beignets avec des boissons sucrées. Au bout de quelques mois, Brice se plaint de fatigue en classe et prend du poids. \textbf{Question :} quelles recommandations simples permettraient à Brice d'inverser la tendance ? C'est à cette question que toute cette section répond.
\end{exemplebox}

\courssub{Les recommandations de l'OMS en matière d'activité physique}

\textbf{Une observation de départ.} Dans les villes camerounaises, les occasions spontanées de bouger diminuent : moto-taxis pour de courts trajets, écrans omniprésents, jeux de rue remplacés par les jeux vidéo. Or le corps humain est « construit » pour le mouvement : le cœur, les muscles, les os et même le cerveau ont besoin d'une sollicitation régulière pour fonctionner correctement. L'Organisation mondiale de la santé (OMS) a donc établi, à partir de centaines d'études épidémiologiques, des recommandations chiffrées selon l'âge.

\begin{definition}[Activité physique modérée et activité physique soutenue]
Une \cle{activité physique} est tout mouvement produit par les muscles squelettiques qui augmente la dépense d'énergie au-dessus du niveau de repos. On distingue :
\begin{itemize}
\item une activité \cle{modérée} : elle accélère nettement la respiration et le rythme cardiaque, mais \textbf{permet encore de parler} sans essoufflement excessif. Exemple : la marche rapide ;
\item une activité \cle{soutenue} (ou intense) : la respiration devient forte et \textbf{parler devient difficile}. Exemples : la course, un match de football disputé, la corde à sauter rapide.
\end{itemize}
Ce « test de la parole » est un moyen simple, sans appareil, d'évaluer l'intensité de son effort.
\end{definition}

\begin{propriete}[Recommandations de l'OMS pour les 5--17 ans]
Pour les enfants et adolescents de 5 à 17 ans, l'OMS recommande :
\begin{itemize}
\item \textbf{au moins 60 minutes par jour} d'activité physique d'intensité modérée à soutenue ;
\item des activités de \cle{renforcement musculaire} et de \cle{charge osseuse} (sauts, course, jeux de ballon, gymnastique) au moins \textbf{3 fois par semaine} ;
\item la limitation du temps passé en position assise, en particulier \textbf{moins de 2 heures par jour} devant les écrans de loisir (téléphone, télévision, jeux vidéo).
\end{itemize}
Pour les adultes (18--64 ans), la recommandation est d'au moins 150 minutes d'activité modérée par semaine, soit environ 30 minutes par jour en moyenne. Les besoins quotidiens en mouvement sont donc \textbf{deux fois plus élevés chez l'adolescent} que chez l'adulte, car la croissance du squelette, des muscles et du cerveau exige une stimulation régulière.
\end{propriete}

\begin{popfigure}
\begin{center}
\begin{tikzpicture}
\begin{axis}[
ybar, bar width=0.9cm,
width=0.85\linewidth, height=6.5cm,
ylabel={Temps d'activité physique quotidien (min)},
symbolic x coords={Enfant (5--12 ans), Adolescent (13--17 ans), Adulte (18--64 ans)},
xtick=data,
ymin=0, ymax=70,
ytick={0,10,20,30,40,50,60,70},
nodes near coords,
nodes near coords style={font=\small\bfseries, color=popdark},
ymajorgrids=true, grid style={popline!40},
axis line style={popdark},
ylabel style={font=\small},
xticklabel style={font=\small},
]
\addplot[fill=popblue!70, draw=popblue] coordinates {(Enfant (5--12 ans),60)};
\addplot[fill=poporange!75, draw=poporange] coordinates {(Adolescent (13--17 ans),60)};
\addplot[fill=popgreen!70, draw=popgreen] coordinates {(Adulte (18--64 ans),30)};
\end{axis}
\end{tikzpicture}
\end{center}
\legende{Temps minimal quotidien d'activité physique modérée à soutenue recommandé par l'OMS selon la tranche d'âge. Enfants et adolescents : 60 min/jour ; adultes : environ 30 min/jour (150 min/semaine).}
\end{popfigure}

\textbf{Quelles activités comptent ?} Bonne nouvelle : il n'est pas nécessaire d'être licencié dans un club sportif. De nombreuses activités de la vie quotidienne sont de véritables activités physiques :
\begin{itemize}
\item la \textbf{marche rapide} (aller au lycée ou au marché à pied d'un pas soutenu) ;
\item la \textbf{course}, le \textbf{football}, le handball, le basket-ball, l'athlétisme ;
\item la \textbf{danse} (traditionnelle ou moderne), la corde à sauter ;
\item les \textbf{travaux domestiques et champêtres actifs} : balayage énergique, port d'eau, jardinage, travaux des champs ;
\item les jeux de plein air entre camarades.
\end{itemize}

\begin{methode}[Atteindre 60 minutes par jour : les minutes s'additionnent]
L'objectif de 60 minutes peut sembler élevé, mais il n'est pas nécessaire de les faire d'un seul coup. Les durées s'\textbf{additionnent} sur la journée :
\begin{enumerate}
\item Repère les occasions de bouger déjà présentes dans ta journée (trajet maison--lycée, récréations actives, tâches domestiques).
\item Ajoute une ou deux séances courtes d'activité plus intense (sport, danse, corde à sauter).
\item Additionne les durées : dès que le total atteint 60 minutes d'intensité modérée à soutenue, l'objectif quotidien est atteint.
\item Deux à trois fois par semaine, inclus des activités de renforcement (sauts, pompes adaptées, jeux de ballon avec accélérations).
\end{enumerate}
\end{methode}

\begin{exemplebox}[Un exemple chiffré : la journée d'Amina]
30 min de marche rapide (trajets aller-retour au lycée) + 20 min de football à la récréation + 10 min de corde à sauter le soir = \textbf{60 minutes}. Objectif quotidien atteint, sans matériel coûteux ni inscription en salle de sport.
\end{exemplebox}

\begin{attention}[Limiter le temps sédentaire, même si l'on fait du sport]
Faire une heure de sport ne « neutralise » pas 6 heures passées assis devant un écran. La \cle{sédentarité} (position assise ou allongée prolongée en état de veille) est un facteur de risque \emph{indépendant} : elle ralentit le métabolisme des graisses et favorise la résistance à l'insuline. Règle pratique : \textbf{moins de 2 h par jour} d'écrans de loisir, et se lever pour bouger quelques minutes toutes les heures.
\end{attention}

\courssub{Les recommandations alimentaires}

\textbf{Observation.} L'alimentation moderne en milieu urbain s'est transformée : boissons sucrées, beignets, fritures, plats très salés ou très gras sont devenus faciles d'accès, tandis que la consommation de fruits, de légumes et d'eau diminue. Pourtant, les principes d'une alimentation protectrice sont simples et compatibles avec la cuisine camerounaise traditionnelle, souvent très saine (légumes-feuilles, tubercules peu raffinés, poissons, fruits de saison).

\begin{propriete}[Les grandes recommandations alimentaires]
\begin{itemize}
\item Consommer \textbf{au moins 5 fruits et légumes par jour} (environ 400 g au total) : ils apportent fibres, vitamines, minéraux et composés protecteurs, avec peu de calories.
\item \textbf{Limiter les sucres rapides} (boissons sucrées, bonbons, pâtisseries) : ils provoquent des pics de glycémie et favorisent la prise de poids et le diabète de type 2.
\item \textbf{Réduire le sel à moins de 5 g par jour} (environ une cuillère à café, tous aliments confondus, y compris les bouillons cubes) : l'excès de sel favorise l'hypertension artérielle.
\item \textbf{Limiter les graisses saturées} (fritures répétées, viandes très grasses, huile de palme en excès) et privilégier de bonnes sources de lipides en quantité modérée.
\item \textbf{Privilégier l'eau} comme boisson principale (environ 1,5 à 2 L par jour, davantage sous climat chaud) : l'eau n'apporte aucune calorie et aucun sucre.
\end{itemize}
\end{propriete}

\begin{popfigure}
\begin{center}
\begin{tikzpicture}[scale=1.05]
% Assiette : disque divisé
\fill[popgreen!55] (90:2.6) arc (90:270:2.6) -- (0,0) -- cycle;
\fill[popgold!65] (90:2.6) arc (90:0:2.6) -- (0,0) -- cycle;
\fill[poporange!70] (0:2.6) arc (0:-90:2.6) -- (0,0) -- cycle;
\draw[popdark, thick] (0,0) circle (2.6);
\draw[popdark, thick] (0,2.6) -- (0,-2.6);
\draw[popdark, thick] (0,0) -- (2.6,0);
% Légendes internes
\node[font=\small\bfseries, text=popdark, align=center] at (-1.5,0.9) {Légumes et};
\node[font=\small\bfseries, text=popdark, align=center] at (-1.5,0.3) {fruits};
\node[font=\small\bfseries, text=popdark, align=center] at (-1.5,-0.4) {(1/2)};
\node[font=\small\bfseries, text=popdark, align=center] at (1.5,0.9) {Féculents};
\node[font=\small\bfseries, text=popdark, align=center] at (1.5,0.3) {(1/4)};
\node[font=\small\bfseries, text=popdark, align=center] at (1.5,-0.9) {Protéines};
\node[font=\small\bfseries, text=popdark, align=center] at (1.5,-1.5) {(1/4)};
% Verre d'eau
\draw[popblue, thick] (4.1,-0.6) -- (4.35,1.1) -- (5.25,1.1) -- (5.5,-0.6) -- cycle;
\fill[popblue!35] (4.2,-0.4) -- (4.4,0.7) -- (5.2,0.7) -- (5.4,-0.4) -- cycle;
\node[font=\small\bfseries, text=popblue!80!black, align=center] at (4.8,-1.2) {Eau};
% Pointeurs
\draw[thin, popmuted] (-1.5,1.6) -- (-1.5,2.9) node[above, font=\footnotesize, text=popdark, align=center] {ndolè, folong, gombo,\\carottes, mangues, bananes...};
\draw[thin, popmuted] (1.9,1.9) -- (3.1,2.9) node[above, font=\footnotesize, text=popdark, align=center] {riz, igname, manioc,\\maïs, plantain...};
\draw[thin, popmuted] (1.9,-1.9) -- (3.1,-2.9) node[below, font=\footnotesize, text=popdark, align=center] {poisson, viande maigre,\\œufs, haricots, arachides...};
\end{tikzpicture}
\end{center}
\legende{Le modèle de l'« assiette équilibrée » : la moitié de l'assiette est occupée par des légumes et des fruits, un quart par des féculents (de préférence peu raffinés), un quart par des protéines. L'eau est la boisson de référence.}
\end{popfigure}

\begin{exemplebox}[Adapter l'assiette équilibrée au repas camerounais]
Un plat de ndolè avec du miondo peut être rééquilibré : grande portion de ndolè (légumes-feuilles), portion modérée de miondo ou de plantain (féculent), morceau raisonnable de poisson ou de viande (protéines), et un fruit de saison en dessert. Un verre d'eau accompagne le repas à la place d'une boisson sucrée.
\end{exemplebox}

\textbf{La régularité des repas et le petit-déjeuner.} Sauter des repas, en particulier le \cle{petit-déjeuner}, est une erreur fréquente chez les adolescents. Après la nuit (8 à 10 h sans manger), les réserves de glucose de l'organisme sont basses : or le cerveau consomme du glucose en permanence. Un élève qui arrive en classe à jeun présente une baisse de concentration, de la fatigue et des fringales qui le poussent à grignoter des aliments très sucrés ou gras en milieu de matinée. Un petit-déjeuner équilibré combine : un féculent (pain, bouillie de maïs ou de mil), une source de protéines (œuf, lait, arachides), un fruit, et une boisson non sucrée. Trois repas réguliers, à horaires à peu près fixes, évitent les compulsions alimentaires et stabilisent la glycémie au cours de la journée.

\begin{attention}[Le sport ne compense pas une mauvaise alimentation]
On entend parfois : « je peux manger n'importe quoi, je fais du sport ». C'est faux. Le sport ne compense \textbf{pas durablement} une alimentation déséquilibrée : il est beaucoup plus facile d'avaler 500 kcal de fritures et de boisson sucrée en 10 minutes que de les dépenser (il faudrait plus d'une heure de course !). De plus, un excès de sel ou de graisses saturées nuit aux artères même chez un sportif. \textbf{Les deux leviers — alimentation équilibrée et activité physique — doivent agir ensemble.}
\end{attention}

\courssub{Les bénéfices de l'activité physique : une démarche d'investigation}

\textbf{Hypothèse.} L'activité physique régulière améliore la santé sur plusieurs plans à la fois. Comment le vérifier ?

\begin{experience}[Observer les effets d'une séance d'activité modérée]
\textbf{Matériel :} un chronomètre (ou téléphone), un carnet de suivi.
\textbf{Protocole :}
\begin{enumerate}
\item Au repos, assis depuis 5 minutes, mesure ta fréquence cardiaque (nombre de battements en 1 minute, deux doigts sur le poignet). Note-la.
\item Effectue 10 minutes de marche rapide ou de corde à sauter. Mesure immédiatement ta fréquence cardiaque, puis à nouveau 5 minutes après l'arrêt.
\item Répète l'expérience chaque semaine pendant 6 semaines de pratique régulière (au moins 3 séances par semaine).
\end{enumerate}
\textbf{Observations attendues :} la fréquence cardiaque augmente pendant l'effort (le cœur pompe plus vite pour oxygéner les muscles), puis redescend. Au fil des semaines, la fréquence au repos \emph{diminue} et le retour au calme après l'effort est \emph{plus rapide}.
\textbf{Interprétation :} le cœur est un muscle ; entraîné, il devient plus puissant et plus efficace : il pompe plus de sang à chaque battement (volume d'éjection systolique accru), donc bat moins vite au repos. C'est un indice direct de l'amélioration de la santé cardiovasculaire.
\textbf{Conclusion :} l'activité physique régulière renforce le cœur et améliore la condition physique de façon mesurable.
\end{experience}

Au-delà du cœur, les bénéfices d'une pratique régulière sont multiples et documentés :
\begin{itemize}
\item \textbf{contrôle du poids} : l'activité physique augmente la dépense énergétique et préserve la masse musculaire, limitant l'accumulation de graisse ;
\item \textbf{santé du cœur et des vaisseaux} : baisse de la tension artérielle, amélioration du profil lipidique, meilleure sensibilité à l'insuline (protection contre le diabète de type 2) ;
\item \textbf{meilleure concentration scolaire} : l'effort augmente l'oxygénation du cerveau et stimule la sécrétion de molécules favorisant la mémoire et l'attention ; les élèves actifs réussissent en moyenne mieux à l'école ;
\item \textbf{bien-être psychologique} : l'exercice libère des \cle{endorphines} (molécules du plaisir et de l'apaisement), réduit le stress et l'anxiété, améliore le sommeil et l'estime de soi.
\end{itemize}

\begin{savaistu}[Combien de calories brûle-t-on en bougeant ?]
60 minutes d'activité modérée (marche rapide, vélo tranquille, danse) dépensent environ \textbf{200 à 300 kcal} selon le poids et l'intensité — soit l'équivalent énergétique d'un beignet et d'un verre de boisson sucrée. Cela montre deux choses : bouger chaque jour fait une vraie différence sur le long terme ; mais il est plus facile de \emph{ne pas consommer} ces calories superflues que de les éliminer ensuite. D'où l'importance de combiner les deux leviers !
\end{savaistu}

\begin{exoresolu}[Construire la journée santé de Brice]
\textbf{Énoncé.} Brice, 16 ans, passe 4 h par jour devant son téléphone, saute le petit-déjeuner, grignote 2 beignets et une boisson sucrée chaque matinée, et ne pratique aucune activité physique. Propose-lui un programme d'une journée type respectant les recommandations de l'OMS, et justifie chaque modification.
\tcblower
\textbf{Solution détaillée.}
\begin{enumerate}
\item \textbf{Activité physique (objectif : 60 min).} 20 min de marche rapide pour aller au lycée + 20 min de retour à pied + 20 min de football à la récréation = 60 min d'intensité modérée à soutenue. Le mercredi et le samedi, il ajoute des jeux de ballon avec accélérations (renforcement musculaire, 3 fois par semaine).
\item \textbf{Temps d'écran.} Il ramène son temps d'écran de loisir de 4 h à moins de 2 h par jour, en remplaçant une partie par des activités actives ou la lecture.
\item \textbf{Petit-déjeuner.} Bouillie de maïs + un œuf + une banane + un verre d'eau : féculent, protéines, fruit, boisson non sucrée. Fini les fringales de 10 h.
\item \textbf{Alimentation quotidienne.} Il remplace les beignets et la boisson sucrée par un fruit et de l'eau ; il veille à consommer 5 fruits et légumes par jour et limite les bouillons cubes (sel $<$ 5 g/jour).
\end{enumerate}
\textbf{Justification globale :} chaque modification attaque un facteur de risque identifié (sédentarité, excès de sucres rapides, excès de sel, repas sautés). Combinées, elles réduisent le risque de prise de poids, de diabète de type 2 et de maladies cardiovasculaires, tout en améliorant la concentration en classe.
\end{exoresolu}

\begin{aretenir}[L'essentiel de la section]
\begin{itemize}
\item Adolescents (5--17 ans) : \textbf{au moins 60 min/jour} d'activité physique modérée à soutenue, avec renforcement musculaire 3 fois par semaine ; les minutes s'additionnent.
\item Intensité : modérée = on peut encore parler ; soutenue = parler devient difficile.
\item Moins de \textbf{2 h/jour} d'écrans de loisir ; se lever régulièrement.
\item Alimentation : \textbf{5 fruits et légumes par jour}, sel $<$ 5 g/jour, sucres rapides et graisses saturées limités, \textbf{eau} comme boisson principale, assiette = moitié légumes-fruits, quart féculents, quart protéines.
\item Ne jamais sauter le \textbf{petit-déjeuner} ; repas réguliers.
\item Activité physique et alimentation équilibrée sont \textbf{complémentaires} : l'un ne compense pas l'autre.
\end{itemize}
\end{aretenir}

\begin{exercice}[À toi de jouer]
Pendant trois jours, note toutes tes activités physiques (nature et durée) et ton temps d'écran de loisir. Calcule ta moyenne quotidienne d'activité physique. Atteins-tu les 60 minutes recommandées ? Sinon, propose deux modifications concrètes et réalistes de ton emploi du temps pour y parvenir, en précisant l'intensité (modérée ou soutenue) de chaque activité choisie.
\end{exercice}

\coursec{Élaborer un programme de prévention}

Dans la section précédente, nous avons établi les recommandations officielles : une alimentation variée répartie selon la règle 50--55 \% de glucides, 30--35 \% de lipides et 12--15 \% de protéines, et au moins 60 minutes d'activité physique par jour pour l'adolescent. Mais connaître les règles ne suffit pas : comment passer du savoir à l'action ? Comment un élève de Première de Yaoundé ou de Garoua, pris entre les cours, les devoirs et les tentations des beignets-haricots du matin, peut-il transformer ces recommandations en habitudes durables ? C'est toute la question de cette dernière section : \cle{élaborer un programme de prévention}, c'est-à-dire un plan d'action personnel, chiffré, progressif et suivi dans le temps.

\courssub{La démarche générale en quatre étapes}

\textbf{Situation concrète.} Junior, 16 ans, en classe de Première C au lycée de Nkongsamba, se sent essoufflé à l'issue du cross scolaire. Sa taille est de \SI{1,70}{m} pour \SI{82}{kg}, soit un IMC de $82/1{,}70^2 = 82/2{,}89 \approx 28{,}4$. Rapportée aux courbes de corpulence de l'OMS pour les adolescents, cette valeur le situe en \textbf{surpoids}. Sa tante lui conseille de « ne plus rien manger le soir » ; son camarade lui propose un régime « sans féculents » trouvé sur Internet. Que doit-il faire ? La démarche scientifique, ici, s'applique à sa propre santé : on ne décide jamais avant d'avoir mesuré.

\begin{methode}[Les quatre étapes d'un programme de prévention]
Tout programme de prévention efficace suit quatre étapes, dans l'ordre :
\begin{enumerate}
\item \textbf{Le bilan chiffré} : on mesure l'IMC, on évalue la ration alimentaire actuelle (nature et quantité des aliments consommés sur une journée type) et on estime le temps d'activité physique réel (en minutes par jour). Aucune décision ne précède ce diagnostic.
\item \textbf{La fixation d'objectifs réalistes} : on définit des objectifs précis, modérés et datés (méthode SMART, voir plus loin).
\item \textbf{L'élaboration du plan hebdomadaire} : on construit un menu type et un emploi du temps sportif concrets, jour par jour.
\item \textbf{Le suivi et l'ajustement} : on contrôle régulièrement (pesée mensuelle, carnet de bord) et on corrige le plan si nécessaire.
\end{enumerate}
\end{methode}

\begin{popfigure}
\begin{center}
\begin{tikzpicture}[
  etape/.style={rectangle, rounded corners=3pt, draw=popblue, fill=popblueL, thick, text width=4.6cm, align=center, minimum height=1.6cm, font=\small},
  fleche/.style={-{Stealth[length=3mm]}, thick, popdark}
]
\node[etape, align=center] (bilan) at (0,0){\textbf{1. Bilan}\\ IMC, ration alimentaire, temps d'activité actuel};
\node[etape, align=center] (obj) at (6.2,0){\textbf{2. Objectifs}\\ réalistes, mesurables, datés (SMART)};
\node[etape, align=center] (plan) at (6.2,-2.6){\textbf{3. Plan hebdomadaire}\\ menus équilibrés + séances d'activité physique};
\node[etape, align=center] (suivi) at (0,-2.6){\textbf{4. Suivi}\\ pesée mensuelle, carnet de bord, ajustements};
\draw[fleche] (bilan) -- (obj);
\draw[fleche] (obj) -- (plan);
\draw[fleche] (plan) -- (suivi);
\draw[fleche] (suivi) -- (bilan) node[midway, above, font=\small\itshape, text=popmuted]{réajustement};
\end{tikzpicture}
\end{center}
\legende{Organigramme de la démarche en quatre étapes. Le retour de l'étape 4 vers l'étape 1 montre que le programme est cyclique : le suivi permet de réévaluer le bilan et d'ajuster les objectifs.}
\end{popfigure}

\begin{attention}[Commencer par mesurer, pas par décider]
L'erreur la plus fréquente consiste à sauter l'étape du bilan : se lancer dans un régime « au hasard », copier le programme d'un ami ou suivre un conseil entendu à la radio. Or un programme efficace pour un élève en surpoids serait inadapté, voire dangereux, pour un élève déjà mince ou pour un sportif de haut niveau. \textbf{Toute décision doit s'appuyer sur des données chiffrées : IMC, apports énergétiques estimés, minutes d'activité réelles.} Rappelons aussi que chez l'adolescent, l'IMC s'interprète sur les courbes de corpulence de l'OMS en fonction de l'âge et du sexe, et non avec les seuls seuils de l'adulte (18,5 ; 25 ; 30) : en cas de doute, on se réfère au personnel de santé.
\end{attention}

\courssub{Fixer des objectifs SMART}

Un objectif vague (« je veux maigrir », « je vais faire du sport ») échoue presque toujours, car on ne peut ni le suivre ni vérifier qu'il est atteint. La méthode \cle{SMART} impose cinq critères à tout bon objectif :

\begin{definition}[Objectif SMART]
Un objectif est dit \cle{SMART} s'il est :
\begin{itemize}
\item \textbf{S}pécifique : il précise exactement l'action attendue (« marcher 40 min » et non « bouger plus ») ;
\item \textbf{M}esurable : il comporte une grandeur chiffrée (minutes, kilogrammes, portions) ;
\item \textbf{A}tteignable : il tient compte des contraintes réelles (emploi du temps, matériel disponible) ;
\item \textbf{R}éaliste : il reste modéré et compatible avec la santé (pas de perte de poids brutale) ;
\item \textbf{T}emporellement défini (daté) : il fixe une échéance (« d'ici fin décembre »).
\end{itemize}
\end{definition}

\begin{exemplebox}[Un objectif SMART pour Junior]
Objectif vague : « perdre du poids ». Objectif SMART : « perdre \SI{0,5}{kg} par mois pendant 6 mois, en supprimant les boissons sucrées (qui apportaient environ \SI{600}{kJ} par jour, soit l'équivalent d'une bouteille de soda de \SI{33}{cL}) et en marchant 40 minutes chaque jour entre la maison et le lycée ». Vérifions : l'objectif est spécifique (deux actions précises), mesurable (\SI{0,5}{kg}/mois, 40 min/jour), atteignable (la marche domicile-lycée existe déjà), réaliste (\SI{0,5}{kg}/mois est une perte douce, sans danger pour la croissance) et daté (6 mois).
\end{exemplebox}

\begin{attention}[Danger des régimes drastiques chez l'adolescent]
Les régimes très restrictifs, le jeûne prolongé ou la suppression totale d'une famille d'aliments (féculents, graisses) sont \textbf{dangereux chez l'adolescent en pleine croissance} : ils provoquent des carences (fer, calcium, protéines), une fonte musculaire, des troubles de la concentration en classe, et favorisent à terme une reprise de poids rapide (effet « yo-yo »). Chez l'adolescent, on ne cherche presque jamais une perte de poids rapide : on privilégie des \textbf{ajustements progressifs} — réduire les boissons sucrées, diminuer les fritures, augmenter les portions de légumes et l'activité physique — qui permettent souvent de stabiliser le poids pendant que la taille augmente, ce qui fait baisser l'IMC naturellement.
\end{attention}

\courssub{Construire le plan hebdomadaire : menus et séances}

\textbf{Principe.} Le plan traduit les objectifs en actions concrètes, jour par jour. Il comporte deux volets indissociables : l'alimentation et l'activité physique.

\textbf{Le volet alimentaire.} On construit un menu journalier à partir d'\textbf{aliments locaux, disponibles et peu coûteux}, en respectant la répartition énergétique 50--55 \% de glucides, 30--35 \% de lipides, 12--15 \% de protéines, et la présence de fruits, de légumes et d'eau à chaque journée. Exemple de journée type :

\begin{center}
\begin{tabularx}{\linewidth}{|l|X|X|}
\hline
\rowcolor{popblueL} \textbf{Repas} & \textbf{Composition (aliments locaux)} & \textbf{Rôle nutritionnel} \\
\hline
Petit-déjeuner & Bouillie de mil enrichie au lait caillé, une banane, un verre d'eau & Glucides lents (énergie durable), protéines et calcium du lait, vitamines et potassium de la banane \\
\hline
Déjeuner & Riz (portion modérée), sauce feuilles (folong) au poisson fumé, une tranche d'avocat, eau & Glucides, protéines et fer (poisson, feuilles), lipides de bonne qualité (avocat) \\
\hline
Goûter & Une orange ou des arachides grillées (petite poignée) & Vitamine C ou protéines et lipides végétaux \\
\hline
Dîner & Couscous de maïs, ndolé aux arachides et crevettes, salade de tomates, eau & Glucides, protéines, fibres et vitamines des légumes \\
\hline
\end{tabularx}
\end{center}

Les pièges à corriger dans ce menu type : limiter l'huile rouge et les fritures (lipides), remplacer les boissons sucrées par de l'eau, et ne pas doubler les portions de féculents (riz + couscous le même jour, par exemple).

\textbf{Le volet activité physique.} L'objectif est d'atteindre en moyenne \textbf{60 minutes par jour}, en cumulant toutes les formes d'activité : marche domicile-lycée, cours d'EPS, jeux de ballon, tâches actives, séances planifiées.

\begin{propriete}[Régularité et bénéfices santé (admise)]
Pour les bénéfices à long terme sur le c\oe ur, le poids et la glycémie, la \cle{régularité} de l'activité physique compte davantage que son intensité : pratiquer une activité modérée \textbf{au moins 5 jours sur 7} est plus protecteur qu'une séance très intense unique suivie de six jours de sédentarité. (Propriété admise ; elle justifie la planification hebdomadaire plutôt que les « défis » ponctuels.)
\end{propriete}

\begin{popfigure}
\begin{center}
\begin{tikzpicture}[font=\scriptsize]
\foreach \j/\jour in {0/Lundi,1/Mardi,2/Mercredi,3/Jeudi,4/Vendredi,5/Samedi,6/Dimanche}{
  \node[draw=popdark, fill=popblueL, minimum width=2.05cm, minimum height=0.55cm, anchor=north west] at (\j*2.1,0) {\textbf{\jour}};
}
\foreach \i/\repas in {0/{Petit-déj.},1/{Déjeuner},2/{Dîner},3/{Activité}}{
  \node[draw=popdark, fill=poporangeL, minimum width=1.9cm, minimum height=1.05cm, anchor=north east, align=center] at (-0.1,-0.6-\i*1.1) {\repas};
}
\foreach \j in {0,...,6}{
  \foreach \i in {0,...,3}{
    \node[draw=popline, fill=white, minimum width=2.05cm, minimum height=1.05cm, anchor=north west, align=center, text width=1.95cm] at (\j*2.1,-0.6-\i*1.1) {};
  }
}
\node[align=center, text width=1.95cm, font=\tiny] at (1.025,-1.125) {Bouillie de mil + banane};
\node[align=center, text width=1.95cm, font=\tiny] at (3.125,-1.125) {Bouillie + \oe uf};
\node[align=center, text width=1.95cm, font=\tiny] at (5.225,-1.125) {Bouillie de mil + banane};
\node[align=center, text width=1.95cm, font=\tiny] at (7.325,-1.125) {Pain + avocat};
\node[align=center, text width=1.95cm, font=\tiny] at (9.425,-1.125) {Bouillie + lait caillé};
\node[align=center, text width=1.95cm, font=\tiny] at (11.525,-1.125) {Beignets (rare) + bouillie};
\node[align=center, text width=1.95cm, font=\tiny] at (13.625,-1.125) {Bouillie + goyave};
\node[align=center, text width=1.95cm, font=\tiny] at (1.025,-2.225) {Riz + folong au poisson};
\node[align=center, text width=1.95cm, font=\tiny] at (3.125,-2.225) {Couscous + ndolé};
\node[align=center, text width=1.95cm, font=\tiny] at (5.225,-2.225) {Fufu + eru};
\node[align=center, text width=1.95cm, font=\tiny] at (7.325,-2.225) {Riz + sauce arachide + poulet};
\node[align=center, text width=1.95cm, font=\tiny] at (9.425,-2.225) {Haricots + banane vapeur};
\node[align=center, text width=1.95cm, font=\tiny] at (11.525,-2.225) {Poulet DG (fête)};
\node[align=center, text width=1.95cm, font=\tiny] at (13.625,-2.225) {Riz + légumes sautés};
\node[align=center, text width=1.95cm, font=\tiny] at (1.025,-3.325) {Igname + sauce graine};
\node[align=center, text width=1.95cm, font=\tiny] at (3.125,-3.325) {Riz + koki};
\node[align=center, text width=1.95cm, font=\tiny] at (5.225,-3.325) {Couscous + okok};
\node[align=center, text width=1.95cm, font=\tiny] at (7.325,-3.325) {Banane malaxée + légumes};
\node[align=center, text width=1.95cm, font=\tiny] at (9.425,-3.325) {Igname + sauce feuilles};
\node[align=center, text width=1.95cm, font=\tiny] at (11.525,-3.325) {Riz + poisson braisé};
\node[align=center, text width=1.95cm, font=\tiny] at (13.625,-3.325) {Couscous + sauce claire};
\node[align=center, text width=1.95cm, font=\tiny, text=popgreen!60!black] at (1.025,-4.425) {Marche 40 min + foot 20 min};
\node[align=center, text width=1.95cm, font=\tiny, text=popgreen!60!black] at (3.125,-4.425) {Marche 40 min + EPS 60 min};
\node[align=center, text width=1.95cm, font=\tiny, text=popgreen!60!black] at (5.225,-4.425) {Marche 40 min + corde 20 min};
\node[align=center, text width=1.95cm, font=\tiny, text=popgreen!60!black] at (7.325,-4.425) {Marche 40 min + foot 30 min};
\node[align=center, text width=1.95cm, font=\tiny, text=popgreen!60!black] at (9.425,-4.425) {Marche 40 min + danse 20 min};
\node[align=center, text width=1.95cm, font=\tiny, text=popgreen!60!black] at (11.525,-4.425) {Match de foot 90 min};
\node[align=center, text width=1.95cm, font=\tiny, text=popgreen!60!black] at (13.625,-4.425) {Repos actif : promenade 30 min};
\end{tikzpicture}
\end{center}
\legende{Modèle de programme hebdomadaire : menus équilibrés à base d'aliments locaux (lignes repas) et séances d'activité physique (ligne du bas). La moyenne hebdomadaire dépasse 60 min/jour : $(60+100+60+70+60+90+30)/7 = 470/7 \approx 67$ min/jour.}
\end{popfigure}

\courssub{Le suivi : mesurer, comparer, ajuster}

Un programme sans suivi est un programme abandonné. Le suivi repose sur trois outils simples :

\begin{itemize}
\item \textbf{La pesée mensuelle} : une fois par mois, le même jour, à la même heure, sur la même balance (au centre de santé ou à la pharmacie). On calcule le nouvel IMC et on le compare à l'objectif. Une pesée plus fréquente n'a pas de sens : les variations quotidiennes du poids (eau, repas) masquent la tendance réelle.
\item \textbf{Le journal alimentaire simplifié} : un carnet où l'on note chaque jour ce qui a été mangé et bu, et les minutes d'activité. Pas besoin de calculs compliqués : la simple écriture rend les excès visibles.
\item \textbf{Le réajustement} : si le poids ne bouge pas après deux mois, on revient au bilan (étape 1) : les portions sont-elles trop grandes ? Les séances sont-elles réellement tenues ? On ajuste alors les portions ou la durée des séances, jamais en supprimant un repas entier.
\end{itemize}

\begin{savaistu}[Le pouvoir du carnet de bord]
Des études en éducation nutritionnelle ont montré que tenir un simple carnet de bord alimentaire pendant \textbf{une seule semaine} suffit souvent à repérer les principaux excès : boissons sucrées quotidiennes, grignotage de biscuits entre les cours, portions de friture sous-estimées. Beaucoup d'adolescents découvrent ainsi qu'ils boivent l'équivalent de 10 à 15 morceaux de sucre par jour en sodas et jus industriels (un morceau de sucre pèse environ \SI{5}{g} ; une bouteille de soda de \SI{50}{cL} en contient à elle seule l'équivalent de 10 à 12 morceaux). Le carnet agit comme un miroir : on ne corrige bien que ce que l'on voit clairement.
\end{savaistu}

\courssub{Le rôle de la famille et de l'école}

L'adolescent ne décide pas seul de ce qui est cuisiné à la maison ni de son emploi du temps scolaire. Le programme de prévention est donc une affaire collective :

\begin{itemize}
\item \textbf{La famille} : elle fait les achats au marché et prépare les repas. Elle peut privilégier les méthodes de cuisson saines (braisé, vapeur, bouilli plutôt que frit), acheter des fruits de saison (mangues, goyaves, oranges) au lieu de biscuits, et réserver l'eau comme boisson principale du repas. Elle peut aussi marcher avec l'adolescent et limiter le temps d'écran familial.
\item \textbf{L'école} : la cantine scolaire peut proposer des repas équilibrés (féculents + légumes + source de protéines) plutôt que des fritures et boissons sucrées ; les cours d'EPS garantissent au moins deux heures d'activité encadrée par semaine ; les clubs de santé et les journées de sensibilisation permettent l'éducation \textbf{par les pairs} : un élève formé convainc souvent mieux ses camarades qu'un adulte.
\item \textbf{Les pairs} : pratiquer une activité en groupe (football, danse, course) augmente considérablement l'assiduité, car l'engagement devient social.
\end{itemize}

\begin{exoresolu}[Construire et évaluer un programme de prévention]
\textbf{Énoncé.} Aïcha, 15 ans, mesure \SI{1,60}{m} et pèse \SI{70}{kg}. Son journal alimentaire d'une semaine révèle : deux bouteilles de soda de \SI{50}{cL} par jour, fritures au goûter quatre jours sur sept, aucune activité physique en dehors de la marche de 10 minutes jusqu'à l'école. (a) Calcule son IMC et interprète-le. (b) Formule un objectif SMART. (c) Propose deux ajustements alimentaires et un plan d'activité atteignant 60 min/jour en moyenne. (d) Indique comment elle suivra son programme.
\tcblower
\textbf{Solution.}
(a) $IMC = \dfrac{70}{1{,}60^2} = \dfrac{70}{2{,}56} \approx 27{,}3$. Pour une adolescente de 15 ans, cette valeur, rapportée aux courbes de corpulence de l'OMS pour l'âge et le sexe, correspond à un \textbf{surpoids} (à confirmer par un personnel de santé). Il y a un risque accru, à terme, de diabète de type 2 et de maladies cardiovasculaires.

(b) Objectif SMART : « D'ici 4 mois, réduire mon IMC sous 26 en remplaçant les sodas par de l'eau et en pratiquant 60 minutes d'activité par jour au moins 5 jours sur 7 » — spécifique (deux actions), mesurable (IMC, minutes, jours), atteignable (pas de matériel requis), réaliste (pas de perte brutale, croissance en cours), daté (4 mois).

(c) Ajustements alimentaires : remplacer les deux sodas quotidiens par de l'eau (suppression d'environ \SI{1700}{kJ}/jour, soit environ 400 kcal) ; remplacer les fritures du goûter par un fruit ou une petite poignée d'arachides, en ne gardant la friture qu'une fois par semaine. Plan d'activité : marche rapide domicile-école aller-retour (20 min), 20 min de corde à sauter ou de danse le soir, cours d'EPS (60 min) deux fois par semaine, match de ballon de 60 min le samedi. Moyenne : bien supérieure à 60 min/jour, avec activité au moins 5 jours sur 7.

(d) Suivi : pesée mensuelle à la pharmacie avec calcul de l'IMC, carnet de bord quotidien (repas + minutes d'activité), bilan complet au bout de deux mois et réajustement des portions ou des séances si l'IMC ne diminue pas. Aïcha associe sa s\oe ur aînée (achats et cuisine) et sa meilleure amie (corde à sauter en binôme) pour tenir dans la durée.
\end{exoresolu}

\begin{aretenir}[L'essentiel de la section]
\begin{itemize}
\item Un programme de prévention suit quatre étapes : \textbf{bilan chiffré} (IMC, ration, activité) $\rightarrow$ \textbf{objectifs SMART} $\rightarrow$ \textbf{plan hebdomadaire} (menus + séances) $\rightarrow$ \textbf{suivi et ajustement}.
\item Les menus se construisent avec des aliments locaux en respectant la répartition 50--55 \% glucides, 30--35 \% lipides, 12--15 \% protéines, avec fruits, légumes et eau chaque jour.
\item L'activité physique doit atteindre \textbf{60 min/jour en moyenne}, au moins 5 jours sur 7 : la régularité prime sur l'intensité.
\item Les régimes drastiques sont dangereux chez l'adolescent : on ajuste progressivement, jamais on ne prive.
\item Le suivi (pesée mensuelle, carnet de bord) et le soutien de la famille, de l'école et des pairs conditionnent la réussite durable.
\end{itemize}
\end{aretenir}

\begin{exercice}[À toi de jouer]
Relève ta taille et ta masse, calcule ton IMC, puis tiens un carnet de bord alimentaire et d'activité physique pendant trois jours. (1) Dresse ton bilan chiffré. (2) Formule un objectif SMART adapté à ta situation. (3) Construis un programme hebdomadaire (menus à base d'aliments de ta région + séances) atteignant 60 min d'activité par jour en moyenne. (4) Précise tes outils de suivi et les personnes de ton entourage que tu associeras au programme.
\end{exercice}

\newpage

\coursec{Activité d'intégration}

\coursec{Lutte contre la mauvaise alimentation et l'inactivité physique}

\courssub{Objectifs de l'activité}

À l'issue de cette activité d'intégration, tu seras capable de mobiliser l'ensemble des savoirs et savoir-faire de la leçon pour résoudre un problème concret de santé publique. Plus précisément, tu devras :

\begin{itemize}
\item calculer et interpréter l'\cle{IMC} d'un adolescent selon les seuils de référence OMS ;
\item estimer l'apport énergétique d'une ration alimentaire à l'aide d'une table de composition simplifiée et le comparer aux besoins recommandés ;
\item repérer des erreurs alimentaires et des comportements sédentaires dans un mode de vie décrit ;
\item élaborer un programme de prévention de 7 jours combinant alimentation équilibrée (aliments locaux) et activité physique régulière (60 min/jour, recommandation OMS) ;
\item évaluer de manière critique le programme d'un autre groupe à l'aide d'une grille (réalisme, équilibre, régularité).
\end{itemize}

\courssub{Situation déclenchante}

Au lycée de Bafoussam, l'infirmière de l'infirmerie scolaire s'inquiète : lors du dernier bilan de santé, plusieurs élèves de Première présentaient une surcharge pondérale, une fatigue chronique et des difficultés de concentration en classe. Elle sollicite ta classe pour mener une enquête et proposer des solutions. Ton groupe reçoit la fiche de \textbf{Junior}, élève fictif de Première C, 16 ans.

\textbf{Fiche de Junior}

\begin{center}
\begin{tabularx}{\linewidth}{|l|X|}
\hline
\rowcolor{popblueL} \textbf{Caractéristique} & \textbf{Données} \\
\hline
Âge / Classe & 16 ans, Première C \\
\hline
Taille & \SI{1,65}{m} \\
\hline
Masse & \SI{78}{kg} \\
\hline
Activité physique & Aucune ; trajet domicile--lycée en moto-taxi \\
\hline
Temps d'écrans & 5 h par jour (téléphone, télévision, jeux vidéo) \\
\hline
Besoins énergétiques recommandés (adolescent de 16 ans, garçon) & \SI{2500}{kcal/jour} \\
\hline
\end{tabularx}
\end{center}

\textbf{Menus types de Junior (journée 1)}

\begin{center}
\begin{tabularx}{\linewidth}{|l|X|}
\hline
\rowcolor{popblueL} \textbf{Repas} & \textbf{Contenu} \\
\hline
Petit-déjeuner & 2 beignets haricots + 1 boule de pâte d'arachide + 1 verre de jus de fruit industriel sucré \\
\hline
Déjeuner & Grande portion de riz blanc (\SI{300}{g} cuit) + sauce arachide très grasse + 2 morceaux de viande de bœuf frits + 1 bouteille de boisson gazeuse (\SI{50}{cL}) \\
\hline
Collation (16 h) & 1 paquet de chips + 1 boisson gazeuse (\SI{33}{cL}) \\
\hline
Dîner & Fufu maïs (\SI{350}{g}) + eru très huilé + 1 morceau de poisson fumé \\
\hline
\end{tabularx}
\end{center}

\textbf{Table de composition simplifiée (valeurs énergétiques moyennes)}

\begin{center}
\begin{tabularx}{\linewidth}{|X|r|}
\hline
\rowcolor{popblueL} \textbf{Aliment / portion} & \textbf{Énergie (kcal)} \\
\hline
Beignet haricots (1 unité) & 150 \\
\hline
Boule de pâte d'arachide (1 unité) & 180 \\
\hline
Jus industriel sucré (1 verre, \SI{25}{cL}) & 120 \\
\hline
Riz blanc cuit (\SI{100}{g}) & 130 \\
\hline
Sauce arachide grasse (1 louche) & 250 \\
\hline
Viande de bœuf frite (1 morceau) & 200 \\
\hline
Boisson gazeuse (\SI{50}{cL}) & 210 \\
\hline
Boisson gazeuse (\SI{33}{cL}) & 140 \\
\hline
Paquet de chips (1 paquet) & 300 \\
\hline
Fufu maïs (\SI{100}{g}) & 110 \\
\hline
Eru très huilé (1 portion) & 320 \\
\hline
Poisson fumé (1 morceau) & 90 \\
\hline
\end{tabularx}
\end{center}

\textbf{Emploi du temps type de Junior} : 7 h--13 h : cours ; 14 h--16 h : devoirs assis ; 16 h--21 h : écrans (téléphone, télévision) ; aucun sport, aucune marche, aucune tâche physique.

\textbf{Problème posé} : Junior est-il en bonne santé ? Son alimentation et son mode de vie respectent-ils les recommandations ? Quel programme de 7 jours peux-tu lui proposer pour prévenir l'obésité, le diabète de type 2 et les maladies cardiovasculaires ?

\courssub{Consignes}

\begin{enumerate}
\item Calcule l'IMC de Junior. À quelle catégorie des seuils OMS appartient-il ? Quels risques sanitaires cela annonce-t-il ?
\item À l'aide de la table de composition, estime l'apport énergétique total de la journée 1. Compare-le aux besoins recommandés (\SI{2500}{kcal/jour}) et calcule le pourcentage d'excès.
\item Identifie, en les justifiant, \textbf{trois erreurs alimentaires} et \textbf{deux comportements sédentaires} de Junior.
\item Propose pour Junior un objectif raisonnable de rééquilibrage (apport énergétique, activité physique) en t'appuyant sur les recommandations OMS.
\item Élabore sur le tableau-modèle fourni par l'enseignant un \textbf{programme de 7 jours} : menus équilibrés à base d'aliments locaux (ndolé, okok, haricots, poisson, fruits de saison, etc.) et 60 min d'activité physique par jour.
\item Présente ton programme à la classe, puis évalue le programme d'un autre groupe à l'aide de la grille : réalisme, équilibre nutritionnel, régularité de l'activité physique.
\end{enumerate}

\courssub{Ressources à mobiliser}

\begin{aretenir}[Formules et repères utiles]
\begin{itemize}
\item \textbf{Indice de masse corporelle (IMC)} : $\mathrm{IMC} = \dfrac{m}{T^2}$ avec $m$ la masse en kg et $T$ la taille en m. Seuils OMS (adulte, repères pour l'adolescent de plus de 16 ans) : maigreur $< 18,5$ ; corpulence normale : $18,5 \leq \mathrm{IMC} < 25$ ; surpoids : $25 \leq \mathrm{IMC} < 30$ ; obésité : $\mathrm{IMC} \geq 30$.
\item \textbf{Pourcentage d'excès (ou de déficit) énergétique} : $\dfrac{\text{apport} - \text{besoins}}{\text{besoins}} \times 100$.
\item \textbf{Besoins de l'adolescent de 16 ans (garçon)} : environ \SI{2500}{kcal/jour}, avec une répartition équilibrée : glucides (50--55 \%), lipides (30--35 \%), protéines (12--15 \%), plus vitamines, sels minéraux, fibres et eau (\SI{1,5}{L} à \SI{2}{L}/jour).
\item \textbf{Recommandation OMS (5--17 ans)} : au moins \textbf{60 minutes d'activité physique modérée à soutenue par jour} (majoritairement aérobie) ; limiter le temps d'écran récréatif à 2 h/jour.
\item \textbf{Risques sanitaires} de la surcharge pondérale et de la sédentarité : obésité, diabète de type 2 (résistance à l'insuline), hypertension artérielle, dyslipidémie, maladies cardiovasculaires précoces.
\end{itemize}
\end{aretenir}

\begin{methode}[Démarche pour évaluer une ration alimentaire]
\begin{enumerate}
\item Lister tous les aliments et boissons consommés dans la journée avec leurs quantités (poids net ou volume).
\item Convertir les quantités en portions unitaires compatibles avec la table de composition (ex. : \SI{300}{g} de riz = 3 portions de \SI{100}{g}).
\item Multiplier le nombre de portions par la valeur énergétique unitaire pour chaque aliment.
\item Additionner les apports de tous les aliments pour obtenir l'apport énergétique total journalier (AET).
\item Comparer l'AET aux besoins énergétiques recommandés (BER) pour l'âge, le sexe et le niveau d'activité physique (NAP).
\item Calculer l'écart : $\text{Écart} = \text{AET} - \text{BER}$ et le pourcentage : $\frac{\text{Écart}}{\text{BER}} \times 100$.
\item Analyser la répartition qualitative (présence de fruits/légumes, sources de protéines, qualité des lipides, hydratation).
\end{enumerate}
\end{methode}

\courssub{Démarche et corrigé commenté}

\begin{exoresolu}[Bilan santé de Junior et programme de prévention sur 7 jours]
Énoncé : à partir de la fiche de Junior (16 ans, \SI{1,65}{m}, \SI{78}{kg}, menus et emploi du temps décrits), réaliser le bilan santé complet et construire un programme de prévention de 7 jours.

\tcblower

\textbf{Étape 1 : Calcul et interprétation de l'IMC}

\[
\mathrm{IMC} = \frac{m}{T^2} = \frac{78}{1{,}65^2} = \frac{78}{2{,}7225} \approx 28{,}6
\]

L'IMC de Junior est d'environ \textbf{28,6} : il se situe dans l'intervalle $25 \leq \mathrm{IMC} < 30$, ce qui correspond à la catégorie \textbf{surpoids} (pré-obésité), proche du seuil de l'obésité (IMC $\geq$ 30). Sans changement de mode de vie, Junior s'expose à un risque accru d'obésité installée, de diabète de type 2 (résistance à l'insuline favorisée par l'excès de tissu adipeux viscéral), d'hypertension artérielle et de maladies cardiovasculaires (athérosclérose, infarctus, AVC) à l'âge adulte.

\textbf{Étape 2 : Estimation de l'apport énergétique de la journée 1}

\begin{center}
\begin{tabularx}{\linewidth}{|X|r|}
\hline
\rowcolor{popblueL} \textbf{Prise alimentaire} & \textbf{Énergie (kcal)} \\
\hline
Petit-déjeuner : $2 \times 150 + 180 + 120$ & 600 \\
\hline
Déjeuner : Riz \SI{300}{g} ($3 \times 130 = 390$) ; Sauce arachide 250 ; Viande $2 \times 200 = 400$ ; Gazeuse \SI{50}{cL} 210 & 1\,250 \\
\hline
Collation : Chips 300 + Gazeuse \SI{33}{cL} 140 & 440 \\
\hline
Dîner : Fufu \SI{350}{g} ($3{,}5 \times 110 = 385$) ; Eru 320 ; Poisson 90 & 795 \\
\hline
\textbf{Total journalier} & \textbf{3\,085} \\
\hline
\end{tabularx}
\end{center}

Apport énergétique total (AET) : environ \SI{3085}{kcal} pour des besoins énergétiques recommandés (BER) de \SI{2500}{kcal}.

\[
\text{Excès} = \frac{3085 - 2500}{2500} \times 100 = \frac{585}{2500} \times 100 = 23{,}4\ \%
\]

Junior consomme environ \textbf{23,4 \% d'énergie en trop} chaque jour. L'énergie non dépensée est stockée sous forme de triglycérides dans les adipocytes du tissu adipeux (lipogenèse) : c'est le mécanisme physiologique de la prise de masse grasse.

\textbf{Étape 3 : Erreurs alimentaires et comportements sédentaires}

Trois erreurs alimentaires majeures (parmi d'autres) :
\begin{enumerate}
\item \textbf{Excès de sucres rapides et de calories « vides »} : deux boissons gazeuses et un jus industriel par jour, soit environ \SI{470}{kcal} de sucres ajoutés sans fibres, sans vitamines ni satiété durable, favorisant les pics d'insulinémie et la lipogenèse hépatique.
\item \textbf{Excès de lipides d'origine culinaire (graisses de cuisson)} : fritures (beignets, viande frite) et sauces très huilées (eru, sauce arachide) qui alourdissent fortement la ration lipidique bien au-delà des 30--35 \% recommandés, avec des acides gras saturés et potentiellement des acides gras trans (huiles de friture réchauffées).
\item \textbf{Déséquilibre des portions et absence de densité nutritionnelle} : portions de féculents raffinés excessives (\SI{300}{g} riz + \SI{350}{g} fufu = \SI{650}{g} féculents cuits), absence totale de crudités, de légumes verts (hors sauce) et de fruit frais entier dans la journée $\rightarrow$ déficit en fibres, vitamines (C, B9), potassium, magnésium et antioxydants.
\end{enumerate}

Deux comportements sédentaires caractérisés :
\begin{enumerate}
\item \textbf{Temps d'écran récréatif excessif} : 5 h/jour, soit 2,5 fois la recommandation maximale de 2 h/jour (OMS), favorisant le grignotage, la réduction du métabolisme de base et les troubles du sommeil.
\item \textbf{Inactivité physique totale} : NAP (Niveau d'Activité Physique) = 1,2 (sédentaire strict). Trajet motorisé, pas de sport scolaire ou extrascolaire, pas de marche active, tâches ménagères légères uniquement. Dépense énergétique d'activité (DEA) quasi nulle.
\end{enumerate}

\textbf{Étape 4 : Objectifs de rééquilibrage (SMART : Spécifiques, Mesurables, Atteignables, Réalistes, Temporels)}

\begin{itemize}
\item \textbf{Énergétique} : Ramener l'apport vers \SI{2500}{kcal/jour} ($\pm$ 10 \%) en supprimant les boissons sucrées (remplacées par de l'eau), en réduisant les fritures (cuisson vapeur, braisage, papillote), en ajustant les portions de féculents (\SI{150}{g} à \SI{200}{g} cuits par repas principal) et en augmentant la part de légumes et fruits frais.
\item \textbf{Activité physique} : Atteindre \textbf{60 min d'activité physique modérée à soutenue par jour} (marche rapide, football, course lente, corde à sauter, danse, tâches actives : jardinage, nettoyage vigoureux) et limiter les écrans récréatifs à $\leq$ 2 h/jour.
\item \textbf{Pondéral} : Viser une stabilisation de la masse puis une perte progressive (\SI{0,5}{kg} à \SI{1}{kg} par mois maximum), sans régime restrictif dangereux pour la croissance et la puberté (apports en calcium, fer, protéines maintenus).
\end{itemize}

\textbf{Étape 5 : Programme de prévention sur 7 jours (exemple réaliste avec ressources locales)}

\begin{center}
\begin{tabularx}{\linewidth}{|l|X|X|}
\hline
\rowcolor{popblueL} \textbf{Jour} & \textbf{Repas principaux (Déjeuner / Dîner)} & \textbf{Activité physique (60 min cumulées)} \\
\hline
Lundi & \textbf{Déj :} Ndolé (feuilles amères, crevettes séchées, peu d'huile) + Miondo (bâton de manioc, \SI{150}{g}) + Mangue. \newline \textbf{Dîn :} Haricots rouges + Riz (\SI{150}{g}) + Salade tomate/oignon. & 30 min marche rapide (aller-retour lycée) + 30 min football \\
\hline
Mardi & \textbf{Déj :} Haricots blancs + Maïs bouilli (\SI{200}{g}) + Avocat. \newline \textbf{Dîn :} Poisson braisé (bar/tilapia) + Bâton de manioc (\SI{150}{g}) + Banane plantain bouillie (\SI{100}{g}). & 60 min course lente fractionnée (10 min échauffement, 40 min 1' course/1' marche, 10 min retour au calme) \\
\hline
Mercredi & \textbf{Déj :} Okok (feuilles Gnetum, arachide modérée, poisson fumé) + Manioc vapeur (\SI{200}{g}) + Pastèque. \newline \textbf{Dîn :} Omelette (2 œufs, tomates, oignons) + Pain complet (\SI{100}{g}) + Orange. & 30 min corde à sauter (séries de 2 min) + 30 min jeu actif (cache-cache, touche-touche) \\
\hline
Jeudi & \textbf{Déj :} Riz (\SI{150}{g}) + Sauce feuille (feuilles de courge/amarante, peu d'huile) + Poulet grillé + Papaye. \newline \textbf{Dîn :} Soupe de poisson (kwem) + Fufu maïs (\SI{150}{g}) + Ananas. & 60 min marche active domicile--lycée aller-retour (si distance $\leq$ 3 km) ou 45 min marche + 15 min étirements \\
\hline
Vendredi & \textbf{Déj :} Koki (haricots macabo, huile rouge modérée) + Patate douce bouillie (\SI{200}{g}) + Fruit de la passion. \newline \textbf{Dîn :} Salade composée (laitue, tomate, concombre, maïs, thon) + Pain + Yaourt nature. & 45 min football (match) + 15 min étirements / gainage \\
\hline
Samedi & \textbf{Déj :} Sanga (maïs + manioc + feuilles) + Légumes sautés (carottes, haricots verts) + Mangue. \newline \textbf{Dîn :} Brochettes de bœuf maigre + Riz (\SI{150}{g}) + Salade verte + Watermelon. & 60 min tâches actives (jardinage, lessive à la main, nettoyage cour) + jeu de balle \\
\hline
Dimanche & \textbf{Déj :} Poisson braisé + Fufu maïs (\SI{150}{g}) + Légumes vapeur (chou, carottes) + Papaye. \newline \textbf{Dîn :} Bouillie de mil/maïs enrichie (arachide, lait, œuf) + Banane + Pain. & 60 min promenade familiale active (parc, berges, montagne) ou randonnée légère \\
\hline
\end{tabularx}
\end{center}

\textbf{Repères quotidiens communs (non détaillés dans le tableau) :}
\begin{itemize}
\item \textbf{Petit-déjeuner} (ex. : bouillie de mil/soja enrichie + fruit + pain/galette + omelette/fromage) -- \textbf{Collation 10 h / 16 h} (fruit de saison + arachides/amandes \SI{20}{g} ou yaourt) -- \textbf{Hydratation} : Eau comme boisson principale (\SI{1,5}{L} à \SI{2}{L}/jour), pas de boisson sucrée.
\item \textbf{Écrans} : Maximum 2 h loisir/jour, coupure 1 h avant le coucher.
\item \textbf{Sommeil} : 8 à 9 h/nuit (coucher 22 h, lever 6 h).
\end{itemize}

\textbf{Étape 6 : Évaluation croisée par les pairs (Grille d'analyse)}

\begin{center}
\begin{tabularx}{\linewidth}{|l|X|X|}
\hline
\rowcolor{popblueL} \textbf{Critère} & \textbf{Indicateurs de réussite} & \textbf{Appréciation (A / PA / NA)} \\
\hline
\textbf{Réalisme} & Aliments disponibles au marché/localité, budget familial compatible, temps de préparation acceptable, activités faisables sans équipement coûteux. & \\
\hline
\textbf{Équilibre nutritionnel} & Présence quotidienne : protéines (animales/végétales), féculents (portion adaptée), légumes (crus + cuits), fruit frais, source de calcium, eau. Graisses de cuisson maîtrisées. & \\
\hline
\textbf{Régularité de l'AP} & 60 min/jour \textbf{chaque jour} (pas seulement week-end), variété (endurance, renforcement, souplesse), progression possible. & \\
\hline
\textbf{Adaptation à Junior} & Prise en compte des goûts, horaires (cours, devoirs), trajet, environnement (espace pour courir, sécurité). & \\
\hline
\end{tabularx}
\end{center}
\emph{Légende : A = Atteint ; PA = Partiellement Atteint ; NA = Non Atteint.}
\end{exoresolu}

\begin{attention}[Pièges fréquents au Probatoire et en vie courante]
\begin{itemize}
\item \textbf{IMC} : Ne pas oublier de \textbf{mettre la taille au carré} ($\mathrm{IMC} = m/T^2$, pas $m/T$). Utiliser la taille en \textbf{mètres} (pas en cm). L'IMC est un \emph{indicateur de dépistage}, pas un diagnostic clinique : seul un personnel de santé confirme l'état nutritionnel (tour de taille, plis cutanés, bilan sanguin).
\item \textbf{Calcul de la ration} : Convertir \textbf{systématiquement} les quantités réelles en portions de la table (\SI{300}{g} riz = 3 $\times$ \SI{100}{g}). Ne pas confondre poids cru / poids cuit (la table donne souvent le poids cuit pour les féculents).
\item \textbf{Pourcentage d'excès} : La formule est $\frac{\text{AET} - \text{BER}}{\text{BER}} \times 100$. Un excès de 23 \% ne signifie pas que l'apport est 123 \% des besoins (c'est la même chose), mais attention à ne pas écrire $\frac{\text{BER}}{\text{AET}}$.
\item \textbf{Programme « trop parfait »} : Une ration à \SI{1500}{kcal} ou 3 h de sport/jour est \textbf{irréaliste et dangereux} chez un adolescent en croissance (risque de troubles du comportement alimentaire, fatigue, blessures, arrêt de croissance). La grille d'évaluation pénalise le manque de réalisme.
\item \textbf{Boissons} : « Jus de fruit » $\neq$ « Fruit pressé maison sans sucre ». Les jus industriels et boissons gazeuses sont des \textbf{calories liquides} sans satiété : première cible de la réduction.
\item \textbf{Vocabulaire} : « Surcharge pondérale » = terme clinique général (IMC > 25). « Surpoids » = catégorie OMS 25--29,9. « Obésité » = IMC $\geq$ 30. Ne pas les confondre dans la rédaction.
\end{itemize}
\end{attention}

\begin{savaistu}[Contexte camerounais et enjeux de santé publique]
\begin{itemize}
\item \textbf{Transition nutritionnelle} : Au Cameroun, la coexistence de la sous-nutrition (retard de croissance, carences en fer/vitamine A) et de la sur-nutrition (surpoids, obésité, diabète, HTA) dans les mêmes ménages est un défi majeur. L'urbanisation rapide (Yaoundé, Douala, Bafoussam) favorise l'accès aux aliments ultra-transformés (boissons sucrées, pâtes instantanées, huiles bon marché) et la sédentarité (moto-taxi, écrans, insécurité pour le sport en plein air).
\item \textbf{Stratégie nationale} : Le Ministère de la Santé Publique (MINSANTE) et le Ministère des Sports (MINSEP) promeuvent le « Sport pour tous » et la lutte contre les Maladies Non Transmissibles (MNT). L'infirmerie scolaire est un maillon clé du dépistage précoce (IMC, tension artérielle, glycémie).
\item \textbf{Aliments locaux protecteurs} : Les feuilles (ndolé, okok, eru, folong, kpwem), les légumineuses (haricots, niébé, soja), les tubercules (manioc, macabo, patate, igname), les fruits de saison (mangue, papaye, ananas, pastèque, safou, avocat) et les poissons (braisés, fumés, séchés) sont des piliers d'une alimentation équilibrée \emph{si} les modes de cuisson limitent l'huile et le sel.
\item \textbf{Journée mondiale du diabète (14 novembre)} et \textbf{Journée mondiale de l'obésité (4 mars)} : Occasions de campagnes de dépistage et d'éducation dans les lycées.
\end{itemize}
\end{savaistu}

\courssub{Bilan de l'activité}

Cette activité t'a permis de mettre en évidence plusieurs points essentiels pour ton examen et ta vie de citoyen. D'abord, un \textbf{bilan santé repose sur des données chiffrées objectives} : l'IMC et l'estimation de la ration alimentaire transforment une impression subjective (« Junior mange beaucoup ») en un constat scientifique argumenté (IMC de 28,6 ; excès énergétique d'environ 23,4 \%). Ensuite, tu as vérifié par le calcul que la \textbf{surcharge pondérale résulte d'un bilan énergétique déséquilibré} : apports (3085 kcal) supérieurs aux dépenses (besoins 2500 kcal + DEA quasi nulle), aggravés par la sédentarité (5 h d'écrans, aucune activité physique volontaire). Enfin, tu as constaté qu'il est possible de construire un \textbf{programme de prévention réaliste, culturellement ancré et économiquement accessible} : les aliments camerounais traditionnels permettent une alimentation équilibrée à condition de maîtriser les portions, les graisses de cuisson et les boissons sucrées, tandis que 60 minutes d'activité quotidienne (marche, jeux de ballon, corde à sauter, tâches ménagères actives) ne demandent aucun équipement coûteux.

Tu es désormais capable de :
\begin{enumerate}
\item \textbf{Évaluer une ration alimentaire} par rapport aux besoins recommandés (Savoir-faire 1) ;
\item \textbf{Interpréter l'IMC} et le relier aux risques sanitaires (Savoir-faire 2) ;
\item \textbf{Élaborer un programme simple} combinant alimentation équilibrée et activité physique (Savoir-faire 3).
\end{enumerate}
Trois compétences directement mobilisables pour toi-même, ta famille, ta communauté, et exigibles au Probatoire.

\newpage

\coursec{Méthodes et exercices résolus}

\coursec{Lutte contre la mauvaise alimentation et l'inactivité physique}

\courssub{Équilibre alimentaire et besoins nutritionnels de l'adolescent}

\begin{definition}[Équilibre alimentaire]
L'\cle{équilibre alimentaire} est un état où la \cle{ration alimentaire} (ensemble des aliments consommés en 24~h) couvre exactement les \cle{besoins nutritionnels} de l'organisme, sans excès ni déficit, en quantité (énergie) et en qualité (macronutriments, micronutriments, eau, fibres). Il se construit sur plusieurs jours, pas nécessairement sur un seul repas.
\end{definition}

\begin{definition}[Besoins énergétiques et ANC]
Les \cle{Apports Nutritionnels Conseillés (ANC)} pour un adolescent de 15 à 18 ans (fin de croissance, activité intellectuelle intense) sont des valeurs de référence de population :
\begin{itemize}
\item \textbf{Énergie totale} : \SI{\approx 2500}{kcal/j} (garçon) ; \SI{\approx 2200}{kcal/j} (fille). À majorer de \SI{10}{à 20}{\%} pour une activité physique soutenue (agriculture, sport de compétition).
\item \textbf{Protéines} : \SI{12}{à 15}{\%} de l'apport énergétique total (AET) $\approx$ \SI{0,8}{à 1,0}{g/kg/j}. Rôle : construction, réparation, enzymes, anticorps.
\item \textbf{Lipides} : \SI{30}{à 35}{\%} de l'AET. Privilégier acides gras insaturés (huiles végétales, poisson) ; limiter saturés (gras animaux, fritures) et \emph{trans} (produits industriels).
\item \textbf{Glucides} : \SI{50}{à 55}{\%} de l'AET. Privilégier \cle{glucides complexes} (amidon : riz, maïs, manioc, plantain, igname, pain) ; limiter \cle{sucres rapides} (boissons sucrées, confiseries, pâtisseries) à $< 10\,\%$ de l'AET (OMS).
\item \textbf{Fibres} : \SI{\ge 25}{g/j} (légumes feuilles, légumineuses, fruits, céréales complètes). Rôle : transit, satiété, modulation glycémique.
\item \textbf{Eau} : \SI{1,5}{à 2}{L/j} (boisson + eau des aliments). Seule boisson indispensable.
\end{itemize}
\end{definition}

\begin{propriete}[Équivalences énergétiques des macronutriments]
Pour calculer la valeur énergétique d'une ration, on utilise les coefficients d'Atwater (valeurs moyennes) :
\[
\begin{array}{c|c}
\text{Macronutriment} & \text{Énergie (kcal/g)} \\ \hline
\text{Glucides} & 4 \\
\text{Protéines} & 4 \\
\text{Lipides} & 9 \\
\text{Alcool (éthanol)} & 7 \quad (\text{à exclure chez l'adolescent})
\end{array}
\]
\end{propriete}

\begin{popfigure}
\begin{tikzpicture}[
  box/.style={draw=popdark, thick, rounded corners, fill=popcream, minimum width=3cm, minimum height=1.2cm, align=center},
  arrow/.style={-Stealth, thick, popdark},
  labelbox/.style={draw=popblue, thick, rounded corners, fill=popblueL, minimum width=3.5cm, minimum height=1cm, align=center}
]
\node[box, align=center] (entree) at (0,0){APPORTS\\Aliments \& Boissons\\(Glucides, Lipides, Protéines,\\Alcool)};
\node[labelbox, align=center] (stock) at (5,0){STOCKAGE\\Tissu adipeux (graisses)\\\& Glycogène (foie, muscles)};
\node[box, align=center] (sortie) at (10,0){DÉPENSES\\Métabolisme de base (60-70\%)\\+ Activité physique (20-30\%)\\+ Thermogenèse (10\%)};

\draw[arrow] (entree.east) -- node[above, font=\small] {Énergie ingérée} (stock.west);
\draw[arrow] (stock.east) -- node[above, font=\small] {Énergie mobilisée} (sortie.west);
\draw[arrow, poporange] (entree.north) -- ++(0,0.8) -| node[above, pos=0.25, font=\small, poporange, align=center]{Excès chronique\\ $\rightarrow$ Prise de poids} (stock.north);
\draw[arrow, popgreen] (sortie.north) -- ++(0,0.8) -| node[above, pos=0.75, font=\small, popgreen, align=center]{Déficit chronique\\ $\rightarrow$ Perte de poids} (stock.north);

\node[draw=popdark, dashed, thick, rounded corners, fit=(entree)(sortie), inner sep=5pt, label={[font=\bfseries, popdark]above:BILAN ÉNERGÉTIQUE QUOTIDIEN}] {};
\end{tikzpicture}
\legende{Bilan énergétique : l'équilibre du poids corporel résulte de l'égalité entre apports (alimentation) et dépenses (métabolisme de base + activité physique + thermogenèse). Un excès chronique d'apports conduit au stockage sous forme de graisses (obésité).}
\end{popfigure}

\begin{savaistu}[Contexte camerounais : Enquête STEPS 2013 et transition nutritionnelle]
Au Cameroun, l'enquête STEPS (OMS/Ministère de la Santé) a révélé une \cle{transition nutritionnelle} rapide : coexistence de la dénutrition (retard de croissance, anémie) et de la surcharge pondérale (surpoids, obésité), surtout en milieu urbain (Douala, Yaoundé). En 2013, \SI{\approx 25}{\%} des adultes étaient en surpoids/obèses. Les adolescents sont exposés à la « malbouffe » (beignets-haricots, boissons sucrées, fritures) et à la sédentarité (écrans, transport motorisé). Le Programme National de Lutte contre les Maladies Non Transmissibles (PNLMNT) promeut l'alimentation saine (produits locaux : ndolé, eru, foléré, fruits de saison, poisson, arachides) et l'activité physique.
\end{savaistu}

\courssub{Évaluer une ration alimentaire par rapport aux besoins recommandés}

\begin{methode}[Évaluer une ration alimentaire (démarche en 6 étapes)]
\begin{enumerate}
\item \textbf{Identifier les besoins de référence (ANC)} : selon l'âge, le sexe, l'activité physique. Pour un lycéen de 16-18 ans : Garçon \SI{\approx 2500}{kcal/j}, Fille \SI{\approx 2200}{kcal/j}. Répartition macronutriments : P \SI{12-15}{\%}, L \SI{30-35}{\%}, G \SI{50-55}{\%}.
\item \textbf{Inventorier les aliments consommés} sur 24~h (3 repas + collation si existante) avec \textbf{quantités précises} (poids net ou volume). Utiliser des tables de composition (ex. Table Ciqual, tables africaines FAO).
\item \textbf{Calculer l'apport en macronutriments (g)} pour chaque aliment, puis sommer : $P_{\text{tot}}$, $L_{\text{tot}}$, $G_{\text{tot}}$.
\item \textbf{Calculer l'apport énergétique total (AET)} en kcal :
$E_P = 4 \times P_{\text{tot}}$ ; $E_L = 9 \times L_{\text{tot}}$ ; $E_G = 4 \times G_{\text{tot}}$ ; $\text{AET} = E_P + E_L + E_G$.
\item \textbf{Calculer la répartition \%} : $\%P = \frac{E_P}{\text{AET}}\times 100$, $\%L = \frac{E_L}{\text{AET}}\times 100$, $\%G = \frac{E_G}{\text{AET}}\times 100$.
\item \textbf{Comparer, interpréter et corriger} :
\begin{itemize}
\item AET vs Besoin : Excès ($> 110\,\%$), Équilibré ($90-110\,\%$), Déficit ($< 90\,\%$).
\item Répartition vs ANC : chaque macronutrient dans sa fourchette ?
\item Qualité : présence fruits/légumes (fibres, vitamines), eau, limitation sel/sucres ajoutés.
\item Proposer des substitutions concrètes et accessibles localement.
\end{itemize}
\end{enumerate}
\end{methode}

\begin{exoresolu}[La ration de Marlyse est-elle équilibrée ?]
Marlyse, 16 ans, élève en Première C à Douala, consomme en une journée : 80~g de protéines, 90~g de lipides et 300~g de glucides. Ses besoins sont estimés à \SI{2200}{kcal/jour}.
\begin{enumerate}
\item Calcule l'apport énergétique total de sa ration.
\item Calcule la part (en \%) de chaque macronutriment dans cet apport.
\item Conclus sur l'équilibre de cette ration et propose une correction.
\end{enumerate}
\tcblower
\textbf{1. Apport énergétique total.}

On applique les équivalences énergétiques (coefficients d'Atwater) :
\begin{align*}
E_{\text{protéines}} &= 80 \times 4 = \SI{320}{kcal}\\
E_{\text{lipides}} &= 90 \times 9 = \SI{810}{kcal}\\
E_{\text{glucides}} &= 300 \times 4 = \SI{1200}{kcal}
\end{align*}
\[E_{\text{total}} = 320 + 810 + 1200 = \SI{2330}{kcal}\]

\textbf{2. Répartition en pourcentage.}
\begin{align*}
\%\,\text{Protéines} &= \dfrac{320}{2330}\times 100 \approx 13,7\,\%\\
\%\,\text{Lipides} &= \dfrac{810}{2330}\times 100 \approx 34,8\,\%\\
\%\,\text{Glucides} &= \dfrac{1200}{2330}\times 100 \approx 51,5\,\%
\end{align*}
Vérification : $13,7 + 34,8 + 51,5 = 100\,\%$.

\textbf{3. Conclusion et correction.}

\begin{itemize}
\item \textbf{Quantitatif} : AET = \SI{2330}{kcal} vs Besoin = \SI{2200}{kcal}. Excès de \SI{130}{kcal} ($+5,9\,\%$). Situation \textbf{acceptable} (proche de l'équilibre), mais à surveiller sur le long terme (risque de prise de poids lente).
\item \textbf{Qualitatif (Répartition)} : Protéines 13,7~\% $\in [12-15]$\% \checkmark ; Lipides 34,8~\% $\in [30-35]$\% \checkmark (mais limite haute) ; Glucides 51,5~\% $\in [50-55]$\% \checkmark.
\item \textbf{Analyse critique} : La part des lipides est à la limite supérieure. L'énoncé ne précise pas l'apport en fibres, vitamines, eau.
\end{itemize}
\textbf{Corrections proposées} : Remplacer une partie des graisses de cuisson (huile de palme raffinée, margarine) par des huiles non chauffées (arachide, coton) ou réduire les quantités ; privilégier cuissons vapeur, grillade, papillote ; augmenter les glucides complexes (riz, maïs, tubercule) et \textbf{imposer} fruits/légumes à chaque repas (ndolé, eru, gombo, ananas, papaye) pour fibres et micronutriments ; boire \SI{1,5}{L} d'eau/jour.
\end{exoresolu}

\begin{exercice}[Évaluation de la ration de Kevin]
Kevin, 17 ans, Première TI, garçon actif (marche 30 min/jour + cours), besoins \SI{2600}{kcal/j}. Sa ration journalière :
\begin{itemize}
\item Petit-déj : Pain (100g) + Beurre (20g) + Chocolat au lait (20g) + Lait (200ml).
\item Déjeuner : Riz blanc (250g cuit) + Sauce arachide (huile 30g, pâte arachide 40g, viande 50g) + Plantain frit (100g).
\item Goûter : Beignet (50g) + Boisson sucrée (300ml).
\item Dîner : Pâtes (200g cuites) + Sauce tomate/huile (15g) + Poisson grillé (80g) + Mangue (100g).
\end{itemize}
\textit{Données moyennes (g/100g ou 100ml) : Pain (P8, L3, G55) ; Beurre (L82) ; Chocolat lait (P7, L30, G55) ; Lait entier (P3, L3, G5) ; Riz cuit (P2,5, G25) ; Arachide grasse (P25, L50, G10) ; Viande maigre (P20, L5) ; Plantain frit (L15, G30) ; Beignet (P5, L20, G50) ; Boisson sucrée (G10) ; Pâtes cuites (P5, G25) ; Huile (L100) ; Poisson (P20, L2) ; Mangue (G15).}
\begin{enumerate}
\item Calcule les apports totaux en P, L, G (en g) et l'AET (kcal).
\item Calcule les pourcentages. Compare aux ANC.
\item La ration est-elle équilibrée ? Propose 3 corrections concrètes.
\end{enumerate}
\end{exercice}

\courssub{L'IMC : calcul, interprétation et lien avec les risques sanitaires}

\begin{definition}[Indice de Masse Corporelle (IMC)]
L'\cle{IMC} est un indicateur simple de la corpulence, valide pour les adultes et les adolescents en fin de croissance (après 18 ans, ou 16-17 ans si croissance terminée). Il ne distingue pas masse grasse et masse musculaire.
\[\mathrm{IMC} = \dfrac{m}{t^2}\]
$m$ : masse en \textbf{kilogrammes (kg)} ; $t$ : taille en \textbf{mètres (m)}. Unité : \SI{}{kg/m^2}.
\end{definition}

\begin{popfigure}
\begin{tikzpicture}[
  zone/.style={draw=popdark, thick, minimum height=1.2cm, align=center, font=\small},
  label/.style={font=\bfseries\small, text=popdark}
]
% Zones IMC
\node[zone, fill=popblue!20, minimum width=2.5cm, align=center] (m1) at (0,0){Maigreur\\$< 18,5$};
\node[zone, fill=popgreen!20, minimum width=3.5cm, anchor=west, align=center] (n) at (m1.east){Normale\\$18,5 - 24,9$};
\node[zone, fill=poporange!20, minimum width=3cm, anchor=west, align=center] (s) at (n.east){Surpoids\\$25 - 29,9$};
\node[zone, fill=poppink!20, minimum width=3cm, anchor=west, align=center] (o1) at (s.east){Obésité I\\$30 - 34,9$};
\node[zone, fill=poppink!40, minimum width=3cm, anchor=west, align=center] (o2) at (o1.east){Obésité II\\$35 - 39,9$};
\node[zone, fill=poppink!60, minimum width=2.5cm, anchor=west, align=center] (o3) at (o2.east){Obésité III\\$\ge 40$};

% Flèche échelle
\draw[popdark, thick, -Stealth] (-0.5,-1.2) -- (17.5,-1.2) node[right, font=\small] {IMC (kg/m$^2$)};
\foreach \x/\val in {0/10, 2.5/15, 5.5/20, 8.5/25, 11.5/30, 14.5/35, 17/40}
  \draw (\x,-1.1) -- (\x,-1.3) node[below, font=\tiny] {\val};

% Exemples marqueurs
\draw[popblue, very thick] (1.25,0.6) -- (1.25,0.9) node[above, popblue, font=\tiny] {Brice (16,7)};
\draw[poporange, very thick] (7.5,0.6) -- (7.5,0.9) node[above, poporange, font=\tiny] {Alain (28,4)};
\draw[poppink, very thick] (11.5,0.6) -- (11.5,0.9) node[above, poppink, font=\tiny] {M. Etoundi (32,1)};
\draw[popgreen, very thick] (4.5,0.6) -- (4.5,0.9) node[above, popgreen, font=\tiny] {Cible (22)};
\end{tikzpicture}
\legende{Classification OMS de l'IMC (adulte/adolescent post-pubère) avec positionnement des cas cliniques du cours. Zone verte = risque minimal. Zones orange/rose = risque accru cardiométabolique. Zone bleue = risque carences/infections.}
\end{popfigure}

\begin{methode}[Calculer et interpréter l'IMC (sans erreur)]
\begin{enumerate}
\item \textbf{Convertir la taille en mètres} : \SI{165}{cm} = \SI{1,65}{m}. \textbf{(Erreur n°1 : diviser par 165 au lieu de $1,65^2$)}.
\item \textbf{Calculer $t^2$} (surface corporelle approximative).
\item \textbf{Diviser} $m$ par $t^2$. Arrondir au dixième (\SI{1}{chiffre après la virgule}).
\item \textbf{Écrire l'unité} : \SI{}{kg/m^2}.
\item \textbf{Situer dans la classification OMS} (tableau ci-dessous) et \textbf{nommer la catégorie}.
\item \textbf{Relier aux risques sanitaires} (voir tableau).
\end{enumerate}
\end{methode}

\begin{center}
\begin{tabularx}{\linewidth}{|l|X|X|}\hline
\rowcolor{popblueL} \textbf{IMC (kg/m$^2$)} & \textbf{Catégorie (OMS)} & \textbf{Risques sanitaires principaux} \\ \hline
$< 17,0$ & Maigreur sévère (Dénutrition) & Carences (fer, vit A, Zn, Protéines), immunodépression, ostéoporose, aménorrhée, mortalité accrue \\ \hline
$17,0 - 18,4$ & Maigreur modérée & Fatigue, baisse concentration, fragilité infectieuse, retard cicatrisation \\ \hline
$18,5 - 24,9$ & \textbf{Corpulence normale} & \textbf{Risque minimal} (référence) \\ \hline
$25,0 - 29,9$ & \textbf{Surpoids} & Risque accru : Diabète T2, HTA, dyslipidémie, stéatose hépatique, arthrose \\ \hline
$30,0 - 34,9$ & \textbf{Obésité classe I (modérée)} & Risque élevé : Diabète T2, HTA, coronaropathie, AVC, apnée du sommeil \\ \hline
$35,0 - 39,9$ & \textbf{Obésité classe II (sévère)} & Risque très élevé : complications multiples, insuffisance veineuse, risque chirurgical \\ \hline
$\ge 40,0$ & \textbf{Obésité classe III (massive/morbide)} & Risque vital, espérance de vie réduite, prise en charge spécialisée obligatoire \\ \hline
\end{tabularx}
\end{center}

\begin{exoresolu}[L'IMC de deux élèves à Yaoundé]
Alain : $m = \SI{82}{kg}$, $t = \SI{170}{cm}$ ; Brice : $m = \SI{54}{kg}$, $t = \SI{180}{cm}$.
\begin{enumerate}
\item Calcule l'IMC de chaque élève (arrondi au dixième).
\item Interprète chaque résultat et indique les risques sanitaires associés.
\end{enumerate}
\tcblower
\textbf{1. Calcul des IMC.}

Conversion : Alain $t = \SI{1,70}{m}$ ; Brice $t = \SI{1,80}{m}$.

Alain : $\mathrm{IMC}_A = \dfrac{82}{1,70^2} = \dfrac{82}{2,89} \approx \SI{28,4}{kg/m^2}$.

Brice : $\mathrm{IMC}_B = \dfrac{54}{1,80^2} = \dfrac{54}{3,24} \approx \SI{16,7}{kg/m^2}$.

\textbf{2. Interprétation et risques.}

\begin{itemize}
\item \textbf{Alain (28,4)} : \cle{Surpoids} ($25 \le \mathrm{IMC} < 30$). Risque majoré de \cle{diabète de type 2}, \cle{hypertension artérielle (HTA)}, \cle{dyslipidémie} (cholestérol LDL haut, HDL bas), stéatose hépatique non alcoolique (NASH). Conseil : rééquilibrage alimentaire (moins fritures, sodas, plus légumes) + activité physique \SI{\ge 150}{min/sem} (marche rapide, football).
\item \textbf{Brice (16,7)} : \cle{Maigreur sévère} ($\mathrm{IMC} < 17$). Dénutrition probable. Risques : anémie ferriprive, carences vitaminiques (A, D, B12), ostéopénie, baisse défenses immunitaires (infections répétées), fatigue chronique, troubles de la concentration. Causes à explorer : ration insuffisante, parasitose intestinale (vers), tuberculose, VIH, troubles du comportement alimentaire (anorexie), hyperthyroïdie. Conseil : augmentation progressive apports énergétiques/protéiques (haricots, arachides, poisson, œufs, lait) + consultation médicale pour étiologie.
\end{itemize}
\textbf{Conclusion} : Aucun n'est en zone normale. Prise en charge différenciée urgente.
\end{exoresolu}

\courssub{Conséquences de la sédentarité et de la mauvaise alimentation : mécanismes physiopathologiques}

\begin{experience}[Modélisation de la chaîne causale : du mode de vie à la maladie]
\textbf{Objectif} : Comprendre comment l'association « Excès énergétique + Sédentarité » conduit aux pathologies chroniques.
\textbf{Schéma explicatif (à reproduire)} :
\end{experience}

\begin{popfigure}
\begin{tikzpicture}[
  node/.style={draw=popdark, thick, rounded corners, fill=popcream, minimum width=3.2cm, minimum height=1.1cm, align=center, font=\small},
  patho/.style={draw=poppink, thick, rounded corners, fill=poppinkL, minimum width=3.5cm, minimum height=1.1cm, align=center, font=\small},
  arrow/.style={-Stealth, thick, popdark},
  link/.style={-Stealth, poporange, thick, dashed}
]
% Niveau 1 : Causes
\node[node, align=center] (alim) at (0,3){MAUVAISE\\ALIMENTATION\\Excès Lipides \& Sucres rapides\\Défaut Fibres \& Micronutriments};
\node[node, align=center] (sed) at (6,3){SÉDENTARITÉ\\Dépense physique $<$ \SI{60}{min/j}\\Temps écran $>$ \SI{4}{h/j}\\Transport motorisé exclusif};

% Niveau 2 : Bilan
\node[node, fill=poporangeL, align=center] (bilan) at (3,1.5){BILAN ÉNERGÉTIQUE POSITIF\\Chronique\\Apports $>$ Dépenses};

% Niveau 3 : Conséquence immédiate
\node[node, fill=popgoldL, align=center] (obes) at (3,0){SURCHARGE PONDÉRALE / OBÉSITÉ\\Hypertrophie \& Hyperplasie adipocytes\\Tissu adipeux viscéral (android) \\IMC $\ge 25$};

% Niveau 4 : Mécanismes pathologiques
\node[patho, align=center] (diab) at (-3,-1.5){DIABÈTE TYPE 2\\Résistance à l'insuline (IR)\\Hyperinsulinémie compensatrice\\Épuisement $\beta$ pancréatiques\\Glycémie à jeun $\ge \SI{1,26}{g/L}$};
\node[patho, align=center] (cvd) at (9,-1.5){MALADIES CARDIOVASCULAIRES\\Dyslipidémie (LDL$\uparrow$, HDL$\downarrow$, TG$\uparrow$)\\\cle{Athérome} (plaques cholestérol)\\\cle{HTA} (raideur artérielle)\\Infarctus, AVC, Artérite};

% Flèches
\draw[arrow] (alim.south) -- (bilan.north west);
\draw[arrow] (sed.south) -- (bilan.north east);
\draw[arrow] (bilan.south) -- (obes.north);
\draw[arrow] (obes.south west) -- (diab.north east);
\draw[arrow] (obes.south east) -- (cvd.north west);

% Liens transverses
\draw[link] (diab.east) -- node[above, font=\tiny, poporange] {Facteurs de risque partagés} (cvd.west);

% Annotations mécanismes
\node[font=\tiny, popdark, align=left, anchor=north west] at (-5.5,-2.5) {
\textbf{Mécanisme Diabète T2 :}\\
Adipocytes hypertrophiés $\rightarrow$ sécrètent cytokines pro-inflammatoires (TNF-$\alpha$, IL-6) \& acides gras libres $\rightarrow$ inhibent signalisation insuline (récepteur tyrosine kinase) dans muscle/foie/adipocytes $\rightarrow$ \cle{Résistance à l'insuline} $\rightarrow$ Glucose reste dans sang.
};
\node[font=\tiny, popdark, align=left, anchor=north east] at (11.5,-2.5) {
\textbf{Mécanisme Athérome :}\\
Excès LDL-cholestérol + HTA + Inflammation $\rightarrow$ dépôt cholestérol dans intima artère $\rightarrow$ \cle{plaques d'athérome} $\rightarrow$ rétrécissement lumière (sténose) $\rightarrow$ ischémie (angor, infarctus) ou rupture $\rightarrow$ thrombose (AVC, IDM).
};
\end{tikzpicture}
\legende{Chaîne causale unifiée : Mauvaise alimentation + Sédentarité $\rightarrow$ Bilan énergétique positif $\rightarrow$ Obésité (surtout viscérale) $\rightarrow$ Double chemin pathologique : Résistance à l'insuline (Diabète T2) et Athérosclérose (Maladies Cardiovasculaires). L'obésité viscérale est le lien central (tissu adipeux = organe endocrinien inflammatoire).}
\end{popfigure}

\begin{propriete}[Diabète de type 2 (DT2) vs Diabète de type 1 (DT1) — Distinction cruciale pour le Probatoire]
\begin{center}
\begin{tabular}{|l|c|c|}\hline
\rowcolor{popblueL} \textbf{Caractère} & \textbf{Diabète Type 1 (DT1)} & \textbf{Diabète Type 2 (DT2)} \\ \hline
\textbf{Âge typique} & Enfant / Jeune adulte & Adulte $> 40$ ans (mais $\downarrow$ âge : adolescents obèses) \\ \hline
\textbf{Mécanisme principal} & \textbf{Auto-immunité} : Destruction cellules $\beta$ pancréatiques $\rightarrow$ \textbf{Déficit absolu en insuline} & \textbf{Résistance à l'insuline} (récepteur/post-récepteur) + \textbf{Déficit relatif/secrétoire tardif} \\ \hline
\textbf{Lien avec obésité/sédentarité} & \textbf{Aucun} (facteurs génétiques + environnementaux viraux) & \textbf{Majeur} : Obésité viscérale cause n°1 (80-90\% des cas) \\ \hline
\textbf{Traitement initial} & \textbf{Insuline} obligatoire à vie & Hygiène de vie (régime + sport) $\pm$ Antidiabétiques oraux (Metformine) $\pm$ Insuline tardive \\ \hline
\textbf{Cétose} & Fréquente (acidocétose) & Rare (sauf stress sévère) \\ \hline
\end{tabular}
\end{center}
\textbf{Au Probatoire} : Pour un sujet obèse/sédentaire hyperglycémique, on écrit \textbf{« Résistance à l'insuline »} (et non « manque d'insuline »). Le diagnostic biologique : Glycémie à jeun $\ge \SI{1,26}{g/L}$ (ou $\ge \SI{7,0}{mmol/L}$) à deux reprises, ou HbA1c $\ge 6,5\,\%$.
\end{propriete}

\begin{exoresolu}[Le cas de M. Etoundi — Raisonnement causal complet]
M. Etoundi, 45 ans, employé de bureau à Bafoussam, mesure \SI{1,72}{m} et pèse \SI{95}{kg}. Il déjeune chaque jour de plats très gras accompagnés de boissons sucrées, ne pratique aucune activité physique et se déplace uniquement en voiture. Sa glycémie à jeun est de \SI{1,40}{g/L}.
\begin{enumerate}
\item Calcule son IMC et interprète-le.
\item À l'aide des données, explique le diagnostic probable et le mécanisme de la maladie.
\end{enumerate}
\tcblower
\textbf{1. IMC.}
\[\mathrm{IMC} = \dfrac{95}{1,72^2} = \dfrac{95}{2,9584} \approx \SI{32,1}{kg/m^2}\]
Interprétation : \cle{Obésité modérée (classe I)} ($30 \le \mathrm{IMC} < 35$). Risque cardiométabolique élevé.

\textbf{2. Diagnostic et mécanisme physiopathologique.}

\begin{itemize}
\item \textbf{Diagnostic probable} : \cle{Diabète de type 2}. Justification : Glycémie à jeun \SI{1,40}{g/L} $> \SI{1,26}{g/L}$ (seuil OMS) + contexte obésité + âge + mode de vie.
\item \textbf{Chaîne causale détaillée} :
\begin{enumerate}
\item \textbf{Excès énergétique chronique} : Apports (graisses saturées, sucres rapides) $\gg$ Dépenses (métabolisme de base seulement, PAS d'activité physique).
\item \textbf{Stockage adipocytaire} : Excès $\rightarrow$ hypertrophie adipocytes viscéraux (graisse abdominale). Tour de taille probable $> \SI{102}{cm}$ (homme).
\item \textbf{Adiposopathie (maladie du tissu adipeux)} : Adipocytes hypertrophiés deviennent dysfonctionnels : sécrètent moins d'adiponectine (insulinosensibilisante) et plus de \textbf{cytokines pro-inflammatoires} (TNF-$\alpha$, IL-6, CRP) + libèrent massivement des \textbf{acides gras libres (AGL)} dans la veine porte $\rightarrow$ foie.
\item \textbf{Résistance à l'insuline (RI) hépatique \& périphérique} : AGL hépatiques $\rightarrow$ gluconéogenèse accrue + inhibition signalisation insuline (voie IRS/PI3K/Akt) $\rightarrow$ le glucose ne rentre plus dans cellules (muscle, adipeux) $\rightarrow$ \textbf{Hyperglycémie}.
\item \textbf{Hyperinsulinémie compensatrice} : Pancréas surstimulé $\rightarrow$ sécrète plus d'insuline (hyperinsulinémie à jeun) pour vaincre la RI. À terme, épuisement des $\beta$-cellules $\rightarrow$ déficit insulinique relatif.
\end{enumerate}
\end{itemize}
\textbf{Conclusion} : Le mode de vie (alimentation déséquilibrée + inactivité physique) est la \textbf{cause racine} de l'obésité, elle-même cause directe du diabète de type 2 par le mécanisme de la résistance à l'insuline induite par l'adiposopathie viscérale.
\end{exoresolu}

\courssub{Recommandations officielles : Alimentation et Activité Physique (OMS / PNNS Cameroun)}

\begin{aretenir}[Recommandations clés à connaître par cœur pour le Probatoire]
\begin{itemize}
\item \textbf{Alimentation (Repères PNNS adaptés Cameroun)} :
\begin{enumerate}
\item \textbf{3 repas/jour} + 1 collation si besoin (pas de grignotage).
\item \textbf{5 fruits et légumes/jour} (\SI{\ge 400}{g}) : variés, de saison, crus et cuits (ndolé, eru, gombo, foléré, ananas, papaye, agrumes, carotte, chou).
\item \textbf{Féculents à chaque repas} : Privilégier complets/locaux (riz étuvé, maïs, mil, sorgho, manioc, igname, plantain, patate douce) $\rightarrow$ Glucides complexes + fibres.
\item \textbf{Protéines} : 1 à 2 portions/jour. Alterner : Poisson (riche Oméga-3, abordable : maquereau, sardine, capitaine), Légumineuses (haricots, niébé, soja, arachides), Viandes maigres (poulet, bœuf dégraissé), Œufs. Limiter charcuteries, abats, viandes grasses.
\item \textbf{Matières grasses} : Huiles végétales non chauffées (arachide, coton, tournesol, rouge non raffinée modérément) ; limiter beurre, margarine, huile de friture, graisses cachées (pâtisseries, plats préparés).
\item \textbf{Produits laitiers} : 2 à 3/j (lait, yaourt, fromage) pour Ca, Vit D. Si intolérance : lait fermenté, soja enrichi.
\item \textbf{Eau} : \SI{1,5}{à 2}{L/j} (seule boisson à volonté). Zéro alcool < 18 ans.
\item \textbf{Sel} : $< \SI{5}{g/j}$ (1 c. à café) ; limiter bouillons cubes (Maggi, Jumbo), poissons séchés salés, sauces industrielles. Utiliser épices locales (piment, gingembre, ail, oignon, persil, celeri).
\item \textbf{Sucres ajoutés} : $< 10\,\%$ AET (idéal $< 5\,\%$) $\approx$ \SI{25}{g/j} (6 c. à café). Supprimer sodas, jus industriels, bonbons, gâteaux industriels.
\end{enumerate}
\item \textbf{Activité Physique (OMS 2020)} :
\begin{enumerate}
\item \textbf{Adolescents (15-17 ans)} : \textbf{\SI{\ge 60}{min/jour}} d'activité \textbf{modérée à soutenue} (essoufflement léger à marqué). Majoritairement \textbf{endurance} (marche rapide, course, vélo, natation, football, danse, corde à sauter) + \textbf{renforcement musculaire/os} 3 fois/sem (port de charges, pompes, gainage, travaux champs).
\item \textbf{Adultes} : \textbf{\SI{\ge 150}{min/sem}} modérée \textbf{OU} \SI{75}{min/sem} soutenue + renforcement 2 fois/sem.
\item \textbf{Sédentarité} : Limiter temps assis/écrans loisirs $< \SI{2}{h/j}$ (hors scolaire/pro). Couper par pauses actives toutes les heures.
\end{enumerate}
\end{itemize}
\end{aretenir}

\courssub{Élaborer un programme de prévention personnalisé (Alimentation + Activité Physique)}

\begin{methode}[Construire un programme de prévention réaliste et mesurable]
\begin{enumerate}
\item \textbf{Diagnostic initial} : IMC + Tour de taille (si $> \SI{94}{cm}$ H / $> \SI{80}{cm}$ F $\rightarrow$ risque métabolique) + Enquête alimentaire (24h ou fréquence) + Niveau d'activité (IPAQ court) + Temps écrans + Antécédents familiaux.
\item \textbf{Objectifs SMART} : Spécifiques, Mesurables, Atteignables, Réalistes, Temporels.
Ex : « Perdre \SI{0,5}{à 1}{kg/mois} » (pas « maigrir vite ») ; « Marcher \SI{30}{min} 5 fois/sem » ; « Manger 3 fruits/jour ».
\item \textbf{Volet Alimentaire} : Planifier menus types (J1 à J7) respectant ANC. Liste de courses. Cuisson saine. Gestion fringales (fruit, yaourt, arachides nature).
\item \textbf{Volet Activité Physique} : Choisir activités \textbf{aimées} et \textbf{accessibles} (marche, danse, foot, corde, travaux champs, ménage dynamique). Planifier créneaux fixes dans agenda. Progressivité : \SI{10}{min} $\rightarrow$ \SI{30}{min} $\rightarrow$ \SI{60}{min}.
\item \textbf{Volet Comportemental} : Sommeil \SI{8}{à 9}{h} (régule ghréline/leptine) ; Gestion stress (respiration, prière, méditation) ; Soutien familial/entourage.
\item \textbf{Suivi \& Réévaluation (Mensuel)} : Pesée (même balance, même conditions) $\rightarrow$ Recalcul IMC $\rightarrow$ Tour de taille $\rightarrow$ Journal alimentaire/activité $\rightarrow$ Ajustement programme. Si échec 2 mois : Consultation médecin/nutritionniste.
\end{enumerate}
\end{methode}

\begin{exoresolu}[Programme de prévention pour Sandrine (Première TI, 17 ans)]
\textbf{Données} : $t = \SI{1,60}{m}$, $m = \SI{70}{kg}$, Écran \SI{6}{h/j}, Sport 0, Alimentation : beignets + sodas quotidiens.
\begin{enumerate}
\item Diagnostic IMC et analyse.
\item Objectif à 3 mois.
\item Programme hebdomadaire détaillé (Tableau).
\item Indicateurs de suivi.
\end{enumerate}
\tcblower
\textbf{1. Diagnostic.}
$\mathrm{IMC} = \dfrac{70}{1,60^2} = \dfrac{70}{2,56} = \SI{27,3}{kg/m^2}$ $\rightarrow$ \cle{Surpoids}.
Causes : Sédentarité majeure (écrans 6h, sport 0) + Déséquilibre alimentaire (excès lipides fritures, sucres rapides sodas, défaut fibres/légumes/protéines nobles).

\textbf{2. Objectif à 3 mois (Réaliste).}
Perte de \SI{2}{à 3}{kg} (cible \SI{67}{à 68}{kg}, IMC $\approx 26$). Stabiliser puis descendre vers zone normale (\SI{47,5}{à 63,5}{kg}) sur 6-12 mois. \textbf{Pas de régime restrictif} (danger croissance, troubles alimentaires).

\textbf{3. Programme Hebdomadaire Type.}
\begin{center}
\begin{tabularx}{\linewidth}{|l|X|X|X|}\hline
\rowcolor{popblueL} \textbf{Jour} & \textbf{Activité Physique (Cible $\ge$ 45-60 min)} & \textbf{Alimentation (Corrections ciblées)} & \textbf{Comportement Écrans/Sommeil} \\ \hline
Lundi & Marche rapide \SI{45}{min} (aller/retour école ou soir) & Petit-déj complet (pain + œuf + fruit) ; 0 beignet & Éteindre écrans \SI{1}{h} avant coucher \\ \hline
Mardi & Danse / Football / Corde à sauter \SI{60}{min} & Déjeuner : Riz + Sauce légumes (ndolé/eru) + Poisson grillé ; 0 soda & Devoirs sur table (pas lit) ; Pause active \SI{5}{min}/h \\ \hline
Mercredi & Marche \SI{30}{min} + Gainage \SI{10}{min} (3 series) & Goûter : Arachides grillées (\SI{30}{g}) + Mangue ; Eau \SI{1,5}{L} & Temps écran loisir $< \SI{2}{h}$ \\ \hline
Jeudi & Jogging léger / Vélo \SI{45}{min} & Dîner léger : Pâtes complètes + Légumes + Haricots ; Pas de friture & Lecture / Discussion famille \\ \hline
Vendredi & Sport collectif école/club \SI{60}{min} & Repas fête modéré : 1 part gâteau max, sinon fruit & Coucher \SI{22}{h}30 \\ \hline
Samedi & Aide ménage actif / Champ / Marche nature \SI{60}{min} & Préparer menus semaine (batch cooking) ; Courses liste & Activités sociales réelles \\ \hline
Dimanche & Repos actif : Promenade famille \SI{30}{min} & Repas équilibré plaisir ; Hydratation & Planifier semaine ; Coucher tôt \\ \hline
\end{tabularx}
\end{center}

\textbf{Justifications pédagogiques} :
\begin{itemize}
\item \textbf{Activité} : Variée (endurance, force, ludique), $\approx$ \SI{300}{min/sem} $>$ \SI{150}{min} rec. adulte, proche \SI{420}{min} (60min/j) rec. ado. Progression intégrée.
\item \textbf{Alimentation} : Substitutions concrètes (beignet $\rightarrow$ fruit/arachide ; soda $\rightarrow$ eau ; friture $\rightarrow$ grillé/vapeur ; riz blanc $\rightarrow$ riz étuvé/maïs). Produits locaux, peu coûteux. 5 fruits/légumes visés.
\item \textbf{Écrans/Sommeil} : Règle « 1h sans écran avant dormir » améliore sommeil $\rightarrow$ régule appétit.
\end{itemize}

\textbf{4. Suivi Mensuel (Tableau de bord).}
\begin{center}
\begin{tabular}{|c|c|c|c|c|}\hline
\rowcolor{popblueL} \textbf{Mois} & \textbf{Masse (kg)} & \textbf{IMC} & \textbf{Tour taille (cm)} & \textbf{Action si stagnation} \\ \hline
M0 (Début) & 70,0 & 27,3 & \textit{à mesurer} & — \\ \hline
M1 & 69,0 & 26,9 & $\downarrow$ & Continuer \\ \hline
M2 & 68,2 & 26,6 & $\downarrow$ & Continuer \\ \hline
M3 & 67,5 & 26,3 & $\downarrow$ & \textbf{Objectif atteint} : Maintenir + viser -1 kg/mois suivant \\ \hline
\end{tabular}
\end{center}
Si pas de perte à M2 : Réduire portions féculents \SI{20}{\%}, augmenter marche \SI{15}{min}, vérifier sommeil/stress, consulter infirmier scolaire / médecin.
\end{exoresolu}

\courssub{Exploiter des données épidémiologiques ou expérimentales (Savoir-faire 5)}

\begin{methode}[Analyser un tableau ou un graphique de données de santé publique]
\begin{enumerate}
\item \textbf{Identifier les variables} : Lignes (groupes), Colonnes (indicateurs), Unités (\%, n, moyenne, IMC, glycémie...).
\item \textbf{Extraire les valeurs brutes} et calculer les \textbf{fréquences, pourcentages, moyennes} demandés.
\item \textbf{Comparer les groupes} (exposés vs non exposés ; avant vs après ; garçons vs filles).
\item \textbf{Formuler la relation observée} : « Le groupe [X] présente une fréquence de [maladie] \textbf{supérieure/inférieure} de [facteur X] par rapport au groupe [Y] ». Utiliser « \textbf{est associé à} » (corrélation), pas « cause » sans mécanisme.
\item \textbf{Interpréter biologiquement} : Mobiliser les mécanismes vus en cours (bilan énergétique, résistance insuline, athérome).
\item \textbf{Conclure sur la prévention} : Quelle action prioritaire ? (Primaire : éducation ; Secondaire : dépistage ; Tertiaire : prise en charge).
\end{enumerate}
\end{methode}

\begin{exoresolu}[Enquête surpoids/activité dans un lycée de Garoua (400 élèves)]
\begin{center}
\begin{tabular}{|l|c|c|}\hline
\rowcolor{popblueL} \textbf{Groupe} & \textbf{Sport $\ge$ 3 h/sem.} & \textbf{Sport $<$ 1 h/sem.} \\ \hline
Effectif total & 150 & 250 \\ \hline
Élèves en surpoids ou obèses (IMC $\ge$ 25) & 12 & 75 \\ \hline
\end{tabular}
\end{center}
\begin{enumerate}
\item Calcule la fréquence du surpoids/obésité dans chaque groupe (en \%). 
\item Compare les deux groupes. Quelle relation observes-tu ?
\item Propose une explication biologique et une mesure de prévention primaire.
\end{enumerate}
\tcblower
\textbf{1. Fréquences (Prévalence).}
\begin{align*}
\text{Groupe Actif } (\ge 3\text{h/sem}) &: \dfrac{12}{150} \times 100 = \mathbf{8\,\%} \\
\text{Groupe Sédentaire } (< 1\text{h/sem}) &: \dfrac{75}{250} \times 100 = \mathbf{30\,\%}
\end{align*}

\textbf{2. Comparaison et Relation.}
La prévalence du surpoids/obésité est de \textbf{30\% chez les sédentaires} contre \textbf{8\% chez les actifs}.
Le risque relatif (RR) = $\frac{30}{8} = 3,75$. Un élève sédentaire a \textbf{3,75 fois plus de risque} d'être en surpoids/obèse qu'un élève actif.
Il existe une \cle{corrélation négative forte} entre le niveau d'activité physique et la fréquence du surpoids/obésité dans cette population.

\textbf{3. Explication biologique et Prévention.}
\textbf{Mécanisme} : L'activité physique augmente la dépense énergétique de repos (masse musculaire) et d'effort, améliore la sensibilité à l'insuline (translocation GLUT4 indépendante de l'insuline), favorise l'oxydation des lipides et limite l'hypertrophie des adipocytes viscéraux. La sédentarité fait l'inverse : bilan positif, stockage graisseux, inflammation bas bruit, résistance à l'insuline.
\textbf{Prévention primaire} : Intégrer \SI{45}{min} d'EPS obligatoire \textbf{quotidienne} (pas hebdomadaire) dans l'emploi du temps ; aménager cours de récréation pour jeux actifs ; sensibiliser parents sur marche école-maison ; interdire vente boissons sucrées/fritures dans l'enceinte scolaire (arrêté municipal).
\end{exoresolu}

\begin{exercice}[Synthèse : Cas de Paul, 16 ans, Première C]
Paul mesure \SI{1,75}{m} pour \SI{85}{kg}. Il ne fait pas de sport, mange au restaurant scolaire (souvent frites + viande grasse + pain blanc), goûte pain-chocolat + soda, dîne pâtes + sauce tomate industrielle + fromage. Il dort \SI{6}{h}/nuit, écrans \SI{4}{h}/jour. Antécédents : père diabétique T2, mère HTA.
\begin{enumerate}
\item Calcule son IMC. Classe-le. Quels sont ses risques immédiats et à 10 ans ?
\item Sa glycémie à jeun est de \SI{1,10}{g/L}. Est-il diabétique ? Comment interpréter cette valeur (norme : $< \SI{1,00}{g/L}$) ? Quel terme médical utilise-t-on ?
\item Calcule ses besoins énergétiques approximatifs (garçon 16 ans, sédentaire). Sa ration journalière apporte \SI{3200}{kcal}. Quel est le bilan ? Combien de kg de graisse stocke-t-il par mois si 1 kg graisse $\approx \SI{7700}{kcal}$ ?
\item Propose un programme de prévention sur 1 semaine (tableau Jour / Activité / Repas clés / Objectif sommeil/écran) avec 3 mesures alimentaires précises et 3 mesures d'activité précises.
\item Le directeur veut réduire le surpoids au lycée. À partir de l'enquête de Garoua, quel argument chiffré donnes-tu pour imposer 1h d'EPS quotidienne ?
\end{enumerate}
\end{exercice}

\begin{corrige}[Exercice : Cas de Paul]
\textbf{1.} $\mathrm{IMC} = \frac{85}{1,75^2} = \frac{85}{3,0625} = \SI{27,8}{kg/m^2}$ $\rightarrow$ \textbf{Surpoids}. Risques immédiats : essoufflement à l'effort, douleurs articulaires (genoux), fatigue, stigmatisation. Risques 10 ans (si inchangé) : Diabète T2 (antécédent paternel + surpoids), HTA (antécédent maternel), dyslipidémie, stéatose hépatique, syndrome métabolique, infarctus/AVC précoces.
\textbf{2.} Glycémie \SI{1,10}{g/L} : \textbf{Pas diabétique} (seuil $\ge \SI{1,26}{g/L}$). Mais $> \SI{1,00}{g/L}$ $\rightarrow$ \cle{Hyperglycémie modérée à jeun} (ou \textbf{Prédiabète} / Intolérance au glucose). Signal d'alarme majeur : résistance à l'insuline installée. Réversible par hygiène de vie.
\textbf{3.} Besoin $\approx \SI{2500}{kcal/j}$ (garçon 16 ans sédentaire). Excès = \SI{700}{kcal/j}. Sur 30 jours = \SI{21000}{kcal}. Stockage graisse = $\frac{21000}{7700} \approx \SI{2,7}{kg/mois}$ ! Prise de poids rapide prévisible.
\textbf{4.} Programme (extrait) : Lundi : Marche 45min / Midi: Riz+légumes+poisson / 0 soda / Coucher 22h30. Mardi : Corde 20min+Gainage 10min / Goûter: Fruit+Arachides / Écran $<$2h. Mercredi : Foot 60min / Dîner léger légumineuses. Jeudi : Marche 30min / Petit-déj complet. Vendredi : Danse 45min / 0 friture. Samedi : Ménage actif 60min / Prépa menus. Dimanche : Randonnée 45min / Repos.
Mesures alimentaires : 1) Remplacer soda par eau/citronnade sans sucre. 2) Imposer légumes verts à midi et soir. 3) Remplacer pain blanc/frites par plantain bouilli/riz étuvé/igname.
Mesures activité : 1) Marche aller-retour école (\SI{30}{min}). 2) Inscription club sport (foot/danse) 2x/sem. 3) Corde à sauter \SI{10}{min} matin avant douche.
\textbf{5.} Argument : « L'enquête de Garoua (400 élèves) montre que la prévalence du surpoids/obésité est de \textbf{30\% chez les élèves faisant $<$ 1h de sport/sem} contre \textbf{8\% chez ceux faisant $\ge$ 3h/sem} (RR = 3,75). Passer l'EPS de 2h/sem à 5h/sem (quotidienne) ferait basculer la majorité des élèves dans le groupe « protecteur » ($\ge$ 3h), divisant potentiellement par 3,75 le nombre de nouveaux cas de surpoids. C'est un investissement santé publique rentable (coût maladie évité). »
\end{corrige}

\begin{attention}[Les 5 erreurs fatales au Probatoire (SVT - Éducation à la Santé)]
\begin{enumerate}
\item \textbf{IMC : Oublier la conversion cm $\rightarrow$ m.} Écrire $\frac{82}{170^2}$ au lieu de $\frac{82}{1,70^2}$. Résultat absurde ($0,0028$). \textbf{Sanction : 0 pt au calcul.} $\rightarrow$ \textbf{Règle} : Écrire explicitement $t = \SI{1,70}{m}$ avant la formule.
\item \textbf{Équivalences énergétiques : Confondre Lipides (9 kcal/g) et Glucides/Protéines (4 kcal/g).} Inverser les coefficients fausse tout l'AET et les pourcentages. $\rightarrow$ \textbf{Règle} : Mémoriser « 9 pour le gras, 4 pour le sucre et la viande ».
\item \textbf{IMC sans interprétation ni risque.} Donner juste le chiffre (ex: « IMC = 28,4 »). \textbf{Sanction : Note incomplète.} $\rightarrow$ \textbf{Règle} : Toujours écrire : « IMC = 28,4 kg/m$^2$ $\rightarrow$ Surpoids $\rightarrow$ Risque Diabète T2, HTA, CVD ».
\item \textbf{Confondre Diabète Type 1 et Type 2.} Écrire « Manque d'insuline » ou « Pancréas ne marche plus » pour un obèse de 45 ans. $\rightarrow$ \textbf{Règle} : DT2 = \textbf{Résistance à l'insuline} (cellules ne répondent pas) + Hyperinsulinémie initiale. DT1 = \textbf{Déficit absolu en insuline} (auto-immunité).
\item \textbf{Recommandations vagues.} « Manger équilibré, faire du sport ». \textbf{Sanction : 0 pt.} $\rightarrow$ \textbf{Règle} : Chiffrer et contextualiser : « 5 fruits/légumes/jour (ndolé, eru, ananas) », « \SI{1,5}{L} eau/jour », « \SI{60}{min} marche rapide/jour », « Supprimer sodas/beignets quotidiens », « Coucher \SI{22}{h}30, 0 écran 1h avant ».
\end{enumerate}
\end{attention}

\newpage

\coursec{Exercices}

\courssub{Je vérifie mes connaissances}

\begin{exercice}[QCM : Rappels de cours]
Pour chaque affirmation, une seule proposition est correcte. Recopie le numéro et la lettre correspondante.
\begin{enumerate}
\item Le seuil d'IMC définissant l'obésité chez l'adulte selon l'OMS est :
\begin{itemize}
\item[A.] 25 kg/m²
\item[B.] 30 kg/m²
\item[C.] 18,5 kg/m²
\item[D.] 40 kg/m²
\end{itemize}
\item La répartition recommandée de l'apport énergétique total (AET) pour un adolescent est :
\begin{itemize}
\item[A.] Glucides 50-55\%, Lipides 30-35\%, Protéines 10-15\%
\item[B.] Glucides 30-35\%, Lipides 50-55\%, Protéines 10-15\%
\item[C.] Glucides 10-15\%, Lipides 30-35\%, Protéines 50-55\%
\item[D.] Glucides 55-60\%, Lipides 10-15\%, Protéines 30-35\%
\end{itemize}
\item La durée minimale d'activité physique d'intensité modérée à soutenue recommandée par l'OMS pour les adolescents (10-17 ans) est :
\begin{itemize}
\item[A.] 30 minutes par jour
\item[B.] 45 minutes par jour
\item[C.] 60 minutes par jour
\item[D.] 90 minutes par jour
\end{itemize}
\item Parmi les conséquences suivantes, laquelle n'est \textbf{pas} directement favorisée par la sédentarité et la surcharge pondérale ?
\begin{itemize}
\item[A.] Diabète de type 2
\item[B.] Maladies cardiovasculaires (HTA, AVC)
\item[C.] Ostéoporose juvénile
\item[D.] Certains cancers (côlon, sein, endomètre)
\end{itemize}
\end{enumerate}
\end{exercice}

\begin{exercice}[Vrai ou Faux ? Justifie]
Indique si l'affirmation est vraie (V) ou fausse (F) et justifie brièvement ta réponse en mobilisant une notion du cours.
\begin{enumerate}
\item « Un adolescent avec un IMC de 22 kg/m² est obligatoirement en bonne santé métabolique. »
\item « La sédentarité se définit uniquement par l'absence de pratique sportive en club. »
\item « Le diabète de type 2 résulte d'une destruction auto-immune des cellules $\beta$ pancréatiques. »
\item « Une ration alimentaire équilibrée doit apporter exactement la même quantité d'énergie chaque jour. »
\end{enumerate}
\end{exercice}

\begin{exercice}[Définitions et notions clés]
Donne une définition précise et scientifique pour chacun des termes suivants :
\begin{enumerate}
\item Indice de Masse Corporelle (IMC).
\item Équilibre alimentaire (ou ration équilibrée).
\item Sédentarité (au sens physiologique et comportemental).
\item Métabolisme de base (MB).
\end{enumerate}
\end{exercice}

\begin{exercice}[Calculs directs d'IMC et classification]
Trois élèves de Première passent une visite médicale. Leurs mensurations sont :
\begin{itemize}
\item Élève A : masse $m_A = 58$ kg, taille $t_A = 1,65$ m.
\item Élève B : masse $m_B = 85$ kg, taille $t_B = 1,70$ m.
\item Élève C : masse $m_C = 48$ kg, taille $t_C = 1,60$ m.
\end{itemize}
\begin{enumerate}
\item Calcule l'IMC de chaque élève (formule : $\text{IMC} = \frac{m}{t^2}$, résultat arrondi au dixième).
\item Classe chaque élève selon la classification OMS pour adultes (rappel : $<18,5$ = insuffisance ; $18,5-24,9$ = normal ; $25-29,9$ = surpoids ; $\ge 30$ = obésité).
\item Pour chaque élève, propose un conseil de prévention adapté (une phrase).
\end{enumerate}
\end{exercice}

\courssub{Je m'entraîne}

\begin{exercice}[Évaluation d'une ration alimentaire : le déjeuner type]
On souhaite vérifier si le déjeuner suivant couvre correctement les besoins d'un déjeuner (environ 35-40\% de l'AET) pour un adolescent dont l'AET est de \SI{2500}{kcal}.
\textbf{Composition du repas :}
\begin{itemize}
\item \SI{150}{g} de riz cuit
\item \SI{50}{g} d'arachides grillées
\item \SI{20}{mL} d'huile de cuisson (densité $\approx 0,92$ g/mL)
\item \SI{100}{g} de poisson grillé (type capitaine)
\end{itemize}
\textbf{Table de composition (valeurs pour 100 g ou 100 mL) :}
\begin{center}
\begin{tabular}{|l|c|c|c|c|}\hline
Aliment & Eau (g) & Protéines (g) & Lipides (g) & Glucides (g) \\ \hline
Riz cuit & 68 & 2,7 & 0,3 & 28 \\
Arachides grillées & 2 & 25 & 49 & 16 \\
Huile végétale & 0 & 0 & 100 & 0 \\
Poisson grillé & 70 & 24 & 2 & 0 \\ \hline
\end{tabular}
\end{center}
Coefficients énergétiques : Protéines = \SI{4}{kcal/g} ; Lipides = \SI{9}{kcal/g} ; Glucides = \SI{4}{kcal/g}.
\begin{enumerate}
\item Calcule l'apport en protéines, lipides, glucides (en g) et l'énergie totale (en kcal) de ce repas.
\item Calcule le pourcentage de l'AET (\SI{2500}{kcal}) apporté par ce repas. Ce pourcentage est-il dans la norme pour un déjeuner ?
\item Calcule la répartition énergétique (\% Protides, \% Lipides, \% Glucides) de ce repas. Compare aux recommandations (P : 10-15\%, L : 30-35\%, G : 50-55\%). Commente.
\end{enumerate}
\end{exercice}

\begin{exercice}[Prise de masse grasse annuelle liée à une surconsommation]
Un adolescent de 16 ans a un apport énergétique équilibré par rapport à ses besoins, mais il consomme \textbf{chaque jour} en \textbf{supplément} : 2 beignets (haricots) apportant \SI{180}{kcal} chacun et une boisson sucrée industrielle de \SI{33}{cL} apportant \SI{220}{kcal}.
On admet que \SI{1}{kg} de tissu adipeux (masse grasse) stocke environ \SI{7700}{kcal}.
\begin{enumerate}
\item Calcule l'excédent énergétique quotidien (en kcal).
\item Calcule l'excédent énergétique annuel (en kcal). (Année = 365 jours).
\item Estime la prise de masse grasse (en kg) sur un an si cet excédent est entièrement stocké.
\item Ce calcul théorique est-il toujours vérifié dans la réalité ? Justifie en citant au moins deux facteurs physiologiques ou comportementaux.
\end{enumerate}
\end{exercice}

\begin{exercice}[Interprétation de courbes de croissance et IMC]
Le carnet de santé d'un adolescent montre les points suivants pour son IMC au fil de l'âge :
\begin{center}
\begin{tabular}{|c|c|c|c|c|c|}\hline
Âge (ans) & 10 & 12 & 14 & 16 & 17 \\ \hline
IMC (kg/m²) & 17,2 & 19,5 & 23,8 & 27,5 & 28,9 \\ \hline
\end{tabular}
\end{center}
(Référence OMS : Surpoids si IMC > +1 ET ; Obésité si IMC > +2 ET par rapport à la médiane pour l'âge/sexe).
\begin{enumerate}
\item Trace schématiquement l'évolution de l'IMC en fonction de l'âge (axe des abscisses : âge ; axe des ordonnées : IMC).
\item En te basant sur les valeurs de référence approximatives pour un garçon (P50 $\approx$ 16,5 à 10 ans ; 18,5 à 14 ans ; 21,5 à 17 ans ; +1 ET $\approx$ P85 ; +2 ET $\approx$ P97), détermine à partir de quel âge cet adolescent bascule en surpoids puis en obésité.
\item Quels risques métaboliques et cardiovasculaires cette trajectoire fait-elle courir à cet adolescent s'il n'y a pas de prise en charge ?
\item Propose deux axes d'intervention prioritaires (alimentation, activité physique) pour inverser la tendance.
\end{enumerate}
\end{exercice}

\begin{exercice}[Construire un programme de prévention personnalisé]
Fatou, 16 ans, mesure 1,62 m pour 68 kg. Elle ne pratique aucune activité physique hors cours d'EPS (2h/semaine). Son alimentation type : petit-déj : pain + chocolat + lait sucré ; Déjeuner : riz + sauce grasse (huile, viande) + peu de légumes ; Dîner : pâtes + beignets + soda ; Collations : chips, bonbons.
\begin{enumerate}
\item Calcule son IMC et classe-le.
\item Identifie trois erreurs majeures dans son alimentation par rapport aux repères PNNS (Programme National Nutrition Santé) adaptés au Cameroun.
\item Élabore un programme journalier simple pour Fatou comprenant :
\begin{itemize}
\item Un menu équilibré sur une journée (4 repas) avec des aliments locaux accessibles.
\item Un plan d'activité physique atteignant les 60 min/jour (préciser type, durée, intensité).
\end{itemize}
\item Explique en quoi ce programme permet de prévenir l'apparition du diabète de type 2.
\end{enumerate}
\end{exercice}

\courssub{Je me prépare au Probatoire}

\begin{exercice}[Sujet type Probatoire : Cas clinique complet — « Le cas de Stéphane »]
\textbf{Document 1 : Fiche de renseignements}
Stéphane, 17 ans, élève en Première C. Taille : 1,75 m. Masse : 92 kg. Tension artérielle : 135/85 mmHg (norme < 120/80). Tour de taille : 98 cm (seuil risque > 94 cm homme).
\textbf{Document 2 : Rappel alimentaire 24h (jour de classe)}
\begin{itemize}
\item 6h30 : 200 mL lait entier sucré + 100 g pain blanc + 30 g beurre + 20 g confiture.
\item 12h30 (cantine) : 200 g riz blanc + 100 g ndolé (feuilles, beaucoup d'huile, viande grasse, crevettes) + 1 mangue.
\item 16h00 : 1 pain au chocolat (boulangerie) + 1 canette boisson gazeuse (33 cL).
\item 19h30 : 250 g pâtes + sauce tomate industrielle + 100 g fromage râpé + 1 yaourt sucré.
\end{itemize}
\textbf{Document 3 : Bilan sanguin (à jeun)}
Glycémie : 1,15 g/L (norme 0,70 - 1,10). HDL-cholestérol : 0,35 g/L (norme > 0,40). Triglycérides : 1,80 g/L (norme < 1,50).
\textbf{Document 4 : Mode de vie}
Trajet école-domicile en taxi-moto. 0 min de sport hors EPS. Écrans : 4h/jour. Sommeil : 6h/nuit.
\textbf{Questions :}
\begin{enumerate}
\item Calcule l'IMC de Stéphane. Interprète-le selon la classification OMS. Calcule son tour de taille / taille. Que conclut-on sur la répartition de l'adiposité ?
\item À partir du Document 2, estime l'apport énergétique approximatif de la journée (utilise les coefficients : P=4, L=9, G=4 ; estime les quantités de macronutriments pour chaque aliment). Compare aux besoins estimés (\SI{2800}{kcal} pour garçon 17 ans actif). Commente l'équilibre alimentaire (répartition P/L/G, fibres, densité énergétique).
\item Analyse le Document 3. Quel syndrome métabolique se dessine ? Relie les anomalies biologiques à l'hyperinsulinisme compensateur et à la résistance à l'insuline.
\item En synthèse, Stéphane cumule des facteurs de risque. Cite-en 4 (anthropométriques, biologiques, comportementaux) et nomme la pathologie majeure qui le menace à moyen terme.
\item Rédige une fiche de conseils personnalisée pour Stéphane (10 lignes max) structurée en 3 volets : Alimentation (3 mesures concrètes), Activité physique (objectif progressif), Hygiène de vie (sommeil, écrans).
\end{enumerate}
\end{exercice}

\begin{exercice}[Sujet type Probatoire : Analyse de documents — « Transition nutritionnelle au Cameroun »]
\textbf{Document 1 : Évolution de la prévalence de l'obésité (Enquête DHS Cameroun)}
\begin{center}
\begin{tabular}{|c|c|c|}\hline
Année & Femmes 15-49 ans (\% obèses IMC $\ge$ 30) & Hommes 15-49 ans (\% obèses IMC $\ge$ 30) \\ \hline
2004 & 10,5 & 3,2 \\
2011 & 15,2 & 5,8 \\
2018 & 20,1 & 8,4 \\ \hline
\end{tabular}
\end{center}
\textbf{Document 2 : Témoignage}
« Avant, au village, on mangeait du manioc, des feuilles, du poisson fumé. On marchait 5 km pour aller au champ. Ici à Douala, je mange du riz, de la sauce arachide très grasse, des beignets à tous les coins de rue. Je prends le bendy-skin pour 500 m. Je dors tard, stressé par le travail. En 3 ans, j'ai pris 18 kg. » — Paul, 32 ans, cadre bancaire.
\textbf{Document 3 : Mécanisme physiopathologique du diabète de type 2 (schéma simplifié)}
Surcharge nutritionnelle $\rightarrow$ Hypertrophie adipocytes $\rightarrow$ Inflammation bas bruit (cytokines : TNF-$\alpha$, IL-6) $\rightarrow$ Résistance à l'insuline (foie, muscle, tissu adipeux) $\rightarrow$ Hyperinsulinisme compensateur $\rightarrow$ Épuisement cellules $\beta$ $\rightarrow$ Hyperglycémie chronique.
\textbf{Questions :}
\begin{enumerate}
\item Exploite le Document 1 : Calcule l'augmentation relative (\% d'augmentation) de l'obésité chez les femmes et chez les hommes entre 2004 et 2018. Quelle tendance observes-tu ? Propose deux hypothèses explicatives liées à la transition nutritionnelle et épidémiologique au Cameroun.
\item À partir du Document 2, identifie les changements dans les 3 piliers : Alimentation, Activité physique, Environnement/Comportement. Comment ces changements expliquent-ils le bilan énergétique positif de Paul ?
\item En t'appuyant sur le Document 3, explique le lien entre l'obésité abdominale (adiposité viscérale) et l'installation de l'hyperglycémie de type 2. Utilise les termes : adipokines, résistance à l'insuline, néoglucogenèse hépatique, captage musculaire du glucose.
\item Le diabète de type 2 est une « maladie silencieuse » au début. Cite 3 complications chroniques (micro et macrovasculaires) et l'examen de dépistage annuel recommandé pour la néphropathie.
\item En tant qu'agent de santé publique, propose 3 mesures de prévention primaire (populationnelle) adaptées au contexte camerounais pour freiner cette épidémie (ex : fiscalité, aménagement urbain, éducation, régulation publicitaire).
\end{enumerate}
\end{exercice}

\begin{exercice}[Sujet type Probatoire : Rédaction guidée — « Lettre à un ami »]
\textbf{Consigne :} Ton ami Armand, 17 ans, IMC 29 kg/m², ne fait pas de sport, mange « n'importe quoi » (fast-food, sodas), dort 5h/nuit. Il te dit : « De toute façon, le diabète c'est pour les vieux, et moi je me sens bien. »
Rédige un message de prévention (environ 15 lignes, style courrier/mail) pour le convaincre de changer.
\textbf{Contraintes obligatoires :} Tu dois impérativement utiliser et expliquer correctement les 4 notions suivantes dans ton texte :
\begin{enumerate}
\item La notion d'IMC et son interprétation (seuil surpoids/obésité).
\item La notion d'équilibre alimentaire (qualité, répartition, densité nutritionnelle).
\item Le mécanisme physiopathologique menant au diabète de type 2 (résistance à l'insuline, épuisement pancréatique).
\item La recommandation OMS des 60 min/jour d'activité modérée à intense et ses bénéfices métaboliques (sensibilisation à l'insuline, dépense énergétique).
\end{enumerate}
Ton texte sera évalué sur la rigueur scientifique, la clarté du message et le ton bienveillant mais ferme.
\end{exercice}

\newpage

\coursec{Corrigés des exercices}

\courssub{Je vérifie mes connaissances}

\begin{corrige}[Exercice 1 — QCM : Rappels de cours]
\begin{enumerate}
\item \textbf{B.} Selon la classification de l'OMS pour les adultes, l'obésité est définie par un IMC $\ge 30$ kg/m². Le seuil de 25 kg/m² correspond au surpoids, 18,5 kg/m² à la maigreur, et 40 kg/m² à l'obésité morbide (classe III).
\item \textbf{A.} Les recommandations nutritionnelles (ANC, PNNS) pour l'adolescent prévoient une répartition de l'apport énergétique total (AET) en : Glucides 50--55\% (dont sucres ajoutés $<10\%$), Lipides 30--35\% (dont AGS $<10\%$), Protéines 10--15\%.
\item \textbf{C.} L'OMS recommande aux adolescents (10--17 ans) au moins \textbf{60 minutes par jour} d'activité physique d'intensité modérée à soutenue, majoritairement d'endurance, et des activités de renforcement musculaire et osseux au moins 3 fois par semaine.
\item \textbf{C.} L'ostéoporose juvénile est principalement liée à un déficit de pic de masse osseuse (apport calcique/vitamine D insuffisant, absence de contrainte mécanique \emph{portante}, troubles hormonaux). La sédentarité (manque de contrainte mécanique) est un facteur de risque, mais la \textbf{surcharge pondérale} (excès de masse grasse) est plutôt protectrice sur la densité minérale osseuse par effet de charge mécanique et sécrétion d'œstrogènes par le tissu adipeux. En revanche, obésité, diabète de type 2, maladies cardiovasculaires et certains cancers (côlon, sein, endomètre) sont des conséquences directes et majeures de la sédentarité et de la surcharge pondérale.
\end{enumerate}
\end{corrige}

\begin{corrige}[Exercice 2 — Vrai ou Faux ? Justifie]
\begin{enumerate}
\item \textbf{Faux.} Un IMC de 22 kg/m² se situe dans la plage « normale » (18,5--24,9), mais l'IMC ne distingue pas la masse grasse de la masse maigre. Un adolescent peut avoir un IMC normal mais une adiposité élevée (obésité métabolique de poids normal, phénotype TOFI -- \emph{Thin Outside Fat Inside}) avec résistance à l'insuline, dyslipidémie ou stéatose hépatique. La santé métabolique s'évalue aussi par le tour de taille, la tension artérielle, le profil lipidique et la glycémie.
\item \textbf{Faux.} La sédentarité se définit physiologiquement par un \textbf{niveau de dépense énergétique très faible} ($\le 1,5$ MET) en position assise ou allongée pendant les heures d'éveil (temps d'écran, transport motorisé, travail bureau). Elle est distincte de l'inactivité physique (non-respect des recommandations d'AP). Un adolescent peut faire 2h de sport en club (actif) mais rester assis 10h/jour (sédentaire).
\item \textbf{Faux.} C'est la définition physiopathologique du \textbf{diabète de type 1} (maladie auto-immune). Le diabète de type 2 résulte d'une \textbf{résistance à l'insuline} (insulinorésistance) périphérique (foie, muscle, tissu adipeux) suivie, à long terme, d'un épuisement des cellules $\beta$ pancréatiques (insulinopénie relative).
\item \textbf{Faux.} L'équilibre alimentaire s'apprécie \textbf{sur la semaine}, pas sur une journée unique. Les besoins énergétiques fluctuent quotidiennement selon l'activité physique, la croissance, la température, le stress. Une ration équilibrée couvre les besoins \emph{moyens} sur plusieurs jours, avec une répartition qualitative stable (P/L/G, fibres, micronutriments).
\end{enumerate}
\end{corrige}

\begin{corrige}[Exercice 3 — Définitions et notions clés]
\begin{enumerate}
\item \textbf{Indice de Masse Corporelle (IMC) :} Indice anthropométrique simple estimant la corpulence d'un adulte, calculé par la formule $\text{IMC} = \frac{m}{t^2}$ où $m$ est la masse en kilogrammes (kg) et $t$ la taille en mètres (m). Unité : kg/m². Chez l'enfant/adolescent, l'interprétation nécessite les courbes de référence âge/sexe (z-scores ou percentiles).
\item \textbf{Équilibre alimentaire (ration équilibrée) :} Ration alimentaire qui, sur une période donnée (généralement la semaine), \textbf{couvre intégralement les besoins nutritionnels} (énergie, macronutriments, micronutriments, fibres, eau) sans excès, tout en respectant la répartition recommandée des apports énergétiques (Glucides 50--55\%, Lipides 30--35\%, Protéines 10--15\%), la diversité alimentaire (5 groupes d'aliments), la distribution sur 3 à 4 repas/jour et le plaisir gustatif.
\item \textbf{Sédentarité (sens physiologique et comportemental) :} Ensemble des comportements d'éveil caractérisés par une \textbf{dépense énergétique $\le 1,5$ MET} (équivalent métabolique) en position assise, semi-allongée ou allongée (temps d'écran passif, transport motorisé, travail scolaire/bureau, lecture). Elle s'évalue par le \textbf{temps passé assis} (objectif : $< 2$h/jour loisir écran chez l'adolescent) et se distingue de l'inactivité physique (insuffisance d'AP modérée à intense).
\item \textbf{Métabolisme de base (MB ou BMR -- Basal Metabolic Rate) :} Quantité minimale d'énergie (en kcal/j ou kJ/j) dépensée par l'organisme au \textbf{repos complet} (veille, jeûne de 12h, thermoneutralité 22°C, sans stress) pour maintenir les fonctions vitales (ventilation, circulation, activité cérébrale, turnover cellulaire, homéostasie ionique). Il représente 60--70\% de la dépense énergétique totale (DET). Principaux déterminants : masse maigre (muscle), âge, sexe, thyroïde, génétique.
\end{enumerate}
\end{corrige}

\begin{corrige}[Exercice 4 — Calculs directs d'IMC et classification]
\begin{enumerate}
\item \textbf{Calculs (arrondi au dixième) :}
\begin{align*}
\text{Élève A : } \text{IMC}_A &= \frac{58}{1,65^2} = \frac{58}{2,7225} = 21,30 \approx \mathbf{21,3\ kg/m^2} \\
\text{Élève B : } \text{IMC}_B &= \frac{85}{1,70^2} = \frac{85}{2,89} = 29,41 \approx \mathbf{29,4\ kg/m^2} \\
\text{Élève C : } \text{IMC}_C &= \frac{48}{1,60^2} = \frac{48}{2,56} = 18,75 \approx \mathbf{18,8\ kg/m^2}
\end{align*}
\item \textbf{Classification OMS (adultes) :}
\begin{itemize}
\item Élève A (21,3) : \textbf{Normal} (18,5--24,9).
\item Élève B (29,4) : \textbf{Surpoids} (25,0--29,9) -- limite obésité.
\item Élève C (18,8) : \textbf{Normal} (18,5--24,9) -- limite insuffisance pondérale.
\end{itemize}
\item \textbf{Conseils de prévention adaptés :}
\begin{itemize}
\item Élève A : « Maintenir cette corpulence par une alimentation variée, équilibrée et la pratique régulière de 60 min/jour d'activité physique modérée à intense. »
\item Élève B : « Consulter l'infirmier(e) scolaire ou un médecin pour un bilan ; réduire les apports en graisses ajoutées et sucres rapides, augmenter les légumes/fruits ; viser 60 min/jour d'AP (marche rapide, sport loisir) pour stabiliser puis faire baisser l'IMC. »
\item Élève C : « Surveiller la courbe de croissance (poids/taille) ; s'assurer d'un apport énergétique et protéique suffisant pour la croissance pubertaire (3 repas + 1 collation équilibrés) ; ne pas restreindre l'alimentation sans avis médical. »
\end{itemize}
\end{enumerate}
\end{corrige}

\courssub{Je m'entraîne}

\begin{corrige}[Exercice 5 — Évaluation d'une ration alimentaire : le déjeuner type]
\textbf{Données :} AET = 2500 kcal. Repas visé = 35--40\% AET = 875--1000 kcal.
\textbf{1. Apports en macronutriments et énergie du repas :}
\begin{center}
\begin{tabular}{|l|c|c|c|c|c|c|}\hline
Aliment & Quantité & Protéines (g) & Lipides (g) & Glucides (g) & Énergie (kcal) \\ \hline
Riz cuit & 150 g & $150 \times 2,7/100 = 4,05$ & $150 \times 0,3/100 = 0,45$ & $150 \times 28/100 = 42,0$ & $4,05\times4 + 0,45\times9 + 42\times4 = 16,2 + 4,05 + 168 = 188,25$ \\
Arachides grillées & 50 g & $50 \times 25/100 = 12,5$ & $50 \times 49/100 = 24,5$ & $50 \times 16/100 = 8,0$ & $12,5\times4 + 24,5\times9 + 8\times4 = 50 + 220,5 + 32 = 302,5$ \\
Huile végétale & 20 mL ($\approx 18,4$ g) & 0 & $18,4 \times 100/100 = 18,4$ & 0 & $18,4 \times 9 = 165,6$ \\
Poisson grillé & 100 g & $100 \times 24/100 = 24,0$ & $100 \times 2/100 = 2,0$ & 0 & $24\times4 + 2\times9 = 96 + 18 = 114$ \\ \hline
\textbf{TOTAL} & & \textbf{40,55 g} & \textbf{45,35 g} & \textbf{50,0 g} & \textbf{770,35 kcal} \\ \hline
\end{tabular}
\end{center}

\textbf{2. Pourcentage de l'AET :}
\[
\frac{770,35}{2500} \times 100 = \mathbf{30,8\%}
\]
Ce pourcentage est \textbf{inférieur à la norme} pour un déjeuner (35--40\% de l'AET). L'apport énergétique est insuffisant pour couvrir les besoins de l'après-midi (risque de fringale, grignotage sucré/gras).

\textbf{3. Répartition énergétique du repas :}
\begin{align*}
\% \text{Protéines} &= \frac{40,55 \times 4}{770,35} \times 100 = \frac{162,2}{770,35} \times 100 = \mathbf{21,1\%} \\
\% \text{Lipides} &= \frac{45,35 \times 9}{770,35} \times 100 = \frac{408,15}{770,35} \times 100 = \mathbf{53,0\%} \\
\% \text{Glucides} &= \frac{50,0 \times 4}{770,35} \times 100 = \frac{200}{770,35} \times 100 = \mathbf{26,0\%}
\end{align*}
\textbf{Comparaison aux recommandations (P 10--15\% / L 30--35\% / G 50--55\%) :}
\begin{itemize}
\item \textbf{Protéines :} 21,1\% $\gg$ 15\% (excès notable, dû au poisson + arachides).
\item \textbf{Lipides :} 53,0\% $\gg$ 35\% (excès majeur, dû aux arachides 49\% lipides + huile de cuisson 100\% lipides). Densité énergétique du repas très élevée.
\item \textbf{Glucides :} 26,0\% $\ll$ 50\% (déficit marqué, portion de riz trop faible par rapport aux sources grasses).
\end{itemize}
\textbf{Commentaire :} Ce repas est \textbf{très déséquilibré} : trop gras, trop protéiné, pas assez glucidique (complexes), et probablement pauvre en fibres (riz blanc, pas de légumes). Il favorise la prise de poids et ne procure pas de satiété durable. Corrections : augmenter le riz (200--250g), supprimer l'huile ajoutée (cuisson vapeur/poêle antiadhésive), réduire arachides (30g), ajouter une portion significative de légumes (feuilles, tomates, carottes) et un fruit.
\end{corrige}

\begin{corrige}[Exercice 6 — Prise de masse grasse annuelle liée à une surconsommation]
\begin{enumerate}
\item \textbf{Excédent énergétique quotidien :}
$2 \times 180\ \text{kcal} + 220\ \text{kcal} = 360 + 220 = \mathbf{580\ kcal/jour}$.
\item \textbf{Excédent énergétique annuel :}
$580\ \text{kcal/jour} \times 365\ \text{jours} = \mathbf{211\,700\ kcal/an}$.
\item \textbf{Prise de masse grasse théorique :}
$\frac{211\,700\ \text{kcal}}{7\,700\ \text{kcal/kg}} = \mathbf{27,5\ kg}$ (arrondi au dixième).
\item \textbf{Réalisme du calcul :} \textbf{Non}, ce calcul théorique maximaliste n'est \textbf{jamais vérifié} dans la réalité. Deux facteurs majeurs (au moins) :
\begin{enumerate}
\item \textbf{Thermogénèse adaptative et augmentation du métabolisme de base :} La prise de poids augmente la masse corporelle (masse grasse + masse maigre de soutien), ce qui élève la dépense énergétique de repos (MB $\propto$ masse maigre) et le coût énergétique des déplacements. L'excédent net diminue donc au fil du temps.
\item \textbf{Régulation homéostatique de l'appétit et de la satiété :} La prise de masse grasse augmente la sécrétion de leptine (hormone anorexigène) qui, chez un sujet sain, inhibe la faim et augmente la dépense (via SNS). Inversement, la restriction ou le surpoids modifie la sensibilité à l'insuline/leptine. Le corps ne stocke pas passivement 100\% de l'excédent ; une partie est dissipée en chaleur (thermogénèse sans frisson, NEAT accru involontaire).
\end{enumerate}
Autres facteurs : coût énergétique de la synthèse des triglycérides (lipogenèse $\approx$ 25\% de l'énergie excédentaire), variations de l'activité physique non exercice (NEAT), maladies intercurrentes.
\end{enumerate}
\end{corrige}

\begin{corrige}[Exercice 7 — Interprétation de courbes de croissance et IMC]
\begin{enumerate}
\item \textbf{Trace schématique de l'évolution de l'IMC :}
\begin{popfigure}
\begin{tikzpicture}
\begin{axis}[
    width=\linewidth, height=8cm,
    xlabel={Âge (ans)}, ylabel={IMC (kg/m²)},
    xmin=9.5, xmax=17.5, ymin=15, ymax=32,
    xtick={10,12,14,16,17}, ytick={16,18,20,22,24,26,28,30},
    grid=major, grid style={popline},
    axis lines=middle,
    enlargelimits=0.05,
    legend style={at={(0.02,0.98)}, anchor=north west, fill=white, draw=popline}
]
% Reference curves (approximate P50, P85, P97 for boys)
\addplot[popmuted, dashed, thick, domain=10:17, samples=50] {16.5 + (21.5-16.5)/(17-10)*(x-10)}; \addlegendentry{$P_{50}$ (Médiane)}
\addplot[poporange, dashed, thick, domain=10:17, samples=50] {18.5 + (26.5-18.5)/(17-10)*(x-10)}; \addlegendentry{$P_{85}$ (+1 ET $\approx$ Surpoids)}
\addplot[poppink, dashed, thick, domain=10:17, samples=50] {20.5 + (30.5-20.5)/(17-10)*(x-10)}; \addlegendentry{$P_{97}$ (+2 ET $\approx$ Obésité)}
% Patient data
\addplot[popblue, very thick, mark=*, mark size=3pt] coordinates {(10,17.2) (12,19.5) (14,23.8) (16,27.5) (17,28.9)};
\addlegendentry{Adolescent}
\end{axis}
\end{tikzpicture}
\legende{Évolution de l'IMC de l'adolescent (points bleus) par rapport aux courbes de référence OMS garçons ($P_{50}$, $P_{85}$, $P_{97}$ approximatives).}
\end{popfigure}

\item \textbf{Analyse par rapport aux références (garçon) :}
\begin{itemize}
\item \textbf{10 ans :} IMC 17,2. $P_{50} \approx$ 16,5. $P_{85} \approx$ 18,5. $\rightarrow$ \textbf{Normal} (entre $P_{50}$ et $P_{85}$).
\item \textbf{12 ans :} IMC 19,5. $P_{50} \approx$ 17,5. $P_{85} \approx$ 20,0. $\rightarrow$ \textbf{Normal} (proche $P_{85}$).
\item \textbf{14 ans :} IMC 23,8. $P_{50} \approx$ 18,5. $P_{85} \approx$ 22,0. $P_{97} \approx$ 25,0. $\rightarrow$ \textbf{Surpoids} (IMC $>$ $P_{85}$). \textbf{Basculement vers 13--14 ans.}
\item \textbf{16 ans :} IMC 27,5. $P_{50} \approx$ 20,5. $P_{97} \approx$ 27,0. $\rightarrow$ \textbf{Obésité} (IMC $>$ $P_{97}$). \textbf{Basculement vers 15--16 ans.}
\item \textbf{17 ans :} IMC 28,9. $P_{50} \approx$ 21,5. $P_{97} \approx$ 28,5. $\rightarrow$ \textbf{Obésité} confirmée.
\end{itemize}

\item \textbf{Risques métaboliques et cardiovasculaires :}
Trajectoire ascendante continue franchissant les seuils de surpoids puis d'obésité avant 18 ans $\rightarrow$ risque très élevé de :
\begin{itemize}
\item \textbf{Syndrome métabolique} précoce : insulinorésistance, dyslipidémie athérogène (TG $\uparrow$, HDL $\downarrow$), HTA, stéatose hépatique non alcoolique (NAFLD/NASH).
\item \textbf{Diabète de type 2} d'apparition jeune (avant 25 ans), plus agressif.
\item \textbf{Complications macrovasculaires précoces} : athérosclérose coronaire, AVC, artériopathie des membres inférieurs.
\item \textbf{Retentissement psycho-social} : stigmatisation, dépression, troubles du comportement alimentaire.
\end{itemize}

\item \textbf{Deux axes d'intervention prioritaires :}
\begin{enumerate}
\item \textbf{Alimentation :} Réduire drastiquement la densité énergétique : supprimer boissons sucrées/jus industriels, limiter graisses ajoutées (huile, beurre, sauces grasses), augmenter volume légumes (fibres, satiété) à chaque repas, privilégier céréales complètes/tubercule (ignome, plantain, maïs) vs riz/pâtes blanches, répartir sur 3 repas + 1 collation structurée.
\item \textbf{Activité physique / Sédentarité :} Atteindre \textbf{60 min/jour} d'AP modérée à intense (marche rapide vélo, foot, danse, natation) + renforcement musculaire 3$\times$/sem. Réduire temps assis/écrans $<$ 2h/jour (hors devoirs). Intégrer la marche active pour les trajets quotidiens.
\end{enumerate}
\end{enumerate}
\end{corrige}

\begin{corrige}[Exercice 8 — Construire un programme de prévention personnalisé]
\begin{enumerate}
\item \textbf{IMC de Fatou :}
$\text{IMC} = \frac{68}{1,62^2} = \frac{68}{2,6244} = 25,9\ \text{kg/m}^2$.
\textbf{Classification : Surpoids} (25--29,9). Proche du seuil d'obésité (30).
\item \textbf{Trois erreurs majeures (repères PNNS/Cameroun) :}
\begin{enumerate}
\item \textbf{Excès de densité énergétique et graisses cachées :} Sauces très grasses (huile rouge, arachide, viande grasse) à midi ; beignets frits + fromage + pâtes le soir ; chips/bonbons en collation. Apport lipidique $\gg$ 35\% AET, mostly AGS/AGT.
\item \textbf{Pauvreté en fibres, micronutriments et densité nutritionnelle :} Quasi absence de légumes (peu au déjeuner, aucun le soir), fruits absents, céréales raffinées (pain blanc, riz blanc, pâtes blanches). Apport en fibres $<$ 15g/j (cible $>$ 25g), carences en K, Mg, vitamines antioxydantes.
\item \textbf{Consommation excessive de sucres libres et grignotage non structuré :} Lait sucré + confiture + pain au chocolat + soda + yaourt sucré + bonbons = sucres ajoutés $\gg$ 10\% AET (OMS). Collations ultra-transformées (chips, bonbons) remplacent des aliments nutritifs. Pas de structure « 4 repas » claire.
\end{enumerate}
\item \textbf{Programme journalier pour Fatou (aliments locaux accessibles) :}
\begin{center}
\begin{tabularx}{\linewidth}{|l|X|}\hline
\textbf{Repas} & \textbf{Menu équilibré (exemple)} \\ \hline
\textbf{Petit-déjeuner} (7h) & Bouillie de maïs/soja/arachide \textbf{peu sucrée} (200 mL) + \textbf{1 fruit de saison} (banane, papaye, ananas) + 1 tranche pain complet/manioc + eau/boisson non sucrée. \\ \hline
\textbf{Déjeuner} (13h) & \textbf{Riz complet / Maïs / Fonio} (150--200g cuit) + \textbf{Ndolé / Feuille d'okok / Folong} (cuisiné \textbf{peu d'huile}, poisson fumé/viande maigre/crevettes) + \textbf{Salade crue} (tomate, concombre, carotte râpée, vinaigrette légère) + \textbf{1 fruit}. \\ \hline
\textbf{Collation} (16h) & \textbf{1 fruit} (mangue, goyave, orange) + \textbf{30g arachides grillées non salées} OU 1 galette de maïs/manioc maison + eau. \\ \hline
\textbf{Dîner} (19h30) & \textbf{Légumes vapeur/sauté} (chou, carotte, haricot vert) + \textbf{Poisson grillé / Œufs / Poulet sans peau} (100g) + \textbf{Plantain bouilli / Igname / Patate douce} (150g) + Yaourt nature/fromage blanc 0\% + eau. \\ \hline
\end{tabularx}
\end{center}
\textbf{Plan d'activité physique (objectif 60 min/jour) :}
\begin{itemize}
\item \textbf{Quotidien (30 min) :} Marche active aller-retour école/domicile (remplace taxi-moto si distance $<$ 2 km) OU corde à sauter / danse / footing 30 min le soir.
\item \textbf{3 fois/semaine (45--60 min) :} Sport loisir structuré (football, basketball, handball, danse traditionnelle, natation) -- intensité modérée à soutenue (essoufflement, transpiration).
\item \textbf{EPS scolaire (2h/sem) :} Participation active maximale.
\item \textbf{Règle écrans :} $\le$ 1h/jour loisir ; pause active 5 min toutes les 45 min devoirs.
\end{itemize}
\textbf{Total hebdo :} $\approx$ 420 min (60 min/jour moyen).

\item \textbf{Prévention du diabète de type 2 :} Ce programme agit sur les trois leviers physiopathologiques :
\begin{enumerate}
\item \textbf{Perte de masse grasse viscérale} (déficit énergétique modéré + AP) $\rightarrow$ baisse sécrétion adipokines pro-inflammatoires (TNF-$\alpha$, IL-6) $\rightarrow$ amélioration sensibilité à l'insuline.
\item \textbf{Activité physique régulière} $\rightarrow$ translocation $GLUT_4$ indépendante de l'insuline dans le muscle $\rightarrow$ captage glucose $\uparrow$, glycémie post-prandiale $\downarrow$, sensibilité insulinique $\uparrow$ (effet 24--48h post-séance).
\item \textbf{Alimentation à faible charge glycémique} (céréales complètes, légumineuses, fibres, peu sucres ajoutés) $\rightarrow$ sécrétion insulinique modérée, pas de pics glycémiques $\rightarrow$ préservation cellules $\beta$ pancréatiques.
\end{enumerate}
\end{enumerate}
\end{corrige}

\courssub{Je me prépare au Probatoire}

\begin{corrige}[Exercice 9 — Sujet type Probatoire : Cas clinique complet — « Le cas de Stéphane »]
\textbf{Barème indicatif : 20 points. Points de vigilance :} Justifier chaque chiffre, utiliser le vocabulaire précis (adiposité androïde, syndrome métabolique, insulinorésistance), lier biologie/clinique/comportement, conseils concrets et progressifs.

\begin{enumerate}
\item \textbf{Anthropométrie (4 pts) :}
\begin{align*}
\text{IMC} &= \frac{92}{1,75^2} = \frac{92}{3,0625} = \mathbf{30,0\ kg/m^2} \quad \text{(Obésité Grade 1, seuil $\ge$ 30)} \\
\text{Ratio Tour de taille / Taille} &= \frac{98}{175} = \mathbf{0,56} \quad (> 0,50)
\end{align*}
\textbf{Conclusion :} Obésité confirmée par l'IMC. Ratio $> 0,5$ (ou tour de taille 98 cm $>$ seuil homme 94 cm) signe une \textbf{adiposité androïde (viscérale / abdominale)}, forme la plus métaboliquement délétère (drainage veineux porte direct vers le foie).

\item \textbf{Bilan alimentaire 24h (6 pts) :}
\textit{Estimation macronutriments (valeurs moyennes standard) :}
\begin{itemize}
\item Pdj : Lait entier 200mL (P 6, L 7, G 10) + Pain blanc 100g (P 8, L 1, G 55) + Beurre 30g (L 24) + Confiture 20g (G 13) $\rightarrow$ $\approx$ P 14g, L 32g, G 78g $\rightarrow$ $\approx$ 630 kcal.
\item Déj : Riz blanc 200g cuit (P 5, L 0,5, G 56) + Ndolé 100g (feuilles + huile 25g $\approx$ L 22,5 + viande grasse 50g P 10 L 15 + crevettes) + Mangue 200g (G 28) $\rightarrow$ $\approx$ P 25g, L 38g, G 84g $\rightarrow$ $\approx$ 820 kcal.
\item Goûter : Pain chocolat 80g (P 6, L 12, G 35) + Soda 33cL (G 35) $\rightarrow$ $\approx$ P 6g, L 12g, G 70g $\rightarrow$ $\approx$ 440 kcal.
\item Dîner : Pâtes 250g cuites (P 9, L 1, G 70) + Sauce tomate indus 100g (L 5, G 10) + Fromage râpé 100g (P 25, L 28) + Yaourt sucré 125g (P 4, L 3, G 16) $\rightarrow$ $\approx$ P 38g, L 37g, G 96g $\rightarrow$ $\approx$ 900 kcal.
\end{itemize}
\textbf{Total journalier estimé :} $\approx$ \textbf{2790 kcal} (proche des 2800 kcal besoins théorique garçon 17 ans \emph{actif}).
\textbf{Commentaire équilibre :}
\begin{itemize}
\item \textbf{Répartition P/L/G :} P $\approx$ 83g (12\%), L $\approx$ 119g (38\%), G $\approx$ 328g (47\%). Lipides \textbf{trop élevés} ($>$35\%), Glucides \textbf{trop bas} ($<$50\%) et surtout \textbf{sucres libres excessifs} (confiture, soda, chocolat, yaourt, sauce indus), Fibres \textbf{très faibles} (riz/pâtes/pain blancs, peu légumes, fruits sucrés seulement).
\item \textbf{Densité énergétique élevée} (graisses cachées, fritures, fromage, chocolat) $\rightarrow$ satiété faible $\rightarrow$ surconsommation facilitée.
\item \textbf{Qualité :} Excès AGS (beurre, fromage, viande grasse, huile de cuisson), excès sucres ajoutés, déficit fibres/vitamines/antioxydants (légumes verts quasi absents).
\end{itemize}

\item \textbf{Analyse biologique et syndrome métabolique (5 pts) :}
\begin{itemize}
\item \textbf{Glycémie à jeun 1,15 g/L :} \textbf{Hyperglycémie à jeun modifiée (HGM) / Prédiabète} (norme $<$ 1,10 ; diabète $\ge$ 1,26).
\item \textbf{HDL-cholestérol 0,35 g/L :} \textbf{Hypo-HDL} (norme $>$ 0,40 g/L homme), facteur de risque CV majeur.
\item \textbf{Triglycérides 1,80 g/L :} \textbf{Hypertriglycéridémie} (norme $<$ 1,50 g/L), marqueur d'insulinorésistance hépatique (surproduction VLDL).
\item \textbf{TA 135/85 mmHg :} \textbf{HTA stade 1} (norme $<$ 120/80).
\end{itemize}
\textbf{Syndrome métabolique (critères IDF) :} Tour de taille $>$ 94 cm (homme) + 3 autres critères (TG, HDL, Glycémie, TA) $\rightarrow$ \textbf{Syndrome métabolique AVÉRÉ}.
\textbf{Lien physiopathologique :} Adiposité viscérale $\rightarrow$ lipolyse accrue $\rightarrow$ flux d'acides gras libres (AGL) vers foie (veine porte) $\rightarrow$ \textbf{néoglucogenèse hépatique} non freinée par l'insuline (résistance hépatique) + surproduction VLDL (TG $\uparrow$) $\rightarrow$ résistance musculaire (captation glucose $\downarrow$ via $GLUT_4$) $\rightarrow$ \textbf{hyperinsulinisme compensateur} (pancréas sursollicité) $\rightarrow$ à terme épuisement cellules $\beta$ $\rightarrow$ diabète type 2 installé.

\item \textbf{Facteurs de risque cumulés et pathologie majeure (3 pts) :}
\begin{enumerate}
\item \textbf{Anthropométrique :} Obésité androïde (IMC 30, TT 98 cm).
\item \textbf{Biologique :} Syndrome métabolique complet (HGM, Hypo-HDL, HTG, HTA).
\item \textbf{Comportemental 1 :} Sédentarité extrême (0 min sport hors EPS, taxi-moto, 4h écrans/jour).
\item \textbf{Comportemental 2 :} Alimentation hypercalorique, hyperlipidique, hyperglycidique (sucres libres), hypofibrée ; sommeil court (6h).
\end{enumerate}
\textbf{Pathologie majeure menaçante à moyen terme :} \textbf{Diabète de type 2} (avec complications micro/macrovasculaires précoces : néphropathie, rétinopathie, coronaropathie, AVC).

\item \textbf{Fiche conseils personnalisée pour Stéphane (2 pts -- 10 lignes max) :}
\begin{itemize}
\item \textbf{Alimentation :} 1) Remplacer riz/pâtes/pain blancs par \textbf{céréales complètes/tubercule (maïs, ignome, plantain)} à chaque repas ; 2) \textbf{Supprimer boissons sucrées/sodas/jus industriels} (boire 1,5--2L eau/jour) ; 3) \textbf{Cuisiner « léger »} : vapeur, papillote, grillade sans huile ajoutée ; 1 c.à.s huile max/repas ; viandes maigres/poissons ; légumes verts $\ge$ 1/2 assiette.
\item \textbf{Activité physique :} Objectif progressif : Sem 1--2 : \textbf{30 min marche rapide/jour} (trajet école) ; Sem 3--4 : + 2 séances 30 min sport plaisir (foot, corde, danse) ; Mois 2 : \textbf{60 min/jour} (endurance + renforcement 2$\times$/sem). Zéro ascenseur, zéro taxi-moto $<$ 1km.
\item \textbf{Hygiène de vie :} \textbf{Coucher 22h30 max} (cible 8h sommeil) ; \textbf{Écrans loisir $<$ 1h/jour} (pas d'écran au lit, pas pendant repas) ; Gestion stress : respiration/cohérence cardiaque 5 min matin/soir.
\end{itemize}
\end{enumerate}
\end{corrige}

\begin{corrige}[Exercice 10 — Sujet type Probatoire : Analyse de documents — « Transition nutritionnelle au Cameroun »]
\textbf{Barème indicatif : 20 points. Points de vigilance :} Calculs exacts (% évolution), lien explicite Doc 2 $\rightarrow$ bilan énergétique, mécanisme moléculaire Doc 3 (adipokines $\rightarrow$ IR $\rightarrow$ néoglucogenèse/captation), complications classées micro/macro, mesures préventives \emph{populationnelles} (pas individuelles) et contextualisées Cameroun.

\begin{enumerate}
\item \textbf{Exploitation Document 1 (4 pts) :}
\begin{align*}
\text{Femmes : } \frac{20,1 - 10,5}{10,5} \times 100 &= \frac{9,6}{10,5} \times 100 = \mathbf{91,4\%} \text{ d'augmentation relative.} \\
\text{Hommes : } \frac{8,4 - 3,2}{3,2} \times 100 &= \frac{5,2}{3,2} \times 100 = \mathbf{162,5\%} \text{ d'augmentation relative.}
\end{align*}
\textbf{Tendance :} Augmentation spectaculaire et continue de la prévalence de l'obésité en 14 ans, \textbf{plus rapide chez les hommes} (décollage d'une prévalence très basse) mais \textbf{niveau absolu resté plus élevé chez les femmes} (facteurs hormonaux, parité, activité physique moindre, perception culturelle).
\textbf{Deux hypothèses explicatives (Transition nutritionnelle/épidémiologique) :}
\begin{enumerate}
\item \textbf{Transition alimentaire :} Passage d'un régime traditionnel (tubercules, feuilles, poisson fumé/viande maigre, faible densité énergétique, riches fibres) à un régime « occidentalisé » urbain : excès riz/pâtes blanches, huiles raffinées (arachide, palme raffinée), sucre ajouté (boissons, pâtisseries), produits ultra-transformés (bouillons cubes, conserves), viande grasse/abats. Disponibilité/accessibilité accrue aliments bon marché à haute densité énergétique.
\item \textbf{Transition activité/environnement :} Urbanisation rapide $\rightarrow$ sédentarité professionnelle (tertiaire, transport motorisé « bendy-skin », taxi-moto), disparition de l'activité physique domestique/agricole (marche, port de charges), insécurité/absence d'espaces publics pour le sport, stress chronique/cortisol, réduction temps sommeil.
\end{enumerate}

\item \textbf{Analyse Document 2 -- Paul (4 pts) :}
\begin{center}
\begin{tabularx}{\linewidth}{|l|X|X|}\hline
\textbf{Pilier} & \textbf{Autrefois (Village)} & \textbf{Aujourd'hui (Douala)} \\ \hline
\textbf{Alimentation} & Manioc, feuilles vertes, poisson fumé (faible DG, riches fibres, protéines maigres) & Riz, sauce arachide \textbf{très grasse}, beignets \textbf{à chaque coin de rue} (haute DG, AGS, sucres, sel) \\ \hline
\textbf{Activité physique} & Marche 5 km/jour (champ, eau, marché) = \textbf{AP de vie intense} & \textbf{Bendy-skin pour 500 m} = \textbf{sédentarité totale}, travail bureau assis \\ \hline
\textbf{Environnement / Comportement} & Rythme solaire, sommeil long, faible stress psychosocial & \textbf{Stress travail}, \textbf{coucher tardif}, sommeil court, environnement \textbf{obésogène} (pub, disponibilité alimentaire 24/7) \\ \hline
\end{tabularx}
\end{center}
\textbf{Bilan énergétique positif :} \textbf{Apports $\uparrow\uparrow$} (densité énergétique $\times$ 2-3, grignotage constant) + \textbf{Dépenses $\downarrow\downarrow$} (METs quotidiens divisés par 3-4) = Excédent journalier massif $\rightarrow$ +18 kg en 3 ans (cohérent avec calcul théorique Ex 6).

\item \textbf{Mécanisme obésité abdominale $\rightarrow$ DT2 (Document 3) (5 pts) :}
L'\textbf{adiposité viscérale} (hypertrophie adipocytes) crée un état d'\textbf{hypoxie tissulaire} et de stress réticuloplasmique $\rightarrow$ infiltration macrophagique + sécrétion d'\textbf{adipokines pro-inflammatoires} (TNF-$\alpha$, IL-6, résistine) et \textbf{baisse d'adiponectine} (insulinosensibilisante).
\begin{enumerate}
\item \textbf{Foie :} TNF-$\alpha$/AGL $\rightarrow$ \textbf{résistance à l'insuline hépatique} $\rightarrow$ l'insuline ne freine plus la \textbf{néoglucogenèse} ni la glycogenolyse $\rightarrow$ \textbf{production hépatique de glucose excessive} (jeûne et post-prandial).
\item \textbf{Muscle :} Résistance au récepteur insulinique/$IRS_1$ $\rightarrow$ défaut translocation \textbf{$GLUT_4$} $\rightarrow$ \textbf{captation musculaire du glucose altérée} $\rightarrow$ hyperglycémie post-prandiale majeure.
\item \textbf{Tissu adipeux :} Lipolyse non inhibée $\rightarrow$ flux AGL constant $\rightarrow$ entretien résistance hépatique/musculaire + stéatose ectopique (foie, pancréas, muscle).
\item \textbf{Pancréas :} Hyperglycémie chronique $\rightarrow$ stimulation continue cellules $\beta$ $\rightarrow$ \textbf{hyperinsulinisme compensateur} (hyperinsulinémie à jeun/post-charge) $\rightarrow$ à long terme \textbf{épuisement/apoptose cellules $\beta$} (glucolipotoxicité, amyloidose) $\rightarrow$ \textbf{insulinopénie relative} $\rightarrow$ \textbf{diabète de type 2 clinique}.
\end{enumerate}

\item \textbf{Complications et dépistage (3 pts) :}
\begin{itemize}
\item \textbf{Microvasculaires :} \textbf{Néphropathie diabétique} (insuffisance rénale terminale), \textbf{Rétinopathie} (cécité), \textbf{Neuropathie périphérique} (douleurs, pied diabétique, troubles érectiles).
\item \textbf{Macrovasculaires :} \textbf{Coronaropathie} (IDM silencieux), \textbf{Accident Vasculaire Cérébral (AVC)}, \textbf{Artériopathie Oblitérante des Membres Inférieurs (AOMI)} (amputation).
\item \textbf{Dépistage annuel néphropathie :} \textbf{Dosage de l'albuminurie urinaire} (recherche de microalbuminurie : ratio Albumine/Créatinines urinaires $<$ 30 mg/g ou 30--300 mg/g) + \textbf{Créatininémie / DFG (formule CKD-EPI)}.
\end{itemize}

\item \textbf{Mesures de prévention primaire populationnelle (Cameroun) (4 pts) :}
\begin{enumerate}
\item \textbf{Fiscalité nutritionnelle (taxe « soda » étendue) :} Instaurer une \textbf{taxe excise significative ($\ge$ 20\%) sur les boissons sucrées (sodas, jus néctars, boissons énergisantes) et aliments ultra-transformés riches en sel/sucre/gras} (nouilles instantanées, biscuits, chips). Affecter les recettes au subventionnement des \textbf{fruits, légumes, légumineuses, céréales complètes locaux} (maïs, mil, sorgho, ignome, plantain, feuilles) pour les rendre plus accessibles aux populations vulnérables.
\item \textbf{Aménagement urbain actif :} Obliger les mairies (Yaoundé, Douala, Bafoussam, Bamenda...) à créer des \textbf{pistes cyclables protégées, trottoirs larges et éclairés, parcs urbains gratuits avec équipements sportifs} (ex: berges du Wouri, lac municipal). Planifier la ville « 15 minutes » (commerces, écoles, santé accessibles à pied/vélo). Sécuriser les abords d'écoles pour la marche/vélo (vélo-bus, pédibus).
\item \textbf{Régulation de l'environnement alimentaire \& Éducation :} \textbf{Interdire la publicité} pour les aliments/boissons à profil nutritionnel défavorable (OMS/AFRO nutrient profile model) \textbf{ciblant les enfants/ados} (TV, radio, réseaux sociaux, abords écoles, événements sportifs). \textbf{Étiquetage nutritionnel frontal obligatoire} (type Nutri-Score adapté contexte local) sur tous les préemballés. \textbf{Normes strictes pour les cantines scolaires} : repas équilibrés (légumes, protéines maigres, féculents complets, fruit, eau), interdiction vendeurs ambulants « beignets/sodas » dans l'enceinte et 100m autour des écoles. Intégration module « Nutrition/AP » obligatoire dans programmes scolaires (primaire $\rightarrow$ secondaire).
\end{enumerate}
\end{enumerate}
\end{corrige}

\begin{corrige}[Exercice 11 — Sujet type Probatoire : Rédaction guidée — « Lettre à un ami »]
\textbf{Barème : 10 points (Rigueur scientifique 4 pts, Clarté/Structure 3 pts, Ton bienveillant/ferme 3 pts). Les 4 notions obligatoires doivent apparaître explicitement et correctement expliquées.}

\begin{exemplebox}[Proposition de réponse type (environ 15 lignes)]
Salut Armand,

Je t'écris parce que je tiens à toi et que je m'inquiète. Ton IMC de 29 kg/m² te place en \textbf{surpoids avéré} (seuil 25), aux portes de l'\textbf{obésité} (seuil 30) ; ce n'est pas qu'une question d'esthétique, c'est un marqueur d'adiposité viscérale qui active l'inflammation silencieuse.

Ton alimentation « n'importe quoi » (fast-food, sodas) casse l'\textbf{équilibre alimentaire} : elle apporte trop de graisses saturées, de sucres libres et de calories vides, mais presque pas de fibres, vitamines ni antioxydants. Cette \textbf{haute densité énergétique et faible densité nutritionnelle} surcharge ton foie et ton pancréas.

Physiologiquement, l'excès de graisse abdominale libère des cytokines (TNF-$\alpha$, IL-6) qui provoquent une \textbf{résistance à l'insuline} : tes muscles et ton foie ne répondent plus à l'insuline, le glucose s'accumule dans le sang. Ton pancréas surcompense (\textbf{hyperinsulinisme}) jusqu'à l'\textbf{épuisement de ses cellules $\beta$} : c'est le mécanisme inexorable du \textbf{diabète de type 2}, qui touche de plus en plus de jeunes de 20 ans. « Se sentir bien » aujourd'hui ne veut pas dire « être en santé » demain ; le diabète est silencieux 5 à 10 ans avant le diagnostic.

La \textbf{recommandation OMS de 60 min/jour d'activité modérée à intense} (marche rapide, foot, danse, vélo) n'est pas une option : elle \textbf{sensibilise tes cellules à l'insuline} (via $GLUT_4$), brûle le glucose et les graisses, protège ton cœur et ton cerveau. Couplée à un sommeil de 8h (pas 5h !), c'est ta meilleure assurance vie.

On peut commencer ensemble cette semaine : 30 min marche ce soir, on remplace le soda par l'eau, on ajoute des légumes au dîner. Je suis là pour t'aider, pas te juger. Réfléchis-y, ton corps te le rendra.

Ton pote,
[Ton Prénom]
\end{exemplebox}

\textbf{Points de vigilance pour la notation :}
\begin{itemize}
\item \textbf{Notion 1 (IMC) :} Valeur 29 citée, seuils 25/30 cités, interprétation « surpoids / risque obésité / adiposité viscérale ».
\item \textbf{Notion 2 (Équilibre) :} Mots-clés : densité énergétique, densité nutritionnelle, fibres, sucres libres, graisses saturées, déséquilibre P/L/G.
\item \textbf{Notion 3 (Physiopatho DT2) :} Mots-clés : adiposité viscérale $\rightarrow$ inflammation/cytokines $\rightarrow$ résistance insuline (foie/muscle) $\rightarrow$ hyperinsulinisme compensateur $\rightarrow$ épuisement cellules $\beta$ $\rightarrow$ hyperglycémie chronique. Pas de confusion avec DT1.
\item \textbf{Notion 4 (AP 60 min) :} 60 min/jour modérée à intense citée, mécanismes : sensibilisation insuline ($GLUT_4$), dépense énergétique, prévention prise de poids, bénéfices cardiovasculaires.
\item \textbf{Ton :} « Je tiens à toi », « on peut commencer ensemble », « pas te juger » = bienveillance. « Ce n'est pas qu'esthétique », « silencieux 5-10 ans », « assurance vie » = fermeté scientifique.
\end{itemize}
\end{corrige}

\newpage

\coursec{Fiche bilan}

\begin{popfigure}
\centering
\begin{center}\tcbox[colback=popblue,colframe=popblue,coltext=white,arc=9pt,boxsep=4pt,left=6pt,right=6pt]{\bfseries\small Lutte contre la mauvaise alimentation et l'inactivité physique}\end{center}
\vspace{-2pt}
\begin{tcbraster}[raster columns=3,raster equal height=rows,raster column skip=6pt,raster row skip=6pt,enhanced,blankest]
\begin{tcolorbox}[enhanced,colback=white,colframe=popgreen,boxrule=0.9pt,arc=3pt,title={1. Besoins nutritionnels de l'adolescent},fonttitle=\bfseries\scriptsize,coltitle=white,colbacktitle=popgreen,left=4pt,right=4pt,top=2pt,bottom=2pt,boxsep=1pt,toptitle=1.5pt,bottomtitle=1.5pt,halign=left]
\begin{itemize}[nosep,leftmargin=*,itemsep=1pt,font=\scriptsize]
\item ANC Énergétiques
\item Macronutriments
\item Micronutriments (Fe, Ca)
\item Eau \& Fibres
\end{itemize}
\end{tcolorbox}
\begin{tcolorbox}[enhanced,colback=white,colframe=poporange,boxrule=0.9pt,arc=3pt,title={2. Évaluation d'une ration},fonttitle=\bfseries\scriptsize,coltitle=white,colbacktitle=poporange,left=4pt,right=4pt,top=2pt,bottom=2pt,boxsep=1pt,toptitle=1.5pt,bottomtitle=1.5pt,halign=left]
\begin{itemize}[nosep,leftmargin=*,itemsep=1pt,font=\scriptsize]
\item Bilan quantitatif
\item Bilan qualitatif
\item Répartition repas
\item Score diversité
\end{itemize}
\end{tcolorbox}
\begin{tcolorbox}[enhanced,colback=white,colframe=poppurple,boxrule=0.9pt,arc=3pt,title={3. IMC : Calcul \& Interprétation},fonttitle=\bfseries\scriptsize,coltitle=white,colbacktitle=poppurple,left=4pt,right=4pt,top=2pt,bottom=2pt,boxsep=1pt,toptitle=1.5pt,bottomtitle=1.5pt,halign=left]
\begin{itemize}[nosep,leftmargin=*,itemsep=1pt,font=\scriptsize]
\item Formule : $P/t^2$
\item Classes OMS
\item Risques santé
\item Limites (muscle)
\end{itemize}
\end{tcolorbox}
\begin{tcolorbox}[enhanced,colback=white,colframe=poppink,boxrule=0.9pt,arc=3pt,title={4. Conséquences pathologiques},fonttitle=\bfseries\scriptsize,coltitle=white,colbacktitle=poppink,left=4pt,right=4pt,top=2pt,bottom=2pt,boxsep=1pt,toptitle=1.5pt,bottomtitle=1.5pt,halign=left]
\begin{itemize}[nosep,leftmargin=*,itemsep=1pt,font=\scriptsize]
\item Obésité (abdominale)
\item Diabète Type 2
\item Maladies Cardio-Vasc.
\item Syndrome métabolique
\end{itemize}
\end{tcolorbox}
\begin{tcolorbox}[enhanced,colback=white,colframe=popgold,boxrule=0.9pt,arc=3pt,title={5. Recommandations Activité Physique},fonttitle=\bfseries\scriptsize,coltitle=white,colbacktitle=popgold,left=4pt,right=4pt,top=2pt,bottom=2pt,boxsep=1pt,toptitle=1.5pt,bottomtitle=1.5pt,halign=left]
\begin{itemize}[nosep,leftmargin=*,itemsep=1pt,font=\scriptsize]
\item OMS : 60 min/j (ado)
\item Intensité modérée à soutenue
\item Renforcement 3×/sem
\item Limiter sédentarité
\end{itemize}
\end{tcolorbox}
\begin{tcolorbox}[enhanced,colback=white,colframe=popdark,boxrule=0.9pt,arc=3pt,title={6. Recommandations Alimentaires},fonttitle=\bfseries\scriptsize,coltitle=white,colbacktitle=popdark,left=4pt,right=4pt,top=2pt,bottom=2pt,boxsep=1pt,toptitle=1.5pt,bottomtitle=1.5pt,halign=left]
\begin{itemize}[nosep,leftmargin=*,itemsep=1pt,font=\scriptsize]
\item Fruits/Légumes 5/j
\item Céréales complètes
\item Limiter sel, sucre, gras saturés
\item Hydratation eau
\end{itemize}
\end{tcolorbox}
\begin{tcolorbox}[enhanced,colback=white,colframe=popmuted,boxrule=0.9pt,arc=3pt,title={7. Programme prévention personnalisé},fonttitle=\bfseries\scriptsize,coltitle=white,colbacktitle=popmuted,left=4pt,right=4pt,top=2pt,bottom=2pt,boxsep=1pt,toptitle=1.5pt,bottomtitle=1.5pt,halign=left]
\begin{itemize}[nosep,leftmargin=*,itemsep=1pt,font=\scriptsize]
\item Diagnostic (IMC, ration)
\item Objectifs SMART
\item Plan alimentaire
\item Plan activité
\item Suivi \& Évaluation
\end{itemize}
\end{tcolorbox}
\end{tcbraster}

\legende{Carte mentale de la leçon NCT06 : 7 piliers pour prévenir les maladies liées au mode de vie.}
\end{popfigure}

\begin{bilan}

\courssub{Définitions clés}
\begin{itemize}
  \item \cle{Équilibre alimentaire} : Ration couvant les besoins en énergie, macronutriments (protides, lipides, glucides), micronutriments (vitamines, minéraux), fibres et eau, répartie sur 3 à 4 prises/jour.
  \item \cle{Besoins nutritionnels (ANC)} : Apports Nutritionnels Conseillés pour couvrir les besoins de 97,5\% de la population ; varient selon âge, sexe, activité physique, état physiologique.
  \item \cle{Sédentarité} : Éveil en position assise/allongée avec dépense énergétique $\le$ 1,5 MET (ex: écrans, transport motorisé). Différent de l'inactivité physique (non-respect recommandations OMS).
  \item \cle{IMC (Indice de Masse Corporelle)} : $IMC = \frac{Masse\ (kg)}{Taille^2\ (m^2)}$. Indicateur de corpulence, pas de composition corporelle.
  \item \cle{Obésité} : Excès de masse grasse (IMC $\ge$ 30 kg/m² adulte ; courbes de référence ado). Facteur de risque majeur maladies chroniques.
  \item \cle{Diabète de type 2 (DT2)} : Hyperglycémie chronique par résistance à l'insuline + défaut de sécrétion. Lié surpoids, sédentarité, alimentation déséquilibrée.
  \item \cle{Maladies Cardio-Vasculaires (MCV)} : Atteintes cœur/vaisseaux (HTA, AVC, IDM). Facteurs modifiables : alimentation, inactivité, tabac, alcool.
\end{itemize}

\courssub{Formules \& Repères numériques à connaître}
\begin{center}
\begin{tabularx}{\linewidth}{|l|X|X|}\hline
\rowcolor{popblueL}
\textbf{Notion} & \textbf{Formule / Valeur de référence} & \textbf{Commentaire / Unités} \\ \hline
IMC & $\displaystyle IMC = \frac{m}{t^2}$ & $m$ en kg, $t$ en m. Unité : kg/m² \\ \hline
Interprétation IMC (Adulte/OMS) & Maigreur $<$ 18,5 \textbar\ Normal 18,5--24,9 \textbar\ Surpoids 25--29,9 \textbar\ Obésité $\ge$ 30 & Chez l'ado : utiliser courbes de référence (percentiles) âge/sexe \\ \hline
Tour de taille (risque abdominal) & Homme $\ge$ 94 cm (risque accru $\ge$ 102) \textbar\ Femme $\ge$ 80 cm (risque accru $\ge$ 88) & Mesuré au niveau du dernier rebord costal \\ \hline
Besoins Énergétiques (ANC ado 15--17 ans) & Garçon $\approx$ 2700--3100 kcal/j \textbar\ Fille $\approx$ 2100--2400 kcal/j & Selon NAP (Niveau Activité Physique) : 1,4 à 2,0 \\ \hline
Répartition macronutriments & Protides 10--15\% \textbar\ Lipides 30--35\% \textbar\ Glucides 50--55\% & Dont sucres ajoutés $<$ 10\% AET (idéal $<$ 5\%) \\ \hline
Fibres (ANC ado) & $\ge$ 25 g/j & Fruits, légumes, céréales complètes, légumineuses \\ \hline
Activité Physique (OMS 5--17 ans) & $\ge$ 60 min/jour intensité modérée à soutenue & + Activités renforcement musculaire/os 3 fois/sem \\ \hline
\end{tabularx}
\end{center}

\courssub{Méthodes express}
\begin{methode}[Évaluer une ration alimentaire (24h ou type)]
  \begin{enumerate}
    \item \textbf{Lister} tous aliments/boissons consommés (poids ou mesures usuelles).
    \item \textbf{Quantifier} apports : Énergie (kcal), Protides/Lipides/Glucides (g), Fibres, Fer, Calcium, Sel (table Ciqual/ANC).
    \item \textbf{Comparer} aux ANC de l'adolescent (sexe, âge, NAP).
    \item \textbf{Analyser} répartition : \% énergie par macronutriment, nombre de portions fruits/légumes, produits laitiers, féculents complets.
    \item \textbf{Conclure} : Équilibrée / Déséquilibrée (excès/défauts) + 2--3 corrections concrètes.
  \end{enumerate}
\end{methode}

\begin{methode}[Interpréter l'IMC et relier aux risques]
  \begin{enumerate}
    \item Calculer IMC = $m/t^2$.
    \item Situer sur grille OMS (adulte) ou courbes percentiles (ado : carnet de santé).
    \item Croiser avec \textbf{tour de taille} (risque métabolique abdominal).
    \item Évaluer risques : Surpoids/Obésité $\rightarrow$ Risque accru DT2, HTA, dyslipidémie, MCV, apnée sommeil, ostéoarticulaire.
    \item \textbf{Nuancer} : Sportif musclé (IMC élevé, masse grasse normale) ; Personne âgée/sarcopénique (IMC normal, masse grasse élevée).
  \end{enumerate}
\end{methode}

\begin{methode}[Élaborer un programme prévention simple (Alimentation + Activité)]
  \begin{enumerate}
    \item \textbf{Diagnostic} : IMC, tour de taille, rappel alimentaire 3j, niveau activité (questionnaire type IPAQ-A), antécédents familiaux.
    \item \textbf{Objectifs SMART} : Spécifiques, Mesurables, Atteignables, Réalistes, Temporels (ex: -2 kg en 2 mois ; marcher 30 min 5×/sem).
    \item \textbf{Plan Alimentaire} : 3 repas + 1 collation ; 1/2 assiette légumes ; féculents complets à chaque repas ; protéines variées ; limiter ULTRA-transformés, boissons sucrées, sel ajouté.
    \item \textbf{Plan Activité} : Quotidien (marche, vélo, escaliers) + 3 séances/sem structurées (cardio + renfo) + limiter temps assis $<$ 2h d'affilée.
    \item \textbf{Suivi} : Carnet alimentaire/activité hebdo ; pesée 1×/sem ; réévaluation à 4--6 semaines ; ajustement.
  \end{enumerate}
\end{methode}

\end{bilan}

\begin{aretenir}[Checklist avant le Probatoire]
  $\square$ Je sais définir l'équilibre alimentaire et citer les 7 groupes d'aliments avec leurs apports principaux.\\
  $\square$ Je connais les ANC énergétiques et en macronutriments pour un adolescent de 1ère (garçon/fille).\\
  $\square$ Je sais calculer un IMC ($m/t^2$) et l'interpréter avec la classification OMS (adulte) ou les courbes de corpulence (ado).\\
  $\square$ Je sais mesurer et interpréter le tour de taille (seuils H/F).\\
  $\square$ Je sais évaluer une ration alimentaire (quantitatif + qualitatif) et proposer des corrections.\\
  $\square$ J'explique le lien physiopathologique : Sédentarité + Déséquilibre alimentaire $\rightarrow$ Surpoids $\rightarrow$ Résistance insuline $\rightarrow$ DT2 / Dyslipidémie $\rightarrow$ Athérosclérose $\rightarrow$ MCV.\\
  $\square$ Je cite les recommandations OMS d'activité physique pour les adolescents (60 min/j modérée à soutenue + renfo 3×/sem).\\
  $\square$ Je cite les repères PNNS / guides alimentaires camerounais : 5 fruits/légumes, féculents complets, limiter sel/sucre/gras saturés/ultra-transformés.\\
  $\square$ Je sais rédiger un programme de prévention personnalisé (Diagnostic $\rightarrow$ Objectifs SMART $\rightarrow$ Plan Alim/Activité $\rightarrow$ Suivi).\\
  $\square$ Je connais les limites de l'IMC (ne distingue pas muscle/gras, ni répartition graisseuse) et l'intérêt du tour de taille.\\
  $\square$ Je sais identifier les idées reçues : « régime restrictif rapide », « sport = maigrir localisé », « produits light = illimités ».\\
  $\square$ Je relie les enjeux locaux : transition nutritionnelle au Cameroun (urbanisation, fast-food, boissons sucrées) et double fardeau malnutrition/obésité.
\end{aretenir}

\findecours

\end{document}
